\quellseite{285}
schicken den befudelten Rock dem Vater heim, und lassen ihn sagen, es habe es ein Thier gethan, wollens keine Schuld noch Wort haben. Also haben die Juden Christum ermordet und verkauft, alle seine Worte vergeben und hinweg geworfen, seine Gnade beschmeißt, und hören noch nicht auf: noch wollen sie unschuldig sein, und nichts gethan haben, stellen sich, als müßten sie nichts darum, geben noch vor, sie seien die lieben Kinder; wollen nicht hören, daß sie Gottes Sohn ermordet haben, sondern sagen, sie haben es einem gethan, der sei ein Mörder und Schalk gewesen.

22. Also muß Gottes Wort und Wahrheit geschmäht werden, darnach sie lassen hingehen und das Maul wischen, wie die Hure, davon Salomo sagt in [den] Sprüchen [Cap. 30, 20.]. Das sind zwei greuliche Sünden, daß sie Gottes Wort lästern und schänden, darnach sagen, sie thun wohl daran; wie Christus sagt im Johanne [Cap. 16, 2.]: „Es wird die Zeit kommen, daß, wer euch tödtet, wird meinen, er thue Gott einen Dienst daran.“

23. Nun, das hat Jakob müssen leiden, wird von seinen Söhnen betrogen, daß er meint, es sei, wie sie sagen. [Es] ist so viel gesagt: Gott stellt sich gegen sie, als wisse er es nicht, läßt sie\footnote{In den alten Ausgaben: „lässet“, das ist, läßt sie.} immer hingehen und Heuchler bleiben, den Rock beflecken und verderben, daß sie meinen, er wisse nicht, was sie gehandelt haben. Denn, wie es uns im Sinne ist, so ist es auch vor Gott; wie ich ihn halte, so habe ich ihn. Darum gehen die Juden nicht anders hin, als haben sie keine Schuld, legen es auf das wilde Thier, das ist, die Heiden, oder sonst böse Leute. Dieweil schweigt Gott stille; doch thut es ihm wehe, und trägt lange Zeit Leide. Welches geschieht durch die Seinen, die es also fühlen; aber zuletzt kommt die Zeit, daß es offenbar wird.

24. Also ist die Figur fein abgemalt, wie es Christo und seinen Christen geht, daß sie müssen geschändet werden von denen, so ihnen die Nächsten sind, und darnach vor Gott kommen, als haben sie wohl gethan; wie sie jetzt aus dem Evangelio Ketzerei machen. Ueber die Plage müssen wir die Schmach dazu haben, als sei es verbieteter Lohn, und sollen sie darum frönen.

25. Daß aber die zwei Brüder nicht gerne sehen, daß Joseph sollte umkommen, sonderlich Ruben, welcher ihn gerne dem Vater hätte wiedergebracht, [das] sind etliche vernünftige Leute, die es gerne gütlich ausrichteten, und schonten; aber es hilft nichts. Das ist dies Capitel, daß eitel Kreuz lehrt, wie die Christen leiden müssen, und daß Evangelium muß geschmäht und verfolgt werden.

\section{Das achtunddreißigste Capitel}

V. 1—30. Es begab sich um dieselbe Zeit, daß Juda hinab zog von seinen Brüdern, und that sich zu einem Mann zu Odollam, der hieß Hira. Und Juda sahe daselbst eines Cananiters Mannes Tochter, der hieß Sua, und nahm sie. Und da er sie beschlief, ward sie schwanger, und gebar einen Sohn, den hieß er Er. Und sie ward aber schwanger, und gebar einen Sohn, den hieß sie Onan. Sie fuhr fort und gebar einen Sohn, den hieß sie Sela. Und er\footnote{Jenaer: sie.} war zu Chesib, da sie ihn gebar.

Aber er war böse vor dem Herrn, darum tödtete ihn der Herr. Da sprach Juda zu Onan: Lege dich zu deines Bruders Weibe, und nimm sie zur Ehe, daß du deinem Bruder Samen erweckest. Aber da Onan wußte, daß der Same nicht sein eigen sein sollte, wenn er sich zu seines Bruders Weibe legte, ließ er es auf die Erde fallen, und verderbete es, auf daß er seinem Bruder nicht Samen gäbe. Da gefiel dem Herrn übel, das er that, und tödtete ihn auch. Da sprach Juda zu seiner Schnur Thamar: Bleibe eine Wittwe in deines Vaters Hause, bis mein Sohn Sela groß wird; denn er gedachte: vielleicht möchte er auch sterben, wie seine Brüder. Also ging Thamar hin, und blieb in ihres Vaters
\quellseite{286}
Hause. Da viel Tage verlaufen waren, starb des Sua Tochter, Juda Weib. Und nachdem Juda ausgetrauert hatte, ging er hinauf, seine Schafe zu scheren, gen Thimnath, mit seinem Hirten Hira von Odollam. Da ward der Thamar angesagt: Siehe, dein Schwäher geht hinauf gen Thimnath, seine Schafe zu scheren. Da legte sie die Wittwenkleider von sich, die sie trug, schleierte und verhüllte sich, und setzte sich vor die Thüre heraus an dem Wege gen Thimnath; denn sie sahe, daß Sela war groß worden, und sie war ihm nicht zum Weibe gegeben. Da sie nun Juda sahe, meinte er, es wäre eine Hure; denn sie hatte ihr Angesicht verdecket, und machte sich zu ihr am Wege, und sprach: Lieber, laß mich bei dir liegen; denn er wußte nicht, daß sie seine Schnur wäre. Sie antwortete: Was willst du mir geben, daß du bei mir liegst? Er sprach: Ich will dir einen Ziegenbock von der Heerde senden. Sie antwortete: So gib mir ein Pfand, bis daß du mir es sendest. Er sprach: Was willst du für ein Pfand, daß ich dir gebe? Sie antwortete: Dein Siegel, und deinen Fächel,\footnote{Fächel = Binde, Schnur (um Hüte, Stiefeln, Kleider).} und deinen Stab, den du in den Händen hast. Da gab er es ihr, und lag bei ihr, und sie ward von ihm schwanger. Und sie machte sich auf, und ging hin, und legte den Schleier ab, und zog ihre Wittwenkleider wieder an. Juda aber sandte den Ziegenbock durch seinen Hirten von Odollam, daß er das Pfand wieder holte von dem Weibe, und er fand sie nicht. Da fragte er die Leute desselbigen Orts, und sprach: Wo ist die Hure, die außen am Wege saß? Sie antworteten: Es ist keine Hure da gewesen. Und er kam wieder zu Juda, und sprach: Ich habe sie nicht gefunden, dazu sagen die Leute desselbigen Orts, es sei keine Hure da gewesen. Juda sprach: Sie habe es ihr, daß wir nicht vielleicht zu Schanden werden, denn ich habe den Bock gesandt, so hast du sie nicht gefunden. Ueber drei Monden ward Juda angesagt: Deine Schnur Thamar hat gehuret, dazu siehe, sie ist von Hurerei schwanger geworden. Juda sprach: Bringet sie hervor, daß sie verbrannt werde. Und da man sie hervor brachte, schickte sie zu ihrem Schwäher, und sprach: Von dem Manne bin ich schwanger worden, deß dies ist. Und sprach: Kennest du auch, wes das Siegel, und der Fächel, und der Stab ist? Juda erkannte es, und sprach: Sie ist gerechter denn ich; denn ich habe sie nicht gegeben meinem Sohne Sela; doch beschlief er sie nicht mehr. Und da sie gebären sollte, wurden Zwillinge in ihrem Leibe erfunden. Und als sie jetzt gebar, that sich eine Hand heraus. Da nahm die Wehmutter, und band einen rothen Faden darum, und sprach: Der wird der erste heraus kommen. Da aber der seine Hand wieder hinein zog, kam sein Bruder heraus; und sie sprach: Warum ist um deinetwillen ein Fach\footnote{Fach = Loch, Riß.} gerissen? Und man hieß ihn Perez.\footnote{Randglosse: Perez. Perez, ein Zerreißer; Serah heißt Ausgang. Hier ist bedeutet, daß die Werthheiligen sich äußerlich stellen, als wollten sie hervor und die Ersten sein, und werden die Letzten, darüber sich ein großes Reißen hebt unter dem Volke Gottes. Aber der rothe Faden um die Hand ist, daß sie fleischliche Heiligkeit wirken, und die rechten Heiligen verwerfen.} Danach kam sein Bruder heraus, der den rothen Faden um seine Hand hatte, und man hieß ihn Serah.

1. Ich habe vor gesagt: wir müssen schier vor ein jegliches Capitel eine eigene Vorrede und Beschönung machen, denn wir sind so zart, daß wir nicht leiden zu reden, noch [zu] hören von menschlicher Geburt, und haben doch daneben getrieben, das greulich zu sagen ist. Es ist wahr, daß dies ist ein eben\footnote{eben = bedeutend, gewaltig, gar sehr. Vgl. Col. 510.} grob Capitel; nun steht es dennoch in der heiligen Schrift, und hat es der Heilige Geist geschrieben, welcher je so reinen Mund und Feder hat, als wir, daß ich es nicht höher zu beschönen weiß, denn also. Hat jemand einen reinern Mund und Ohren, denn er, der mag es lassen stehen; hat er sich es nicht gescheuet noch geschämt zu schreiben, wollen wir es uns auch nicht schämen zu lesen und [zu] hören.

2. Wollte Gott, wir hätten Zucht und Scham gehalten, da wir sie halten sollten, und Unzucht gemieden, wo man sollte; also haben wir es in Schein gewendet. Wo man aus Noth davon reden sollte, haben wir geschwiegen, aber viel ärger getrieben; und wiederum. Der Heilige Geist weiß wohl, was er gemacht hat, so redet er auch von seiner Creatur, wie es geht. Wenn man es hin und her kehrt, so sind wir geschaffen, Frucht zu zeugen und zu tragen; dazu hat er uns gegeben Glieder, Adern, Flüsse, Blut und
\quellseite{287}
Fleisch; wir machen daraus, was wir wollen, so müssen wir Mann und Weib bleiben, und die Natur gehen lassen, wie sie gepflanzt ist. Da sind wir so keusch und züchtig, wollen nichts davon hören; was man sonst treibt, da schert und lacht man von. Das ist die weltliche Weisheit, die alle Gottes Ordnung verkehrt. Ich bin auch nicht geneigt zu Unzucht und Hurerei, sondern zum ehelichen Stande, darum habe ich müssen davon reden und schreiben, den armen Gewissen zu rathen und helfen, und die falsche Keuschheit der Geistlichen niederzulegen. Sollen wir nun ehelich sein, so müssen ja Früchte folgen; sollen aber Früchte folgen, so muß ja Fleisch und Blut seinen natürlichen Raum und Gang haben, oder [es] wird viel ärger.

3. Nun, diese Historie hat Moses hell und grob beschrieben; darum thue die Augen auf, und denke, daß es geschrieben sei, uns zur Lehre, vom Heiligen Geiste. Denn er hätte es nicht dürfen schreiben, wenn er nicht müßte, wie es um den Menschen gethan wäre. Wir sind wohl nicht alle gleich, denn einer hat diese, der andere jene Gaben von Gott; aber wenn Gott nicht hält, ist einer so fromm als der andere; hält er, so stehen wir, und wiederum. Darum läßt er auch solche Historien schreiben, daß wir sehen, wer wir seien. [Wir] dürfen uns nicht besser rühmen, denn dieser Juda; wollte Gott, daß wir alle so fromm wären. Muß er aber das so eben schreiben von dem Patriarchen, der des Stammes unseres Herrn Christi war? Konnte die Schrift sonst keinen zu Schanden machen, denn eben den Höchsten? Ob gleich die andern alle genarrt hätten, sollte er dennoch billig des geschont haben. Was hat der Heilige Geist für Lust dazu, daß er nicht lieber mag von guten Werken schreiben? Unsere Legendeprediger würden sich schämen, wenn sie sollten von einem Heiligen predigen, der nicht eitel köstliche Werke gethan hätte; aber von diesem Heiligen schreibt er doch gar keines, daß, wenn Fleisch und Blut darin fällt, muß es zum Narren darüber werden.

4. Das ist aber die Meinung: Der Mann ist freilich der höchsten Patriarchen einer, sammt seinen Brüdern; aber Gott läßt sie alle sinken, daß sie im Drecke stecken, verrathen und verkaufen ihr Blut und Fleisch, betrieben dazu den alten, elenden Vater so hart. So schämt sich Moses nicht, seine Großväter\footnote{„so“ fehlt in der Jenaer.} zu schmähen, daß genug wäre, wenn er seine Feinde und Heiden so ausrichtete. [Er] macht nichts denn Schälke und Hurentreiber aus ihnen, und sollen Gottes Volk und Christi Großväter sein, dazu die zween, seine Söhne, verpeipelte Buben sind. Der eine, als er sollte seinem Bruder Kinder zeugen, schüttet er es auf die Erde, davon man nicht predigen darf, und geschieht doch, leider, allzuviel mit Mägdlein und Knaben; er aber sagt es dürr heraus, daß er das Weib im Bette hat, aber ehe sie fruchtbar sollte werden, eher beraubt er sie des Samens, daß sie nicht konnte Kinder haben. Es ist wohl grob gepredigt, wir müssen aber doch von den andern groben Dingen auch reden; denn den natürlichen Gang hat Gott ebensowohl geschaffen, als Essen und Trinken. Mann und Weib gehört ehelich zusammen, wer es aber mißbraucht, wird seine Strafe wohl finden.

5. Das ist uns nun, sage ich, alles geschrieben zur Ermahnung, daß wir greifen bei solchen groben Stücken, wie fromm die Natur ist, wenn Gott die Hand abzieht von seinen großen, lieben Heiligen, daß sie eben Buben sind wie wir; also, daß Gott nicht leiden will, daß wir unsere Kräfte und Vermögen rühmen, sondern seine unaussprechliche Barmherzigkeit preisen. Als, daß Juda ein Leben führt, wie ein anderer Heide; noch wird er erhalten im Gottes Reiche, nicht denn durch bloße Gnade und Güte Gottes, der ihm die Sünde schenkt, wiewohl er das Widerspiel verdient hätte; daß vor Gott nichts gelte, denn die Erkenntniß seiner Gnade; wer die sieht, der sieht Gottes Reich; ob er nicht fällt, daß [es] Gottes Huld sei, nicht sein Vermögen; fällt er aber, daß er darum nicht verzage. Also müssen die Heiligen um unsertwillen zu Schanden werden, daß ihr Wesen nichts gelte, sondern allein sein Werk und Gnade. Uns ist nichts angeboren, denn eitel Sünde und Schande; nicht Eine allein, sondern allerlei. Fleisch und Blut taugt nirgend zu, denn zu stehlen, morden, rauben, fluchen und lästern. Darum ist dies ja ein grob, fleischlich Exempel, aber so tröstlich als eines in der Schrift, ob gleich Fleisch und Blut die Nase davor rümpft, und für schambar hält zu hören.

6. Was wollen wir aber zu der Thamar, seiner Schnur, sagen? Sie wußte, daß er ihr Schwäher war, wie besteht sie denn vor Gott, daß sie solches mit Vorsatz thut, und gibt ihm
\quellseite{288}
Ursache zu sündigen, scheuet sich nicht der Freundschaft halben, daß sie seine Söhne gehabt hatte? Denn es war ja verboten, daß Schwur und Schwäher nicht sollten einander berühren; darum er sie auch hernach nicht mehr wollte erkennen, oder zum Weibe haben. Summa, sie hat auch gethan wie eine Hure, wiewohl sie einen Schein vorwendet, und will es schmücken. Ihr erster Mann, Ger, war so schalkhaftig, daß ihn Gott nicht lange leben ließ. Da nahm sie den andern, der war auch ein Erzbube, beraubte sie der Frucht. Der dritte Sohn [Sela] ist nun jung, und der Vater fürchtete sich, er würde auch sterben, wenn er sie nähme. Sie mußte aber, daß ihr der Sohn auch gebührete. So ferne hatte sie Recht und gut Gewissen; als sie aber sahe, daß er ihr ein ander Weib gab, so brauchte sie der Lücke, und machte sich zu ihm, daß er sie nicht kannte. Das war Trügerei, und ein Hurenstück.

7. So muß man sie lassen bleiben Fleisch und Blut, und ihre Werke gering achten, ja, auch übel gethan sein,\footnote{Erlanger: sind.} aber Gottes Gnade darinne sehen. Sie sind beide in großen Sünden. Er thut es in der Meinung, daß sie eine Hure sei; sie aber, daß er ihr Schwäher ist. Das war noch härter; noch begabt sie Gott mit zween Söhnen, und Perez bleibt in der Linie Christi. Das muß man im Evangelio herrlich lesen in aller Welt, und von der Hurerei singen und sagen. Darum wollen wir auch redlich bekennen, daß sie beide grob genarret haben, wiewohl es ihnen Gott geschenkt hat; dazu, daß man sähe, wie Christus gekommen sei um der Sünder willen, ihnen zu helfen, und sich gar nicht schämt, daß er Huren und Buben in seinem Geschlechte hat, und den Keim nicht auslöschen will; so müssen wir ihn auch wohl stehen lassen.

8. Aber damit ist nicht Raum gegeben den Muthwilligen, die da sagen: Nun will ich auch thun, was mich gelüstet. Gott kann wohl durch die Finger sehen mit denen, so da aus Gebrechlichkeit des Fleisches und Blutes fallen, aber Muthwillen und Trotz will er nicht leiden.

9. Weiter sehen wir, daß das Gesetz Mosis sehr alt ist, doch durch Mosen wieder verneuet; als da war, wenn ein Bruder ohne Erben starb, so war der nächste Bruder oder Freund schuldig, das Weib zu nehmen, und ihm Samen zu erwecken [5 Mos. 25, 5.]. [Es] ist ein alter Brauch gewesen, vielleicht von Adam her aufgestommen. So liest man von der Ruth, welche den Boaz nahm; denn er war ihr Nachmann\footnote{Nachmann = der zweite Mann, nachdem der erste gestorben war.} [Ruth 4, 10.].

10. Es ist aber ein seltsam, schwer Gesetze gewesen, und müssen fromme Leute gewesen sein, die es haben können zukommen. Darum lesen wir, daß die Könige und Herren viel Weiber gehabt haben, daß mehrere Theil dabey, daß sie ihnen angestorben sind. [Es] ist den Weibern auch schwer und hart gewesen, daß sie haben zu dem nächsten Freunde gemußt, und oft Aschenbrödel sein.

11. Sonst ist nichts Sonderliches mehr in diesem Capitel, denn das Hauptstück, daß Gott solche schändlichen Historien läßt schreiben, und doch darinne solche edle, tröstliche Lehre von seiner Gnade und Güte vorhält, zu beweisen seine Wunder auch in Sünden. Denn er ist der Gott, der aus Tod Leben, aus Sündern fromme Leute macht, und aus dem, so\footnote{„so“ fehlt in der Erlanger.} verloren ist, das Beste gewinnt.

12. Wenn man nun sollte die Jahre rechnen, ist Juda noch ein junger Mann gewesen, da er gefreiet hat, nämlich, bald darnach, als Joseph verkauft war, welcher nur siebenzehn Jahre alt war; [er] wird auch vielleicht uneins mit den Brüdern worden sein, um derselben Untugend willen, und sich von ihnen geschlagen haben unter die fremden Cananiter; wird nicht viel über achtzehn Jahre gewesen sein, als er das Weib nahm. Denn er hatte drei Söhne, die mannbar sind; so hat das Weib auch bei achtzehn oder zwanzig Jahren müssen sein, als er bei ihr schlief. Das kann man wohl abnehmen, wenn man zurück rechnet nach dem, das Moses hernach [Cap. 41, 46.] schreibt, daß Joseph dreißig Jahre alt war, als er vor Pharao stand. Das sage ich darum, daß man sehe, wie fein Regiment zu der Zeit gewesen ist, daß man die jungen Leute bald zur Ehe gegeben hat, daß desto mehr Unzucht nachbliebe; wiewohl es dennoch geholfen hat, was es konnte. Zuletzt sind noch zwei Kinder da, wie sie geboren werden, welches gehört in die heimliche Deutung, die wollen wir suchen.

13. Des Weibes Name ist Thamar, und heißt
\quellseite{289}
auf deutsch ein Palmbaum. Nun ist unter anderer Art des Baumes, daß er süße Früchte trägt, die man heißt Datteln; so hat das Holz vor andern Bäumen die Art, wenn man Balken daraus macht, und beschwert sie, so beugt es sich empor wider die Last; je mehr man es drückt, je mehr es über sich will. Nun ist in dieser Historie bedeutet das ganze Mysterium und Wesen, so da heißt das Evangelium und Gesetz Gottes. Denn so habe ich gesagt, daß Gott das jüdische Volk sonderlich hat erwählet, daß man heißt die Synagoge, die alte Kirche und Christenheit, die seines Sohnes Braut und sein eigen Volk war; denen hat er gegeben das Gesetz und Priestertum, das Volk täglich zu lehren und unterweisen, daß es fromm würde. Denn also haben wir droben [Cap. 25, § 26. Cap. 31, § 33 ff. Cap. 19, § 23] gehört, daß die Weiber Rebekka, Lea, Rachel, item die zwei Töchter Lot, sind alle Figuren und Vorbilder des Volkes Gottes gewesen, das regiert ist worden durch sein Wort; also, daß der Same, dadurch sie schwanger sind worden, nichts anders ist, denn Gottes Wort, welches die Seele zu sich nimmt durch den Glauben, und bringt davon Früchte der Liebe.

14. So ist nun hier geschrieben, daß sie zween Brüder zur Ehe nimmt, aber den dritten kriegt sie nicht. Darnach kommt sie zu Juda, ihrem Schwäher, unter der Gestalt einer Hure. Also geht es unter dem Volke, bis auf diesen Tag. Der erste Mann war böse vor dem Herrn, darum tödtete er ihn; das sind die Prediger und Lehrer, die sind verzweifelte Buben, sammt dem Volke, so öffentlich wider das Gesetz sündigen, und [es] nicht halten.

15. Der andere aber, der bei ihr liegt, und verderbt es doch, daß sie nicht fruchtbar werde, noch empfangen kann, das sind, die auch das Predigtamt unter den Juden mit dem Scheine haben, aber es wird nicht recht geführt, sondern verderbt, daß keine Frucht davon kommt. Das sind die Heuchler, die wohl den rechten Text Gottes Wortes und des Gesetzes haben, doch nicht so führen, daß es bessere, sondern lassen es fallen, und mengen es unter die Erde, zwingen es auf unser Werk. Denn wenn man das Gesetz recht predigen will, muß man es so treiben, daß es nicht Werke lehre, sondern erkennen, was uns fehlt; wie Paulus Röm. 3, 20. sagt: „Durch das Gesetz kommt Erkenntniß der Sünde“, und daß man die Gewissen davon führe auf das Erkenntniß der Gnade. Das wäre denn Samen recht geführt, davon die Seelen fruchtbar werden. Die nun nicht so führen, sondern auf Menschen Weise und Werk, die kommen nimmermehr zur Gnade, noch zur Frucht. Darum heißt dieser recht Onan, das ist, eine mühselige Predigt, damit nur die Gewissen beschwert und gemartert werden, und niemand genießen kann.

16. Diese zweierlei Leute findet man nun immerdar in dem Volke, entweder die das Gesetz in Wind schlagen, und sich gar in die Schanze geben, oder die es nur zum Scheine führen. Jene sind die ruchlosen, öffentlichen Buben, die andern sind Heuchler und Werkheilige; die werden beide todtgeschlagen, das ist, von Gott verdammt und verworfen.

17. Da es so weit kam, sollte nun der dritte Sohn das Weib nehmen; aber sie kriegte ihn nicht. Da fährt sie zu, und verstellt sich als eine Hure, und empfängt von dem Schwäher ohne sein Wissen und Willen. Das hat Gott geschrieben, daß man es wohl ansehe. Wenn es eine weltliche, heidnische Historie wäre, so lautete es schimpflich und hurrisch, aber weil es Gottes ist, soll man es in allen Ehren halten. Wir haben oben [Cap. 27, 1. 23.] auch gehört, wie Isaak, da er alt und blind war, betrogen ward, daß er seines ersten Sohnes fehlte mit dem Segen, und traf Jakob durch Unwissenheit; welches so viel bedeutet, daß, da Gott das Evangelium ließ ausgehen, ging es so seltsam zu, daß die das Evangelium kriegten, dazu man sichs nicht versah, und wiederum meinte man, die großen Canien und Priestertümer, die Klugen und Heiligen sollten es fassen, so fielen sie davon, und traten herzu die Fischer, schlechte, geringe Leute, und predigten es in die Welt. Nicht, daß Gott nicht gewußt oder versehen habe, daß es so gehen sollte, sondern daß es vor unsern Augen so scheint, als geschehe es ohngefähr, denn er stellt sich so seltsam, daß man meint, er wolle dies thun, so thut er ein anderes, und läßt es so geschehen, als müßte er es nicht.

18. Derhalben, daß sie\footnote{„sie“ fehlt in der Originalausgabe, in der Wittenberger und in der Jenaer.} Juda hier so ungefähr antrifft, als er meint, sie sei eine Hure, bedeutet, daß das Evangelium unter dem Volke allein die trifft, die sich für Sünder halten; als Petrus im Schiffe sprach [Luc. 5, 8.]: „HErr, gehe von mir hinaus, denn ich bin ein sündiger
\quellseite{290}
Mensch.“ Und, wie Christus [Matth. 21, 31.] sagt: „Huren und Buben werden euch vorgehen im Himmelreiche.“ Da kommt GOTT, und gibt das Evangelium, und macht die Leute fruchtbar, die zuvor keine Frucht hatten, daß sie nun Zwillinge tragen. Also sind die Juden betrogen; weil es die nicht wollten fassen und annehmen, die da sollten, kommen die andern dazu, die nichts davon wissen, und meinen, es sei ein ander Weib. Also scheint es nun in allen GOTTes Werken, warum, daß uns also däucht und in jedermanns Augen ist, als habe GOTT gefehlt, und das Evangelium denen gegeben, denen es nicht gehört.

19. Weiter spricht der Text: Er habe sie nicht mehr zu sich genommen noch erkannt; das ist, daß Christus der Synagoge hat Urlaub gegeben, daß man ihr Gesetz und Sitten nicht mehr hält, und ist nun frei. Er hat nun zween Söhne heraus gebracht, da läßt er es bei bleiben. Und sie beweist es mit dem Zeichen, so er ihr gegeben hat, daß sie von ihm schwanger worden sei; das ist, wenn man die Schrift ansieht, so wird man gewahr, daß es so sei, und findet die Zeichen, daß es so gehen sollte. So bleibt sie bei Ehren, und beweist, daß GOTT ein solcher Herr ist, der seine Gnade gibt den armen Sünderlein und Sündern, und verwerfe die großen Heiligen.

20. Daß sich aber über der Geburt ein Riß erhebt im Leibe, und der der Letzte war, zum ersten heraus kommt, und wiederum der Erste zum letzten, ist abermal dieselbige Figur. Die Juden sollten die Ersten sein, reichen zum ersten die Hand hervor, als wollten sie es thun, und jene dahinten lassen; aber da es zum Treffen kommt, so bleiben die Werke dahinten, und fährt der Glaube vor; also, daß der andere Sohn nicht mit der Hand allein, sondern mit seinem ganzen Wesen heraus kommt. Darum ist die Summa: Die Gläubigen müssen scheinen, als seien sie die Letzten, und sind doch wahrhaftig die Ersten; wie wir auch droben [Cap. 25, § 37 ff.] in Esau’s und Jacobs Historie genug gesehen haben.

21. Daß nun jener einen rothen Faden um die Hand hatte, ist die Figur, daß alle Werkheiligen blutdürstig und Wütheriche sind, und Feinde derer, die den rechten Glauben haben. Also haben wir in dieser Figur vorgebildet, was GOTT wirkt durch das Evangelium in seinem Volke, wie er aus den Letzten die Ersten macht, und die andern fahren läßt, und daß er die Prediger nicht leiden will, die Gottes Wort auf Werke führen und verderben es, sondern die es lauter lassen bleiben im Glauben, daß es Frucht schaffen kann.

\section{Das neununddreißigste Capitel}

B. 1—23. Joseph ward hinab in Egypten geführt, und Potiphar, ein egyptischer Mann, des Pharao Hofmeister, kaufte ihn von den Ismaeliten, die ihn hinab brachten. Und der HErr war mit Joseph, daß er ein glücklicher Mann ward, und war in seines Herrn, des Egyptiers, Hause. Und sein Herr sahe, daß der HErr mit ihm war; denn alles, was er that, das ließ der HErr glücklich abgehen durch ihn, also, daß er Gnade fand vor seinem Herrn, und sein Diener ward; der setzte ihn über sein Haus, und alles, was er hatte, that er unter seine Hände. Und von der Zeit an, da er ihn über sein Haus und alle seine Güter gesetzt hatte, segnete der HErr des Egyptiers Haus, um Josephs willen, und war eitel Segen des HErrn in allem, was er hatte, zu Hause und zu Felde. Darum ließ er es alles unter Josephs Händen, was er hatte; und er hatte nichts davon, denn nur das Brot, das er aß. Und Joseph war schön und hübsch von Angesicht. Und es begab sich nach dieser Geschichte, daß seines Herrn Weib ihre Augen auf Joseph warf, und sprach: Schlafe bei mir. Er weigerte sichs aber, und sprach zu ihr: Siehe, mein Herr weiß nicht, was im Hause ist, und alles, was er hat, das hat er unter meine Hände gethan, und hat nichts so groß in dem Hause, daß er vor mir verborgen habe, ohne dich, denn du bist sein Weib. Wie sollte ich denn nun ein so groß Uebel thun, und wider GOTT sündigen? Aber sie trieb solche Worte gegen Joseph täglich; aber er gehorchte ihr nicht, daß er nahe\footnote{„nahe“ fehlt in der Erlanger.} bei ihr schliefe, noch um sie wäre. Es
\quellseite{291}
begab sich der Tage einen, daß Joseph in das Haus ging, sein Geschäft zu thun, und war kein Mensch vom Gesinde des Hauses dabei. Und sie erwischte ihn bei seinem Kleide, und sprach: Schlafe bei mir. Aber er ließ das Kleid in ihrer Hand, und floh, und lief zum Hause heraus. Da sie nun sahe, daß er sein Kleid in ihrer Hand ließ, und hinaus entfloh, rief sie das Gesinde im Hause, und sprach zu ihnen: Sehet, er hat uns den hebräischen Mann herein gebracht, daß er uns zu Schanden mache; er kam zu mir herein, und wollte bei mir schlafen; ich rief aber mit lauter Stimme, und da er hörte, daß ich ein Geschrei machte und rief, da ließ er sein Kleid bei mir, und floh, und lief hinaus. Und sie legte sein Kleid neben sich, bis sein Herr heim kam, und sagte zu ihm eben dieselbigen Worte, und sprach: Der hebräische Knecht, den du uns hereingebracht hast, kam zu mir herein, und wollte mich zu Schanden machen. Da ich aber ein Geschrei machte und rief, da ließ er sein Kleid bei mir, und floh hinaus. Als sein Herr hörte die Rede seines Weibes, die sie ihm sagte, und sprach: Also hat mir dein hebräischer Knecht gethan, ward er sehr zornig. Da nahm ihn sein Herr, und legte ihn in das Gefängniß, da des Königs Gefangene innen lagen, und er lag allda im Gefängnisse. Aber der HErr war mit ihm, und neigte seine Huld zu ihm, und ließ ihn Gnade finden vor dem Amtmann über das Gefängniß, daß er ihm unter seine Hand befahl alle Gefangenen im Gefängnisse, auf daß alles, was da geschah, durch ihn geschehen mußte. Denn der Amtmann über das Gefängniß sahe, daß der HErr mit ihm war in allem, das unter seinen Händen war, und daß der HErr glücklich abgehen ließ, was er that.

1. Im sieben und dreißigsten Capitel, V. 4. ff., haben wir gehört, wie Joseph von seinen Brüdern geneidet und gehaßt ward, und den Fremden verkauft, und in Egypten gebracht. Hier kommt nun Moses wieder auf die Historie, wie es ihm in Egypten gegangen ist. Da hangen ein Capitel oder drei an einander, die wollen wir nach der Historie aushandeln, ehe wir die Deutung darinne zeigen. Also ist gesagt, wie er unschuldig von den Brüdern gehaßt; und verkauft ist wie ein Vieh unter fremder Leute Hände; verlassen von aller Freundschaft und Bekannten. Was er da für Elend erlitten habe, ist nicht alles beschrieben, sondern nur ein Stück oder zwei angezeigt, uns zur Lehre und Troste, wie GOtt mit seinen Heiligen so wunderlich fährt.

2. Wir haben gehört, wie GOtt dem frommen Joseph hatte Zusagung gethan, seinen Samen zu mehren wie den Sand am Meere, stellt sich aber je länger je seltsamer dazu, als habe er es nie gedacht noch geredet, nimmt ihm seinen liebsten Sohn von der rechten Mutter, die nun auch todt war, jezund den Joseph, darnach auch Benjamin, die er für die rechten Erben hält, und nicht möglich ist, daß sie ihm nicht sollten lieb sein. Darum ist er je ein wunderlicher GOtt, greift es so über alle Weise und Vernunft an, daß es niemand glauben konnte. Wenn er es hätte zuvor gesagt, wie es kommen sollte, so wäre es noch zu glauben gewesen; nun läßt er ihn nicht anders glauben noch wissen, denn er sei gewürgt und zerrissen, gleichwie, daß er sollte glauben, wie er sollte der oberste Fürst in Egypten werden.

3. Wie muß er ihm denn nun thun? Er muß sich an GOttes Wort halten, und also denken: Ob mir schon alle Söhne todt wären, dennoch bleibt GOtt wahrhaftig, daß mein Same solle werden, wie die Sterne am Himmel, und Sand am Meere. Also hat er müssen stracks wider den Strom fahren, und alle Vernunft dämpfen, wiewohl es ihm nicht süß abgegangen, sondern gar schwer und bitter worden ist, wie der Text auch meldet.

4. Also läßt GOtt seine Heiligen über die Natur fahren, und doch Natur in ihnen bleiben, nämlich, die große Liebe und Neigung zu dem Sohne; greift sie aber an, und versucht, ob er mehr folgen wolle der Neigung (die er doch selbst gegeben hat und nicht verwirft) denn seinem Worte. Es ist alles gut Ding und GOttes Geschäft, noch will er es so mit uns treiben, daß er sehe, und wir gewahr werden, ob wir ihn lieber haben denn seine Güter; und stärkt also die Seinen, daß sie allein am Worte können halten, sollten sie gleich alle seine Güter fahren lassen.

5. Das sage ich abermal darum, daß man nicht Stöcke aus den Heiligen mache, sondern lasse Mensch und Natur bleiben, das sie ist; wie es unserer Natur wehe thut, wenn einer einen guten Freund verliert; item, wenn uns hungert und dürstet, oder wenn wir sterben sollen. Wiederum
\quellseite{292}
thut es uns wohl, wenn wir gesund sind, Ruhe und genug haben. Was können wir dawider, weil es in uns gepflanzt, und alles Gottes ist? Da liegt es aber an, daß man nicht Gottes Wort um deswillen fahren lasse, und mehr an den Creaturen hange denn an ihm selbst.

6. So siehe nun, wie er den Joseph hat wollen zum trefflichen Manne machen, durch welchen Land und Leuten, und vornehmlich seinem Volke, Vater und Brüdern sollte geholfen werden, daß er ein Ausbund unter allen Brüdern würde; aber ehe er dazu kommt, wird seiner so wohl versucht als er, nämlich ganzer dreizehn Jahre, vom siebenzehnten Jahre an bis in das dreißigste. Dadurch muß Gott lehren, wie er der rechte Vater sei, und unsere Vorschläge und Gedanken gar nichts gelten. Denn da reißt er den Vater und Sohn von einander mit großem Herzeleid und Schmerzen. Der Vater ist des Sohnes beraubt, der Sohn kommt in ein Land, da er die Sprache nicht kann, geschweige, daß er einen Freund sollte wissen, zu dem er sich Gutes versähe. Ueber das, als ihn Gott dennoch einen gnädigen Herrn schafft, da er doch kaum gar satt das Brod hat, kommt er in einen andern Jammer und Noth, um der Frau willen, und liegt länger denn zwei Jahre unschuldig gefangen, als er auf das treulichste gedient und seinem Herrn viel Frommen und Nutzen geschafft hatte.

7. Da siehst du beide, Gottes Weisheit und Regiment, und seine edlen, zarten Gaben, die in dem Joseph leuchten. Gottes Weisheit und Sorge darinne, daß er sich so fremd stellt gegen den Vater und Sohn, und versorgt ihn doch also, daß ihn der Vater nicht so könnte versorgen mit allem Vermögen, und stellt ihn dem Vater wieder unter Augen, in solcher Herrlichkeit, daß er es nimmer also hätte dürfen wünschen. Ist es nicht wunderlich, er wird so schändlich verkauft im siebenzehnten Jahre, dazu gefangen gelegt, und soll über dreizehn Jahre der Nächste des größesten Königs in Egypten werden, als er nun von allen Brüdern, und schier auch vom Vater vergessen ist? Wenn wir schon keine Historie mehr hätten, sollte man je daraus lernen erkennen, was Gott für ein Regiment führt in der Welt, so wunderlich, daß, wo man meint, es sei der Teufel und Tod, da ist er am nächsten.

8. Er meint, er sei verlassen von Gott und von der Welt, so wartet sein Gott, und hat ein Auge auf ihn. [Er] läßt ihn wohl verkauft und gefangen werden, als sei kein Gott bei ihm, aber da die Zeit kommt, setzt er ihn zum höchsten Ehren; daß wohl von ihm redet der weise Mann im Buche der Weisheit [Cap. 10, 13. 14.]: Haec venditum justum non dereliquit, die Weisheit, das ist, Gottes Wort, hat den frommen Joseph, als er verkauft war, nicht gelassen, und hat ihn errettet aus der Untugend, und war bei ihm in der Grube und Banden, bis sie ihm brachte das königliche Scepter, zog mit ihm in der Feinde Hände, bis in Kerker, so lange bis er heraus kam, und seine Weisheit hervor bracht, daß er Herr über Egypten ward.

9. Also haben die Väter diese Historie hoch angesehen, und sich verwundert über der Gottes Weisheit, als auch David im 105. Psalm, V. 16—21., sagt: „Er ließ eine Theurung in das Land kommen, und entzog allen Vorrath des Brodes. Er sandte einen Mann vor ihnen hin, Joseph ward zum Knechte verkauft. Sie zwangen seine Füße in einen Stock, Eisen ging durch seine Seele, bis zur Zeit, daß\footnote{Jenaer: da.} sein Wort kam, und die Rede des HErrn ihn durchläuterte. Da sandte der König hin, und ließ ihn losgeben, der\footnote{der" fehlt in der Erlanger.} Herr über Völker hieß ihn auslassen. Er setzte ihn zum Herrn über sein Haus, zum Herrscher über alle seine Güter.“ Es war verborgen, was Gott mit ihm ausrichten wollte; daß er aber im Sinne sollte haben, durch sein Elend so viel Land und Leute zu erretten, das gedachte niemand. Aber Gott gab ihm Weisheit in das Herz, dadurch er ein mächtiger Herr ward, der viel Leuten konnte nütze sein. Wäre er daheim geblieben, so wäre er\footnote{Erlanger: es.} nichts mehr denn ein Hirte geblieben.

10. So rechne nun gegen einander, er sei elend, und allein im fremden Lande, so sind die elf Brüder daheim, und meinen, sie stehen gar wohl, und er müßte sein Lebetage ein elender, gefangener Knecht bleiben, so fährt es Gott um, und macht einen solchen Herrn aus ihm, wie ihm geträumt hatte, daß ihm nicht allein seine Brüder müssen unterthan sein und zu Fuße fallen, sondern das ganze Land und alles Volk, ausgenommen der König. Siehe, das thut es, wer Gott vertrauen kann und ihm ausharren. Es ist nichts mit Vater und Mutter, daß sie uns
\quellseite{293}
helfen sollten (wiewohl man ihnen soll gehorsam sein); aber der rechte Vater ist allein GOtt, der aus allen Nöthen und Elende zur höchsten Ehren hilft. Also, daß dies Exempel ein recht Vorbild und große Neigung ist zum Glauben. Wo ich hinkomme, in des Vaters Haus oder in das Elend und fremde Länder, daß ich wisse, GOtt sei auch da daheim, sonderlich wo ich keinen Trost und Zuversicht zu den Leuten haben kann, und ganz verlassen bin.

11. Ich wollte auch gerne, wenn mir es heimgeteilt wäre, bei einem solchen Vater sein und bleiben, wie Jakob war, daß ich könnte sagen: Ich bin je bei einem frommen Vater, der Gottes Wort hat, da der Heilige Geist wohnt mit Gottes Fülle und Segen; aber da stößt GOtt den Sohn aus, und führt ihn hinweg in das Land, da nur der Teufel wohnt und regiert, und nichts von GOtt ist, daß er dazu Leib und Leben muß wagen. Das muß ihm mächtig wehe gethan haben, sonderlich weil er so jung vom Vater kam. Doch hat er von ihm gelernt und gefaßt das Wort, das er predigt, wie ihm GOtt verheißen hatte, daß Christus von ihm kommen sollte; dennoch muß er von ihm und dencken: Ich werde dennoch den GOtt nicht verlieren, der hier ist, er wird anderswo auch bei mir sein; wiewohl es ihm wehe gethan hat. Denn Fleisch und Blut wird sich auch geregt haben, daß ihm manchmal vor Elend die Augen übergegangen sind.

12. Solches sollten wir doch auch einmal lernen, wenn der Unglaube nicht das Herzleid und alles Unglück hätte. Was hatte der gute Joseph, als er weggeführt ward? Sie haben ihm das Kleid ausgezogen, und nicht einen Heller gelassen, verkaufen ihn in solch fernes Land, da er mit schwerer Arbeit und lange Zeit nicht mehr verdiente denn Hülle und Fülle, und dem Herrn so groß frommete. So läßt es GOtt gehen, und ihn eine lange Zeit Kummer leiden, doch nicht Hungers sterben. Aber da die Zeit aus ist, gibt er ihm dagegen das ganze Land in die Hand, daß sie alle müssen von ihm gespeist werden, und er so viel Macht kriegt, daß alles, was er thut, gethan ist. Das hat er alles mit dem Anhalten\footnote{„Anhalten“ fehlt in der Wittenberger und Jenaer.} erlangt, daß er GOtt im Glauben und Geduld ausgehalten hat.\footnote{Die Worte: „daß er ... ausgehalten hat“ fehlen in der Erlanger.} Also thäte GOtt noch bei uns, wenn wir so könnten im Glauben beharren. Es ist kein anderer GOtt jetzund, denn der daselbst regierte, wir haben eben denselben allmächtigen Vater, und eben denselben Wort, daß er uns nicht lassen wolle. Das ist nun ein Stück der Anfechtung, die er überwunden hat.

13. Das andere, als er nun zu Gnaden gekommen war, und [es] ihm ein wenig wohl ging, und [er] gesetzt wird über seines Herrn Haus (wiewohl er nicht viel mehr kriegt, denn das Brod), schlägt eine neue Anfechtung her auf der rechten Seite. Da siehe, welch einen trefflichen Geist er hat, und eine seltsame, hohe Tugend der Jungfrauschaft, daß ihn die Frau im Hause, des Königs Hofmeisters, lieb gewinnt, und anhebt, bei ihr zu schlafen. Was hätte er nicht von ihr können zuwege bringen und erwerben, daß sie ihm heimlich geschenkt hätte und herrliche gute Tage gemacht? Dazu sagt der Text: Sie habe es ihm nicht Ein Mal angeboten, sondern täglich getrieben.

14. Das ist ja eine große Tugend, daß der, der so großen Raum, Zeit, Stätte, Person, und dazu Anreizung hat, und konnte Gut und Gnade gewinnen, dennoch sich enthält; dazu ist er ein junger Geselle, und frisches Geblütes, daß der Natur nicht möglich wäre, wider solch täglich Reizen und Geilen\footnote{Geilen = anhaltend bitten.} zu bestehen, weil die Jugend an ihr selbst hitzig ist.

15. Aber er gibt ihr nichts anders zur Antwort, denn das: „Mein Herr weiß nicht, was im Hause ist, und alles, was er hat, das hat er unter meine Hände gethan, und hat nichts so Großes im Hause, das er vor mir verhohlen habe, ohne dich.“ Das muß ja ein treuer Knecht gewesen sein, der alle des Herrn Gut versorgt, so treulich, daß der Herr ihm alles machen läßt, und nicht darnach fragt, was er im Hause habe, möchte dazu das Weib auch schänden, daß er nichts darum müßte; noch sagt er: Du bist Frau, ich bin Knecht, und meinem Herrn Treue schuldig, die will ich halten; und zeucht es noch weiter: „Wie sollte ich ein solch groß Uebel thun, und wider GOtt sündigen?“ Siehe, welch ein großer Geist, wie hoch setzt er GOttes Gebot über alles, das auf Erden ist, und ihm nichts läßt so lieb sein, dawider zu thun, wiewohl er jung ist und so angehetzt wird.
\quellseite{294}

16. Wo sind nun unsere Mönche und Nonnen, die sich ihrer Keuschheit rühmen? Siehe, ob sie nicht dagegen alle zu Schanden werden, weil er so großen Raum und Fug hat, die Frau im Hause, die ihm täglich anliegt und treibt, welches einem jungen Menschen eine schwere und fährliche Anfechtung ist: noch bleibt er keusch bis in das dreißigste Jahr, daß er kein Weib berührt, hat allein Gott vor Augen, und die Treue seines Herrn. [Es] ist ein frommer, aufgerichter, redlicher Geist in ihm, der doch so große Ursache hat zur Buhlerei, und ohne Unterlaß getrieben wird; welches auch möchte Eisen zerschmelzen, wie St. Hieronymus sagt, und bald versehen wäre, daß einer zu Falle käme, ob er gleich nichts Böses Willens hätte, und lange Keuschheit gehalten. Aber er hütet sich gleichwohl, daß er nicht um sie wäre, noch neben ihr im Hause, Gemach, oder Kammer läge. Denn das gehört auch dazu, daß man desto baß die Keuschheit erhalte, wie St. Paulus lehrt, als er [1 Cor. 6, 18.] sagt: „Fliehet die Hurerei.“

17. Es ist nicht leicht zu überwinden, man fliehe denn weit davon; ob du dich gleich fast fastest und fest hältst, ist es doch gefährlich, wenn Mann und Weib bei einander sind; denn Fleisch und Blut bleibt Fleisch und Blut. Darum ist nichts Sichereres, denn ferne von einander, oder [man] bleibe ewig bei einander, sonst geht es schwerlich rein ab. Darum hat er nicht bei ihr wollen sein, noch mit ihr umgehen, daß er nichts von ihr sähe noch hörete. Was geschah aber? Sie lauert darauf, und erwischt ihn einmal, als niemand dabei war. Denn es ist eine jämmerliche Seuche, je mehr man ihr will wehren, je hitziger sie wird. Das ist nun dem Knaben zu nahe, und ein harter Stoß gewesen; aber der Geist wird muthig und stärkt in ihm, daß er es überwindet, und als er nicht anders kann davon kommen, läßt er seinen Mantel dahinten. Der Heilige Geist braucht nicht umsonst so viel Worte darüber, weiß wohl, daß es ein seltsam, groß Exempel ist; denn Fleisch und Blut ist müthend und unsinnig in dieser Anfechtung, sonderlich wenn sie bei einander sind.

18. Da er nun also hinweggegangen und ihr entflöh, und sie sieht, daß er schlecht nicht will, da wird sie rasend und tobend, und denkt ihn um das Leben zu bringen. So geht es zu, das Hurenübel kann nicht gebüßt werden, es komme denn ein ander Uebel darein; wie der Poet sagt: Cedit\footnote{Erlanger: Severit.} amor furiis [Wuth tritt an die Stelle der Liebe]. Weil sie sieht, daß sie verachtet wird, und kann ihren Muthwillen nicht erlangen, kehrt sich der Sinn um, und wird toll und töricht. Das sind die Früchte Fleisches und Blutes, wenn ihnen Gott nicht steuert und wehrt. Vorhin wollte sie ihn vor Liebe fressen, jetzt will sie ihn vor Zorn erwürgen; denkt also: Er will sich verachten, so hast du ihn nun angegriffen, das wird er nachsagen und dich zu Schanden machen, daß du nach ihm geilet,\footnote{geilen nach jemand — ihm in unzüchtiger Weise nachstehen. (Diese Bedeutung fehlt bei Tieck.)} und wird dir lohnen wie einer Hure. Darum greift sie ihn auf der andern Seite an, wird ihm spinnefeind, und bringt diese Klage vor den Herrn: So hast du mich lieb, das meinst du mit mir, hast mir den hebräischen Knecht herein gebracht, daß er mich zu Schanden mache.

19. Das ist ein recht teuflisch Stücklein. Sie glaubt nicht, daß Gott auch bei ihm wäre, meint, sie wolle es schmücken und bergen, und ihn um Ehre, Leib und Leben bringen. Gott schweigt aber still, und läßt den Knaben ohne Schuld hinführen in das Gefängniß, auf das allerdümmlichste, als einen Ehebrecher; sie aber bezeuget mit Ehern, niemand entschuldiget oder hilfet ihm. Sollte es ihm nicht wehe gethan haben, daß ihm so schändlich Unrecht geschah, daß ihm das Herz hätte mögen weinen? Aber er schweigt, und befiehlt es Gott. Doch schlug ihm oft der Jammer unter Augen, daß [es] nicht genug war, daß er vom Vater kommen war, sondern müßte noch weiter in das Elend, und in der argen Hure willen unschuldiglich umkommen. So tief läßt Gott seine Heiligen stecken, wenn er sie will hoch heben.

20. O! wer die Krempel könnte fassen, daß er so gnädig und gut ist, wenn er uns so greulich angreift. Joseph fühlt es auch noch nicht, hübsch aber, Gott werde ihn nicht lassen. Aber wie man es ansieht, so ist es ein unfreundlich Spiel, doch im Grunde so große Gnade, daß er hernach so zu Ehren kommt. Wer weiß auch, wie es der Hure gegangen ist; der Heilige Geist schreibt es nicht, denn er hat nicht Lust und Freude am Unglück. Das zeigt er aber an, daß [es] Gott am allerbesten meint, wenn er uns auf das allerschändlichste zurichtet. Das kann die Welt nicht glauben, denn es ist wider alle
\quellseite{295}
Sinne. Er versucht aber allein den Glauben, ob er festhalten könne, darnach hilft er zukünftig so stark, daß [es] unmöglich wäre einem Menschen zu gedenken.

21. Denn wenn Joseph sollte gegeben sein zu wünschen, hätte er also gewünscht, daß er möchte aus dem Gefängnisse, darinne er unschuldig gelegt war, und wieder heim zum Vater kommen. O! wie ein geringe, schmal und schwach Gebet war das vor Gott, daß er auch darum noch zwei Jahre sitzen muß; denn er wollte ihm viel Größeres geben. Daß er aber ein solcher Herr sollte werden in Egypten, und dazu Vater und Mutter wieder kriegen, hätte er nimmermehr können denken. Das ist, das Paulus [Röm. 8, 26.] sagt: „Wir wissen nicht, was wir bitten sollen, wie es sich gebühret, sondern der Geist vertritt uns selbst mächtiglich mit unaussprechlichem Seufzen.“ Wenn uns Gott gäbe, was wir bitten, das hinderte nur seine reiche Güte, als, daß er Joseph nicht mehr sollte gegeben haben, denn er gebeten hatte, wäre viel zu geringe gewesen, daß er nicht hätte sein Wunder beweisen können. Darum sollen wir im Elende Gott Raum, Zeit und Stätte geben, und nicht [be]stimmen, wie wir es gerne hätten. Denn wir sind Narren, und wissen nicht, was wir bitten sollen. Wenn wir aber stille hielten, und ihm heim gäben nach seinem Gefallen zu machen, so erführen wir, wie überaus reichlich er vergelten kann. Aber die Natur kann es nicht, so thut er es auch nicht, wenn er dich lieb hat; hat er dich aber nicht lieb, so gibt er dir wohl, es ist dir aber nicht gut.

22. Das läßt er uns predigen und schreiben, daß man seine Art und Weise einmal lerne kennen. Wie der 4. Psalm, V. 4., sagt: „Erkennet doch, daß der Herr seine Heiligen wunderlich führet.“ Also sollte er sagen: Wahr ist es, wenn man zu ihm ruft, so erhört er, so oft wir schreien und bitten; aber doch macht er es so wunderlich, daß [es] niemand begreifen kann; wie Joseph gerne aus dem Elende los wäre, schreiet und fleht um Hülfe; wird auch so bald erhört. Aber wie wunderlich geht es zu: Er hofft bald heraus zu kommen, [es] verzeucht sich aber noch wohl zwei Jahre lang. Ist das erhört? Ja, es heißt wunderlich erhört. Es gilt aber nicht, daß man ihm vorschlage, sondern er will Meister sein, und es also machen, daß man sich vor Wunder segne, und spreche: Hätte ich doch nimmermehr geglaubt, daß es also sollte kommen; also, daß du vor Freuden springest, und müssest bekennen, du seiest nicht allein erhört, sondern wunderlich erhört. Das wäre noth, daß wir es einmal sähen, weil die Zeit alle Tage vorhanden ist, da wir es wohl bedürften, daß man nicht zweifele, wir seien gewiß erhört; aber wie es zugehen soll, und er es hinaus führen wolle, soll ihm allein bekannt sein, daß er es also machen könne, daß wir sagen müssen, Gott sei gelobet, daß es nicht gegangen ist nach unserem Sinne.

\section{Das vierzigste Capitel}

V. 1—23. Und es begab sich darnach, daß sich versündigten der Amtmann über die Schenken des Königes zu Egypten, und der Amtmann über die Bäcker, an ihrem Herrn, dem Könige zu Egypten. Und Pharao ward zornig über sie, und ließ sie setzen in des Hofmeisters Haus, in das Gefängniß, da Joseph gefangen lag. Und der Hofmeister setzte Joseph über sie, daß er ihnen dienete, und saßen etliche Tage im Gefängnisse. Und es träumte ihnen beiden, dem Schenken und Bäcker, in Einer Nacht, einem jeglichen ein eigener Traum, und eines jeglichen Traum hatte seine Bedeutung. Da nun des Morgens Joseph zu ihnen hinein kam, und sahe, daß sie saur sahen, fragte er sie, und sprach: Warum seid ihr heute so übel? Sie antworteten: Es hat uns geträumet, und haben niemand, der es uns auslege. Joseph sprach: Auslegen gehört Gott zu; doch erzählet mir es. Da erzählte der oberste Schenke seinen Traum Joseph, und sprach zu ihm: Mir hat geträumet, daß ein Weinstock vor mir wäre, der hatte drei Reben, und er grünete, wuchs und blühete, und seine Trauben wurden reif; und ich hatte den Becher Pharao in meiner Hand, und nahm die Beeren und zerdrückte sie in den Becher, und gab den Becher Pharao in die Hand. Joseph
\quellseite{296}
sprach zu ihm: Das ist seine Bedeutung: Drei Adern sind drei Tage. Aber drei Tage wird Pharao dein Haupt erheben, und dich wieder an dein Amt stellen, daß du ihm den Becher in die Hand gebest nach der vorigen Weise, da du sein Schenke wärest. Aber gedenke meiner, wenn es dir wohl geht, und thue Barmherzigkeit an mir, daß du Pharao erinnerst, daß er mich aus diesem Hause führe. Denn ich bin aus dem Lande der Hebräer heimlich gestohlen, dazu habe ich auch allhier nichts gethan, daß sie mich eingesetzt haben. Da der oberste Bäcker sahe, daß die Deutung gut war, sprach er zu Joseph: Mir hat auch\footnote{„auch“ fehlt in der Erlanger.} geträumet, ich trüge drei geflochtene Körbe auf meinem Haupte, und im obersten Korbe allerlei gebackene Speise dem Pharao, und die Vögel äßen aus dem Korbe auf meinem Haupte. Joseph antwortete und sprach: Das ist die Deutung: Drei Körbe sind drei Tage, und nach dreien Tagen wird dir Pharao deinen Kopf erheben und dich an den Galgen henken, und die Vögel werden dein Fleisch von dir essen. Und es geschah des dritten Tages, da beging Pharao seinen Jahrstag, und er machte eine Mahlzeit allen seinen Knechten, und erhob das Haupt des obersten Schenken, und das Haupt des obersten Bäckers unter seinen Knechten, und setzte den obersten Schenken wieder zu seinem Schenkamte, daß er den Becher reichte in Pharao Hand. Aber den obersten Bäcker ließ er henken, wie ihnen Joseph gedeutet hatte. Aber der oberste Schenke gedachte nicht an Joseph, sondern vergaß seiner.

1. Da will nun Moses schreiben, wie sich’s begeben hat, daß Joseph wieder los worden ist, doch also, daß ihn Gott noch zwei Jahre ließ gefangen bleiben. Hier ist nun etwas von Träumen zu sagen, weil Joseph zwei Träume deutete, und dadurch genoß, daß er los ward.

2. Aber zum ersten ist hier wohl zu sehen, daß in dem Lande nicht so große Blindheit gewesen ist, daß sie nichts sollten von Gott gewußt haben. Als auch im vorigen Capitel, B. 3., angezeigt ist, daß der Potiphar, Josephs Herr, Gott auch erkannt habe, denn er klar also sagt: „Er sahe, daß der Herr mit ihm war, und ihn um seinetwillen segnete“ etc. Das muß ja ein groß, hoch Erkenntniß gewesen sein, denn es ist nicht ein geringer Verstand. Wer das weiß, daß Gott läßt sein Gut machen und zunehmen, und den Segen gibt, der weiß auch, daß Gott ein Schöpfer Himmels und Erden ist, und alle Dinge regiert. Sonderlich muß es etwas gewesen sein, weil die Schrift nicht meldet, daß sie Abgötterei gehabt und getrieben haben; wiewohl es mag mit unter gelaufen sein. Wie wir bisher so viel Erkenntniß behalten haben, daß ein Gott ist, der Himmel und Erden geschaffen hat, und alles gebe, was da wächst; wiewohl dasselbige Erkenntniß daneben verblendet ist mit mancherlei Abgötterei und falschem Vertrauen und Secten, durch welche jenes verfinstert ist, daß wir nicht auf der rechten Bahn geblieben sind.

3. Denn wiewohl wir geglaubt haben, daß es der rechte Gott sei, der alle Dinge schaffe, hat es doch daran gefehlt, daß wir nicht recht angetroffen haben, wenn wir sagen: Hier ist Gott. Denn wenn man ihn so will heften und anbinden an Zeit, Person und Stätte, und ihm eine Farbe anmalen, daß er so und so gesinnt sei, wie wir meinen, so hat man sein gefehlt; als, daß man ihm sonst oder so dienen, dieses oder jenes Werk thun will, das ihm gefallen solle; wie das ganze Geschwürm unter dem Pabstthum von Orden, Regeln und Menschenlehren gewesen ist. Da kommt der Teufel an Gottes Statt; sonst bleibt gleichwohl das gemeine Erkenntniß von dem rechten Gott.

4. Denn das ist auch wahr, daß St. Paulus [Röm. 1, 19.] sagt, daß Gott aller Welt offenbart hat, daß sie müssen wissen, daß ein Gott sei, ist unverborgen gewesen von Anfang der Welt bis ans Ende. Welches man auch eben dabei merkt, daß, wenn die Heiden nicht Wissen hätten von Gott gehabt, so hätten sie ihre Götzen nicht Gott geheißen. Daß ein Gott müßte sein, das wußten sie, und es war recht; aber wenn sie sagten: Das ist Gott, so fehlten sie sobald. Sie könnten aber nicht fehlen, wenn sie nicht wüßten, daß Gott wäre. Da gehört nun ein höheres Erkenntniß zu, daß man ihn recht treffe und nicht fehle. Denn wo man ihn recht treffen soll, da muß sein Wort sein, das uns lehre; unsere Vernunft und Klugheit wird es langsam finden.

5. Nun müssen wir hier die Träume sehen. Droben [Cap. 37, 7. 9.] haben wir auch gehört von Josephs Träumen, die er seinen Brüdern
\quellseite{297}
Predigt über 1 Mose 40, 1—23.


sagte, darinnen er deutete die Historie, die da folgen wird, wie ihm seine Brüder zu Fuße fallen und unterthan sollen sein. So habe ich sonst gesagt, wie man den Träumen glauben solle; aber hier siehst du ein Meisterstück, so der Heilige Geist anzeigt, wie es nicht gewiß ist, auf Träume sich verlassen, weil hier zwei gleiche Träume beschrieben werden, daß Vernunft müßte sagen, sie müßten auch einerlei Deutung haben, so sie doch Joseph gar wider einander deutet. Denn wie in des Schenken Traume drei Trauben drei Tage deuten, und er Pharao den Becher in die Hand gibt, deutet, daß er wieder an das Schenkamt kommen sollte, also sollte auch des Bäckers Traum, daß er drei Körbe mit gebackenem Brode dem Pharao vorträgt, so viel deuten, daß er nach dreien Tagen an sein Amt gesetzt würde. Nun kommt der Geist darein, und deutet gleich das Wiederspiel.

6. Warum ist nun das geschrieben? Eben darum, daß man auf Träume nicht baue. Denn ich habe gesagt: Es mag geschehen, daß Gott einem Menschen einen Traum gebe, der wahr sei, und also geschehe, wie der Traum klingt. Als man liest von vielen Heiligen, und er selbst sagt im vierten Buche Mosis [Cap. 12, 6.]: „Ist jemand ein Prophet des HErrn, dem will ich mich kund machen in einem Gesichte, oder in einem Traume will ich mit dir reden.“ Und im Joel [Cap. 3, 1.]: „Eure Söhne und eure Töchter sollen weissagen, und eure Jünglinge sollen Gesichte sehen, und eure Aeltesten sollen Träume träumen.“

7. Aber hier steht die Glosse im Text, wie es zugehen soll, nämlich, als Joseph spricht: „Auslegen gehört Gott zu.“ Träume hin, Träume her, Auslegen gebührt dir nicht; wem er einen Traum gibt, und mit dem er redet, gibt er auch daneben einen gewissen Verstand zu deuten, wie er soll. Aber damit ist aufgehoben die Vermessenheit, daß niemand auch rechte, göttliche Träume selbst auslege. Laß träumen, was man will; ist es ein rechter Traum, den Gott haben will, wird er es wohl auslegen, entweder durch sich, oder durch einen Menschen, also, daß du es im Herzen fühlest, daß [es] ein rechter Traum sei; thut er es ohne Mittel, so wird er machen, daß dein Gewissen gefangen wird, daß es recht sei.

8. Also schreibt man von einem heidnischen Könige, Attila, als der Bischof Paulinus gefangen und sein Gärtner war, hatte er des Nachts einen Traum, wie er stünde vor Gerichte, und drei Männer ihn verurtheilten. Danach ließ er ohngefahr Paulinum vor sich kommen, und als er ihn ansah, erschrak er, und sprach: Bei Nacht habe ich einen Traum gehabt, daß ihrer drei über mich urtheilten, und siehe, dieser war der eine. Da ward er so blöde und verzagt, daß er zurück fiel, und bald darnach starb. Siehe, das war ein rechter Traum, wiewohl er es selbst nicht wußte; denn es ging ihm durch das Herz, daß er es fühlte als einen Strahl.

9. Dergleichen liest man auch von Monica, St. Augustini Mutter, und andern mehr. Also, daß wenn sie recht sind, sind sie so, daß Gott selbst auslegt, ohne Mittel, oder durch Mittel, daß es das Gewissen wohl fühlt. Darum ist gar großer Unterschied unter Träumen und Auslegen. Lege du nur keinen selbst aus, denn wenn du zweifelst und meinst, daß es auf das Dünken kommt, so laß nur davon. Entweder der Traum ist nicht recht, oder Gott hat ihn nicht ausgelegt. Wenn er auslegt, so wirst du bald fühlen, ob er recht sei. Also deutet Joseph hier den Gefangenen gewiß und ohne Wanken, daß sie es fühlen und annehmen.

10. Warum haben die weit geirrt, die aus diesem Texte haben genommen, Bücher zu schreiben und Träume zu deuten: als der Narr, der ein Buch gemacht hat de somniis Danielis; haben etliche Regeln darauf gegeben, was dies oder jenes deute. Aber es ist umsonst. Es läßt sich [nicht\footnote{„nicht“ fehlt in den Ausgaben, doch hat es die Jenaer als Conjectur am Rande.}] fassen, ist alles falsch, dazu auch verboten im Gesetze Mosis [5 Mos. 13, 3.]: „Du sollst nicht achten auf die Träume“, das ist, du sollst dich\footnote{„dich“ fehlt in der Erlanger.} nicht unterstehen, die Träume gewiß zu deuten, sondern was rechte Träume sind, die soll Gott selbst auslegen. Falsche Träume sind gleich wie falsche Lehre, kriegen wohl einen Schein, daß man meint, es sei etwas; [es] ist aber ungewiß und betrüglich. Aber Gottes Wort macht das Herz gewiß. Also ist es auch mit dem Glauben und falschem Gottesdienste. Jener macht gewiß, dieser\footnote{In den Ausgaben: dies.} bleibt ungewiß. Darum laß Träume Träume bleiben, wenn sie Gott nicht auslegt.

11. So merkte nun, weil Joseph die zwei
\quellseite{298}
gleichen Träume so ungleich deutet (wiewohl der Bäcker meinte, weil er des Schenken Traum so wohl auslegte, seine Deutung sollte auch also klingen), wenn man in göttlichen Sachen handelt, daß da nicht gilt nachahmen. Denn das treibt der Teufel immerdar von Anfang der Welt, daß er will Gottes Affe sein, und ihm nachthun, was er von ihm sieht. Als, da Gott [2 Mos. 25, 3—9.] geboten hat, einen Tabernakel zu bauen, und daselbst [zu] opfern und räuchern, fuhr er hinnach, und wo er einen hübschen, grünen Wald, Garten oder Berg sahe, da richtete er auch ein Opfern und Räuchern an [Jer. 3, 6.], daß sein Gottesdienst mehr mußte gelten, denn den Gott geboten hatte. Also, zu unsern Zeiten, weil Christus [Luc. 18, 1. 21, 36.] befohlen hat, ohne Unterlaß zu bitten, ist der Teufel da, und stiftet alle Winkel voll Kirchen, da man Tag und Nacht plärrt und heult; und des Affenspiels unzählig viel mehr. Item, man findet etliche, die dem heiligen Kreuze nachahmen; weil Christus [Matth. 16, 24.] geboten hat, man soll ihm selbst absagen, und sein Kreuz tragen, fahren sie zu, und suchen ihnen selbst ein Leiden, und legen ihnen selbst viel Marter auf, und wollen groß Wunderwerk treiben, verderben darüber Leib und Seele.

12. Solches Nachahmens ist die Welt allenthalben voll, daß keine größere Verführung und Gotteslästerung zu finden ist. Solchem zuwider läßt Gott allhier den einen anlaufen mit seiner Hoffnung, daß er meint, es soll eben so viel bedeuten, und behält ihm die Deutung allein. Träume wie du willst, so wird es niemand recht deuten, denn der. Ehe denn sie geschehen, magst du wohl ahnen,\footnote{ahnen = Vermuthungen aufstellen.} aber deute daneben, es sei nichts. Gott thut wohl etlichen etwas kund durch Träume, aber es steht bei uns nicht, dieselbigen zu verstehen; wenn wir es selbst deuten, so fehlt es gewißlich. Darum hüte dich vor ahnen und nachspielen.

13. Also haben wir nun ein Stück von Josephs Historie, wie er, um der Keuschheit willen, von der Frau in Angst und Noth gebracht ward, und in Kerker geworfen, und Gott doch daselbst bei ihm war, und gab ihm einen kleinen Trost, daß er Gnade und Gunst kriegte bei dem obersten Meister, daß er ihn nicht in Stock legte, sondern den Gefangenen dienen und pflegen ließ.

14. Da haben wir gesehen, wie gar reich er gewesen ist im Glauben und Gottes Worte, daß er so beständig bleibt in der Keuschheit, so er doch einen hübschen Deckel der Büberei hätte mögen haben, [es] zu treiben, wie er gewollt hätte, und viel Geschenke dazu gehabt. Welches gar eine edle Tugend und hoher Geist gewesen ist, der nicht allein seines Fleisches mächtig war, sondern auch des Weibes und des Teufels selbst. Dazu ihn Gott erleuchtet hat, und einen prophetischen Geist gegeben, Träume auszulegen; dadurch er auch dazu kommt, daß ihn Pharao empor hebt, also, daß allenthalben angezeigt wird, wie Gott seine Heiligen nicht läßt, ob er wohl ein wenig die Hand zurück zieht und läßt sie verstoßen und unterdrückt werden.

\section{Das einundvierzigste Capitel}

V. 1—57. Und nach zwei Jahren hatte Pharao einen Traum, wie er stände am Wasser, und sähe aus dem Wasser steigen sieben feine Rinder und feist von Leibe, und gingen an der Weide auf der Wiese. Nach diesen sähe er andere sieben Rinder aus dem Wasser aufsteigen, die waren ungestalt und mager von Leibe, und traten neben die Rinder an das Ufer am Wasser, und die ungestalten und magern Rinder fraßen die schönen und fetten Rinder. Da erwachte Pharao. Und er schlief wieder ein, und ihm träumete abermal, und sähe, daß sieben Aehren wuchsen aus Einem Halme, voll und gut; darnach sähe er sieben dünne und versengte Aehren aufgehen, und die sieben mageren Aehren verschlangen die sieben großen und vollen Aehren. Da erwachte Pharao, und sähe, daß es ein Traum war. Und da es Morgen ward, war sein Geist betrübt, und schickte aus und ließ rufen alle Wahrsager in Egypten, und alle Weisen, und erzählte ihnen seinen Traum. Aber da war keiner, der sie dem Pharao deuten konnte. Da redete der oberste Schenke zu Pharao, und sprach:
\quellseite{299}
Predigt über 1 Mose 41, 1—57.


Ich gedachte heute an meine Sünde; da Pharao zornig ward über seine Knechte, und mich mit dem obersten Bäcker in das Gefängniß legte, in des Hofmeisters Hause, da träumete uns beiden in Einer Nacht, einem jeglichen sein Traum, des Deutung ihn betraf. Da war bei uns ein hebräischer Jüngling, des Hofmeisters Knecht, dem erzählten wir es, und er deutete uns unsere Träume, einem jeglichen nach seinem Traume. Und wie er uns deutete, so ist es ergangen; denn ich bin wieder an mein Amt gesetzt, und jener ist gehenkt. Da sandte Pharao hin und ließ Joseph rufen, und ließen ihn eilend aus dem Loch. Und er ließ sich bescheren, und wandelte seine Kleider und kam hinein zu Pharao. Da sprach Pharao zu ihm: Mir hat ein Traum geträumet, und ist niemand, der ihn deute. Ich habe aber gehört von dir sagen, wenn du einen Traum hörest, daß du ihn deutest. Joseph antwortete Pharao und sprach: Gott wird Pharao Glück sagen lassen, auch wohl ohne mich. Pharao sagte an zu Joseph: Mir träumete, ich stund am Ufer bei dem Wasser, und sahe aus dem Wasser steigen sieben Kinder, feist vom Leibe und fein, und gingen an der Weide auf der Wiese. Und nach ihnen sahe ich andere sieben Kinder heraus steigen, dünne und fast ungestalt, mager vom Leibe. Ich habe in ganz Egyptenland nicht solche ungestalte gesehen. Und die sieben mageren und ungestalten Kinder fraßen auf die sieben ersten feisten Kinder. Und da sie die hinein gefressen hatten, merkte man es nicht an ihnen, daß sie die gefressen hatten, und waren ungestalt gleichwie vorhin. Da wachte ich auf. Und sahe abermal in meinem Traume sieben Aehren auf Einem Halme wachsen, voll und gut. Darnach gingen auf sieben dürre Aehren, dünne und versenget. Und die sieben dünnen Aehren verschlangen die sieben guten Aehren. Und ich habe es meinen Wahrsagern gesagt, aber sie sagen mir nichts davon. Joseph antwortete Pharao: Beide Träume Pharao sind einerlei, Gott verkündigt Pharao, was er thut. Die sieben guten Kinder sind sieben Jahre, und die sieben guten Aehren sind auch die sieben Jahre. Es ist einerlei Traum. Die sieben mageren Kinder und ungestalt, die nach jenen aufgestiegen sind, das sind sieben Jahre, und die sieben mageren Aehren und versengten, das werden sein sieben Jahre Theurung. Das ist nun, daß ich gesagt habe zu Pharao, daß Gott Pharao zeigt, was er thut. Siehe, sieben Jahre werden kommen mit großer Fülle in ganz Egyptenlande. Und nach denselben werden sieben Jahre Theurung kommen, daß man vergessen wird aller solcher Fülle in Egyptenlande. Und die Theurung wird das Land verzehren, daß man nichts wissen wird von der Fülle im Lande, vor der theuren Zeit, die hernach kommt; denn sie wird fast schwer sein. Daß aber dem Pharao zum andernmal geträumet hat, bedeutet, daß solch Ding von Gott gefertigt\footnote{fertigen = zurichten, bereiten, fördern. In der Bibel: daß solches Gott gewißlich und eilend thun wird.} wird und Gott dasselbe eilend thun wird. Nun sehe Pharao nach einem verständigen und weisen Manne, den er über Egyptenland setze, und schaffe, daß er Amtleute verordne im Lande, und nehme den Fünften in Egyptenlande, in den sieben reichen Jahren, und sammle alle Speise der guten Jahre, die kommen werden, daß sie Getreide aufschichten unter Pharao Gewalt, zur Nahrung, in den Städten, und verwahren es, auf daß man Speise verordnet finde dem Lande in den sieben theuren Jahren, die über Egyptenland kommen werden, daß nicht das Land vor Hunger verderbe. Die Rede gefiel Pharao und allen seinen Knechten wohl. Und Pharao sprach zu seinen Knechten: Wie könnten wir einen solchen Mann finden, in dem der Geist Gottes sei? Und sprach zu Joseph: Weil dir Gott solches alles hat kund gethan, ist keiner so verständig und weise als du. Du sollst über mein Haus sein, und nach deinem Wort soll alle mein Volk sich nähren. Allein des königlichen Stuhls will ich mehr sein denn du. Und sprach: Siehe, ich habe dich über ganz Egyptenland gesetzt. Und that seinen Ring von seiner Hand, und gab ihn Joseph an seine Hand, und kleidete ihn mit weißer Seide, und hing ihm eine güldene Kette an seinen Hals, und ließ ihn auf dem andern Wagen fahren, und ließ vor ihm ausrufen, daß man die Knie vor ihm beugen sollte, als den er gesetzt hätte über ganz Egyptenland. Und Pharao sprach zu Joseph: Ich bin Pharao, ohne deinen Willen soll niemand seine Hand oder seinen Fuß regen in ganz Egyptenland. Und nannte ihn Zaphnath Panea.\footnote{Handlosse: Zaphnath Panea ist egyptisch geredet, und noch unbewußt, was es sei, ohne daß, so viel man spüren kann, heißt es, wie man auf deutsch spricht: Der heimliche nächste Rath.} Und gab ihm ein Weib, Asnath, die Tochter Potiphera
\quellseite{300}
, des Priesters zu On. Also zog Joseph aus, das Land Egypten zu besehen. Und er war dreißig Jahre alt, da er vor Pharao stand, dem Könige in Egypten, und fuhr aus von Pharao, und zog durch ganz Egyptenland. Und das Land that also die sieben reichen Jahre, und sammelten alle Speise der sieben Jahre, so im Lande Egypten waren, und thaten sie in die Städte. Was für Speise auf dem Felde einer jeglichen Stadt umher wuchs, das thaten sie hinein. Also schüttete Joseph das Getreide auf, über die Maße viel, wie Sand am Meere, also, daß er aufhörete zu zählen, denn man konnte es nicht zählen. Und Joseph wurden zwei Söhne geboren, ehe denn die Theurung kam, welche ihm gebar Asnath, Potiphera, des Priesters zu On, Tochter. Und hieß den ersten Manasse; denn Gott, sprach er, hat mich lassen vergessen aller meiner Arbeit, und alle meines Vaters Hauses. Den andern hieß er Ephraim; denn Gott, sprach er, hat mich lassen wachsen in dem Lande meines Elendes. Da nun die sieben reichen Jahre um waren im Lande Egypten, da fingen an die sieben theuren Jahre zu kommen, da Joseph von gesagt hatte. Und es ward eine Theurung in allen Landen, aber in ganz Egyptenlande war Brod. Da nun das Egyptenland auch Hunger litt, schrie das Volk zu Pharao um Brod. Aber Pharao sprach zu allen Egypten: Gehet hin zu Joseph, was euch der sagt, das thut. Als nun über alles Land Theurung war, that Joseph auf alles, was bei ihm war, und verkaufte den Egypten. Also nahm die Theurung überhand im Lande. Und alle Lande kamen in Egypten, zu kaufen bei Joseph; denn die Theurung nahm überhand in allen Landen.

1. Was ich von den zwei Träumen, des Schenken und Bäckers, gesagt habe, soll man auch hier in des Pharao Traum verstehen. Als nun die Zeit kommt, daß Joseph wieder hervor soll, schickt es Gott, ohne sein Bedenken und unversehens, daß er plötzlich empor kommt, und weiß nicht wie, auf daß man je sehe, daß alle unsere Dinge bloß in Gottes Händen stehen, der schafft und wirkt es, wie er will, kommt immer zuvor, ehe wir es gedenken, daß es alles seine Güte und Gnade thut. Joseph, das heilige Gottes-Kind, kommt so schändlich vom Vater, aber Gott ist bei ihm; hätte nie gedacht, daß es dahin sollte kommen, daß er einen guten, gnädigen Herrn kriegte, der ihm alles in die Hände gibt, was er hat im Hause und Hofe. Aber als ihn die böse Hure in das Gefängniß ohne sein Versehen bringt, und [er] nun wiederum empor soll, weiß er abermal nichts davon: daß also Gottes Rath uns immerdar zuvor kommt; wir sind aber Narren, können uns nicht darein richten, und müssen es doch ihm heimgeben. Wie dieser Joseph thut, so thut er wiederum wie ein frommer Gott und Vater, dessen Werke und Rathschläge eitel Heil, Glück und Gnade ist, wiewohl es viel anders scheint. Als auch hier, als er im allertiefsten Unglück steckt, kommt er bald in das höchste Glück. Denn er meinte nicht anders, denn wäre nun gar vergessen, dieweil der oberste Schenke sein vergessen hatte, und nun zwei ganze Jahre nach ihm gesessen war.

2. Man pflegt oft zu fragen: wie weit Glück und Unglück von einander sei? Da muß alle Vernunft sagen, sie wisse es nicht; ist auch wahr, man kann es nicht wissen, es muß geglaubt sein. Denn da liegt Joseph, als der\footnote{Erlanger: er.} nun ewig sollte gefangen bleiben als ein Ehebrecher, aber vor Gott ist bereits beschlossen, daß sich alle Kniee vor ihm beugen müssen. Also sieht der Glaube wohl, wie nahe es bei einander ist. Gleichwie das Leben und Tod in einander geschlossen sind, daß das Leben im Tode ist, daß [es] nicht näher könnte sein, ja, bereits schon Ein Ding ist: also ist wiederum Unglück im größten Glück, in Armuth und Elend Reichthum, Lust und Freude. Wiederum, im sichersten Leben ist in einem Augenblicke der Tod. So, wenn wir sterben sollen, gehen wir\footnote{„wir“ fehlt in der Erlanger.} hin in den Tod, als sollten wir ewig todt liegen: kommt uns der jüngste Tag in einem Hui, daß wir sagen: Nun will ich ewig leben. Solches kann der Glaube fassen, daß Gutes und Böses, Leben und Tod gar in einander geknüpft ist.

3. Das zeigt uns Gott so durch mancherlei Exempel, wie er der Herr sei, der enthinter\footnote{enhinter (enhinder) = zurück.} wirkt und empor hebt, und alles gehen muß nach seinem Rathe, auf daß wir nichts von uns vornehmen zu thun, sondern alles ihm heimgeben zu regieren, davon fast alle Psalmen Davids singen, als [Ps. 94, 11.]: „Der Herr weiß die Gedanken der Menschen, daß sie nichts sind“; wie auch ein jeglicher in seinem Leben erfährt
\quellseite{301}
und bekennen muß, daß nichts nach unserm Rathe und Vorschläge hinaus geht. So folgt nun, wie Gott den Joseph erhöht hat, und dem Könige gleich gemacht, durch die Auslegung seiner Träume.

B. 37. ff. Die Rede gefiel Pharao und allen seinen Knechten wohl. Und Pharao sprach zu seinen Knechten: Wie könnten wir einen solchen Mann finden, in dem der Geist Gottes sei? Und sprach zu Joseph: Weil dir Gott solches alles hat kund gethan etc.

4. Da sieht man aber ein Stück von Gottes Werken. Die Juden meinten, Gott wäre allein bei ihnen; aber da zeigt er an, daß er regiert und versorgt die ganze Welt, wie Joseph selbst aus dem Geiste bekennt: „Gott zeigt dem Pharao, was er thun will.“ Darum muß dieser Pharao nicht ein schlechter Mann sein vor Gott, sondern rechten Verstand und Glauben von Gott gehabt haben, weil er ihm solches selbst offenbart, und sonst keinem. Denn, daß er Gott erkannt habe, sagt der Text klar, weil er ihn in dem Joseph so ehrt und bekennt, daß er Gottes Geist hat; welches kein Ungläubiger gethan hätte.

5. Also haben wir auch droben [Cap. 20, § 21] in Abrahams Historie gehört, von dem Könige Abimelech, daß der Leute viel gewesen sind, die rechten Glauben und Verstand gehabt haben, ob sie wohl nicht aus der Schnur und Linie Abrahams gewesen sind, welchem wiewohl es zugejagt ist, sind doch andere Leute nicht ausgeschlossen, wie Paulus Röm. 3, 29. sagt: „Ist er nicht auch der Heiden Gott? Ja freilich auch der Heiden Gott.“ Denn wenn man auch der Heiden Historien ansieht, als der Römer, sieht man sein, wie Gott gehandelt und regiert hat, ob sie es gleich nicht verstanden haben.

6. Also hält der Pharao Joseph in Ehren, daß man sieht, wie er ein großer Mann ist, darum ihm Gott auch einen solchen Mann gibt. Denn wenn Gott einem Lande helfen will und rathen, muß er Leute dazu geben. So muß nun dies ganze Land des Joseph genießen, also, daß es ein Ausbund vor allen Landen worden ist, welches auch in der Theurung andere Lande konnte speisen, die verdorben waren. So hat Gott sonderlich auf dies Land ein Auge gehabt, und groß begnadet vor allen andern, um Josephs und seines Vaters Jakob willen, weil er sie dahin schicken wollte, daß vielen Leuten auch geistlich geholfen würde, Gottes Wort zu lernen. Das sage ich darum, daß man wisse, wie Gott wunderlich regiert, nimmt die ganze Welt an, daß man in allen Winkeln finde, die ihn angehören.

7. So sagt nun Moses, wie ihn Pharao über das ganze Land gesetzt hat, und ließ gebieten, daß man die Kniee vor ihm beugen sollte im ganzen Lande, und nennt ihn auf Egyptisch צפנת פענח,\footnote{In der Erlanger: Zaphnat Paacenea; in der Wittenberger und der Jenaer: Paenea.} dafür in unserer lateinischen Bibel steht: Salvatorem mundi, als sollte es heißen, ein Heiland der Welt. Ich halte aber, daß irgend ein Jude, der St. Hieronymum unterwiesen hat, so mit ihm gescherzt habe, wie es ihm gefallen hat. Ich halte, er wird ihn geheißen haben seinen Kanzler und obersten, heimlichsten Rath, darum, daß er ihm den Fingerreif von seiner Hand gibt.

8. Dazu gibt er ihm ein Weib, des höchsten Priesters Tochter. Es möchte aber wohl jemand verdrießen, daß er nicht ein Weib von seinem Geblüte nimmt, sondern befreiet\footnote{befreien = sich verheirathen. — Erlanger: befreiet.} sich mit den Heiden. Aber der Priester Potiphera wird ein großer, tapferer Mann gewesen sein. Denn die Königreiche, die so von Gott begnadet sind gewesen, sind von seinen Leuten versehen gewesen. Sie haben auch Priester und Lehrer in Ehren gehalten; wie wir [Cap. 47, 22. 26.] hören werden, daß Joseph, als er das Land beschwerte mit Zinsen, den Fürsten zu geben, ließ er die Priester frei. Darum müssen sie nicht schlechte Leute gewesen sein; ob sie gleich Heiden gewesen sind, haben sie doch etwas einen rechten Verstand gehabt. Also hat der König nach der besten Meße gegriffen, auf daß er Joseph zu gröbern Ehren setzte.

9. Auch war es nicht verboten, aus den Heiden zu freien; denn David, Salomon und andere haben es auch gethan; allein, daß sie es nicht zu viel machten, und durch sie verführet würden; Joseph würde sich sonst auch nicht hinein begeben haben. [Er] ist mit ihnen bekannt worden, und hat mit ihnen ihr Ding geredet und gerathen, und den Mangel gebessert, wie der 105. Psalm, B. 22., sagt: „daß er seine Fürsten züchtigt nach seinem Willen, und seine Aeltesten weise gemacht“; also, daß er beide, weltlich
\quellseite{302}
und geistlich Regiment, in rechten Schwang gebracht hat. Darum ist es auch aus sonderlicher Ordnung Gottes geschehen, daß er das Weib gefrieget hat, daß er sie und seine Kinder unterweiste, und den ganzen Königshof. Er ist zum Meister gesetzt, der da strafte, was straflich war, und alle Dinge besserte, als denn wohl noth ist, wenn gleich das Wort gepredigt wird.

10. Weiter schreibt Moses, wie er zween Söhne zeugt,\footnote{Erlanger: kriegt.} von welchen die Schrift viel redet. Droben hat er die zwölf Patriarchen nach einander erzählt, also, daß Joseph der erste, und Benjamin der zwölfte war; aber nun hernach nimmt er die zwei und macht sie zu Hauptsöhnen, also, daß von Einem Vater, Joseph, zwei Stämme kämen.

11. Den ersten Sohn nennt er Manasse, darum, daß ihn Gott\footnote{Erlanger: hat.} seines Vaters Haus, und alles Elendes hätte vergessen lassen, denn das Wort Manasse heißt vergessen. Wie? ist denn das christlich, daß er sich rühmt, er habe Vater und Mutter vergessen? [Er] will also sagen: Ich sehe, daß mir Gott die Zuversicht hat wollen nehmen, die ich zu meinem Vater habe; denn Gott ist ein Eiferer, will nicht leiden, daß das Herz einen Boden habe, darauf es sich verlasse und stütze, denn allein auf ihn; darum muß ich mich deß erwägen, wenn ich gleich nimmer zum Vater komme, und immer allein bin; ich habe einen Gott, wenn ich den habe, so habe ich alles.

12. Den andern Sohn heißt er Ephraim, das ist, gemachen; als sollte er sagen: Ich bin aus des Vaters Hause geschüpft, habe kein Erbtheil, rein abgetrocknet und versiegen; aber dagegen hat mich Gott hier im Elende wachsen lassen, mehr denn alle meine Freunde.

13. Auch muß man wissen, wie die zwei Lande, Egypten und Canaan, an einander stoßen, und rühren mit einem Winkel zusammen. Canaan liegt hinein gegen Mitternacht, Egypten gegen den Abend, daß sie mit einander grenzen und zusammen treffen, wie zwei Zipfel. Darum ist es nicht eine ferne Reise gewesen, als sie hinein zogen, und Getreide gekauft haben; wie folgt.

\section{Das zweiundvierzigste Capitel}

\begin{quote}
V. 1—38. Da aber Jakob sahe, daß Getreide in Egypten feil war, sprach er zu seinen Söhnen: Warum sehet ihr zu? Siehe, ich höre, es sei in Egypten Getreide feil; ziehet hinab und kaufet uns Getreide, daß wir leben und nicht sterben. Also zogen hinab zehn Brüder Josephs, daß sie in Egypten Getreide kauften, denn Benjamin, Josephs Bruder, ließ Jakob nicht mit seinen Brüdern ziehen; denn er sprach: Es möchte ihm ein Unfall begegnen. Also kamen die Kinder Israel, Getreide zu kaufen, sammt andern, die mit ihnen kamen; denn es war im Lande Canaan auch Theurung. Aber Joseph war der Regent im Lande, und verkaufte Getreide allem Volke im Lande. Da nun seine Brüder zu ihm kamen, fielen sie vor ihm nieder zur Erde auf ihr Antlitz. Und er sahe sie an, und kannte sie, und stellte sich fremde gegen sie, und redete hart mit ihnen, und sprach zu ihnen: Woher kommt ihr? Sie sprachen: Aus dem Lande Canaan, Speise zu kaufen. Aber wiewohl er sie kannte, kannten sie ihn doch nicht. Und Joseph gedachte an die Träume, die ihm von ihnen geträumet hatten, und sprach zu ihnen: Ihr seid Kundschafter, und seid gekommen zu sehen, wo das Land offen ist. Sie antworteten ihm: Nein, mein Herr, deine Knechte sind gekommen, Speise zu kaufen. Wir sind alle eines Mannes Söhne, wir sind redlich, und deine Knechte sind nie Kundschafter gewesen. Er sprach zu ihnen: Nein, sondern ihr seid kommen, zu besehen, wo das Land offen ist. Sie antworteten ihm: Wir, deine Knechte, sind zwölf Brüder, Eines Mannes Söhne, im Lande Canaan, und der jüngste ist noch bei unserm Vater; aber der eine ist nicht mehr vorhanden. Joseph sprach zu ihnen: Das ist es, daß ich euch gesagt habe, Kundschafter seid ihr; darum will ich euch prüfen; bei dem Leben Pharaonis, ihr sollt nicht von dannen kommen, es komme denn her euer jüngster Bruder. Sendet einen unter euch hin, der euren Bruder hole; ihr
\end{quote}
\quellseite{303}
aber sollt gefangen sein. Also will ich prüfen eure Rede, ob ihr mit Wahrheit umgehet, oder nicht. Denn wo nicht, so seid ihr, bei dem Leben Pharaonis, Kundschafter. Und er that sie zusammen in eine Verwahrung drei Tage lang. Am dritten Tage aber sprach er zu ihnen: Wollt ihr leben, so thut also, denn ich fürchte Gott. Seid ihr redlich, so lasset eurer Brüder einen gebunden liegen in eurem Gefängnisse, ihr aber ziehet hin, und bringet die Waare der Nothdurft zu Hause, und bringet euren jüngsten Bruder zu mir, so will ich euren Worten glauben, daß ihr nicht sterben müsset. Und sie thaten also. Sie aber sprachen unter einander: Das haben wir an unserm Bruder verschuldet, daß wir sahen die Angst seiner Seele, da er uns flehete, und wir wollten ihn nicht erhören; darum kommt nun dies Trübsal über uns. Ruben antwortete ihnen, und sprach: Sagte ich es euch nicht, da ich sprach: Versündiget euch nicht an dem Knaben, und ihr wolltet nicht hören? nun wird sein Blut gefordert. Sie wußten aber nicht, daß es Joseph verstand, denn er redete mit ihnen durch einen Dolmetscher. Und er wandte sich von ihnen und weinete. Da er nun sich wieder zu ihnen wandte, und mit ihnen redete, nahm er aus ihnen Simeon, und band ihn vor ihren Augen, und that Befehl, daß man ihre Säcke mit Getreide füllete, und ihr Geld wieder gäbe, einem jeglichen in seinen Sack; dazu auch jeglichem seine Zehrung auf die Reise. Und man that ihnen also. Und sie luden ihre Waare auf ihre Esel, und zogen von dannen. Da aber einer seinen Sack aufthat, daß er seinem Esel Futter gäbe in der Herberge, ward er gewahr seines Geldes, das oben im Sacke lag, und sprach zu seinen Brüdern: Mein Geld ist mir wieder worden, siehe, in meinem Sacke ist es. Da entfiel ihnen ihr Herz, und erschraken unter einander, und sprachen: Warum hat uns Gott das gethan? Da sie nun heim kamen zu ihrem Vater Jakob in das Land Canaan, sagten sie ihm alles, was ihnen begegnet war, und sprachen: Der Mann, der im Lande Herr ist, redete hart mit uns, und hielt uns für Kundschafter des Landes. Und da wir ihm antworteten: Wir sind redlich, und nie Kundschafter gewesen, sondern zwölf Brüder, unsers Vaters Söhne; einer ist nicht mehr vorhanden, und der jüngste ist noch heutiges Tages bei unserem Vater im Lande Canaan, sprach er: Daran will ich merken, daß ihr redlich seid: Einen eurer Brüder lasset bei mir, und nehmet die Nothdurft für euer Haus, und ziehet hin, und bringet euren jüngsten Bruder zu mir; so merke ich, daß ihr nicht Kundschafter, sondern redlich seid; so will ich euch auch euren Bruder geben, und möget im Lande werben. Und da sie die Säcke ausschütteten, fand ein jeglicher sein Bündlein Geldes in seinem Sacke. Und da sie sahen, daß es Bündlein ihres Geldes waren, sammt ihrem Vater erschraken sie. Da sprach Jakob, ihr Vater, zu ihnen: Ihr beraubet mich meiner Kinder; Joseph ist nicht mehr vorhanden, Simeon ist nicht mehr vorhanden, Benjamin wollt ihr hinnehmen; es geht alles über mich. Ruben antwortete seinem Vater und sprach: Wenn ich dir ihn nicht wiederbringe, so erwürge meine zwei Söhne. Gib ihn mir in meine Hand, ich will ihn dir wieder bringen. Er sprach: Mein Sohn soll nicht mit euch hinab ziehen, denn sein Bruder ist todt, und er ist allein überblieben. Wenn ihm ein Unfall auf dem Wege begegnete, da ihr auf reiset, würdet ihr meine grauen Haare mit Schmerzen zur Hölle hinunter treiben.

\section{Das dreiundvierzigste Capitel}

V. 1—34. Die Theurung aber drückte das Land. Und da es gar alle worden, was sie für Waare aus Egypten gebracht hatten, sprach ihr Vater Jakob zu ihnen: Ziehet wieder hin, und kaufet uns ein wenig Speise. Da antwortete ihm Juda und sprach: Der Mann band uns das hart ein, und sprach: Ihr sollt mein Angesicht nicht sehen, es sei denn euer Bruder mit euch. Ist es nun, daß du unsern Bruder mit uns sendest, so wollen wir hinab ziehen, und dir zu essen kaufen. Ist es aber, daß du ihn nicht sendest, so ziehen wir nicht hinab; denn der Mann hat gesagt zu uns: Ihr sollt mein Angesicht nicht sehen, euer Bruder sei denn mit euch. Israel sprach:
\quellseite{304}
Warum habt ihr so übel an mir gethan, daß ihr dem Manne gesagt, wie ihr noch einen Bruder habt? Sie antworteten: Der Mann forschte so genau nach uns und unserer Freundschaft, und sprach: Lebt euer Vater noch? habt ihr auch noch einen Bruder? Da sagten wir ihm, wie er uns fragte. Wie konnten wir so eben wissen, daß er sagen würde: Bringet euren Bruder mit hernieder? Da sprach Juda zu Israel, seinem Vater: Laß den Knaben mit mir ziehen, daß wir uns aufmachen und reisen, und leben, und nicht sterben, beide, wir und du, und unsere Kindlein. Ich will Bürge für ihn sein, von meinen Händen sollst du ihn wieder fordern. Wenn ich dir ihn nicht wiederbringe, und vor deine Augen stelle, so will ich mein Lebelang die Schuld tragen. Denn wo wir nicht hätten verzogen, wären wir schon wohl zweimal wieder gekommen. Da sprach Israel, ihr Vater, zu ihnen: Muß es denn ja also sein, so thut es. Und nehmet von des Landes besten Früchten in eure Säcke, und bringet dem Manne Geschenke hinab, ein wenig Balsam, und Honig, und Würze, und Myrrhen, und Datteln, und Mandeln.\footnote{¹) Randglosse: Diese Namen der Früchte sind noch bisher ungewiß, auch bei den Juden selbst.} Nehmet auch ander Geld mit euch, und das Geld, das euch oben in euren Säcken wieder worden ist, bringet auch wieder mit euch; vielleicht ist ein Irrthum da geschehen. Dazu nehmet euren Bruder, machet euch auf und kommet wieder zu dem Manne. Aber der allmächtige Gott gebe euch Barmherzigkeit vor dem Manne, daß er euch lasse euren andern Bruder und Benjamin. Ich aber muß sein wie einer, der seiner Kinder gar beraubt ist. Da nahmen sie diese Geschenke und ander Geld mit sich, und Benjamin, machten sich auf, zogen in Egypten, und traten vor Joseph. Da sahe sie Joseph mit Benjamin, und sprach zu dem, der über sein Haus war: Führe diese Männer zu Hause, und schlachte und richte zu, denn sie sollen zu Mittage mit mir essen. Und der Mann that, wie ihm Joseph gesagt hatte, und führte die Männer in Josephs Haus. Sie fürchteten sich aber, daß sie in Josephs Haus geführt wurden, und sprachen: Wir sind herein geführt um des Geldes willen, das wir in unsern Säcken vorhin wieder gefunden haben, daß er es auf uns bringe, und fälle ein Urtheil über uns, damit er uns nehme zu eigenen Knechten sammt unsern Eseln. Darum traten sie zu dem Manne, der über Josephs Haus war, und redeten mit ihm vor der Hausthüre, und sprachen: Mein Herr, wir sind vorhin herab gezogen, Speise zu kaufen, und da wir in die Herberg kamen, und unsere Säcke aufthaten, siehe, da war eines jeglichen Geld oben in seinem Sacke, mit völligem Gewichte; darum haben wir es wieder mit uns gebracht, haben auch ander Geld mit uns herab gebracht, Speise zu kaufen. Wir wissen aber nicht, wer uns unser Geld in unsere Säcke gesteckt hat. Er aber sprach: Gehabt euch wohl, fürchtet euch nicht, euer Gott, und eures Vaters Gott, hat euch einen Schatz gegeben in eure Säcke. Euer Geld ist mir worden. Und er führte Simeon zu ihnen heraus, und führte sie in Josephs Haus, gab ihnen Wasser, daß sie ihre Füße wuschen, und gab ihren Eseln Futter. Sie aber bereiteten das Geschenk zu, bis daß Joseph kam auf den Mittag; denn sie hatten gehört, daß sie daselbst das Brod essen sollten. Da nun Joseph zum Hause einging, brachten sie ihm zu Hause das Geschenk in ihren Händen, und fielen vor ihm nieder zur Erde. Er aber grüßte sie freundlich und sprach: Geht es eurem Vater, dem Alten, wohl, von dem ihr mir saget? Lebt er noch? Sie antworteten: Es geht deinem Knechte, unserm Vater, wohl, und lebt noch, und neigten sich und fielen vor ihm nieder. Und er hub seine Augen auf, und sahe seinen Bruder Benjamin, seiner Mutter Sohn, und sprach: Ist das euer jüngster Bruder, da ihr mir von saget? Und sprach weiter: Gott sei dir gnädig, mein lieber Sohn. Und Joseph eilte, denn seines Herzens Grund entbrannte ihm gegen seinem Bruder, und suchte, wo er weinete, und ging in sein Kämmerlein, und weinete daselbst. Und da er sein Angesicht gewaschen hatte, ging er heraus, und hielt sich fest, und sprach: Leget Brod auf. Und man trug ihm besonders auf, und jenen auch besonders, und den Egyptern auch besonders. Denn die Egypter dürfen nicht Brod essen mit den Hebräern; denn es ist ein Greuel vor ihnen. Und man setzte sie gegen ihn, den Erstgebornen nach seiner Erstgeburt, und den Jüngsten nach seiner Jugend. Da verwunderten sie sich unter einander. Und man trug ihnen Bescheiden vor von seinem Tische, aber dem Benjamin ward fünfmal mehr, denn den andern. Und sie tranken, und wurden trunken mit ihm.
\quellseite{305}

\section{Das vierundvierzigste Capitel}

V. 1—34. Und Joseph befahl dem, der über sein Haus war, und sprach: Fülle den Männern ihre Säcke mit Speise, so viel sie führen mögen, und lege jeglichem sein Geld oben in seinem Sack. Und meinen silbernen Becher lege oben in des Jüngsten Sack, mit dem Gelde für die Waare. Der that, wie ihm Joseph hatte gesagt. Des Morgens, da es licht ward, ließen sie die Männer ziehen mit ihren Eseln. Da sie aber zur Stadt hinaus waren, und nicht ferne gekommen, sprach Joseph zu dem, der über sein Haus war: Auf, und jage den Männern nach, und wenn du sie ergreifst, so sprich zu ihnen: Warum habt ihr Gutes mit Bösem vergolten? Ist es nicht das, daraus mein Herr trinkt, und damit er weissagt? Es ist übel gethan, das ihr gethan habt. Und als er sie ergriff, redete er mit ihnen solche Worte. Sie antworteten ihm: Warum redet mein Herr solche Worte? es sei ferne von deinen Knechten, ein solch Ding zu thun. Siehe, das Geld, das wir gefunden oben in unsern Säcken, haben wir wieder gebracht zu dir aus dem Lande Canaan. Und wie sollten wir denn aus deines Herrn Hause gestohlen haben Silber oder Gold? Bei welchem er gefunden wird unter deinen Knechten, der sei des Todes. Dazu wollen wir auch meines Herrn Knechte sein. Er sprach: Ja, es sei, wie ihr geredet habt; bei welchem er gefunden wird, der sei mein Knecht, ihr aber sollt ledig sein. Und sie eilten, und ein jeglicher legte seinen Sack ab auf die Erde, und ein jeglicher that seinen Sack auf. Und er suchte, und hob am Größesten an bis auf den Jüngsten. Da fand sich der Becher in Benjamins Sacke. Da zerrissen sie ihre Kleider, und lud ein jeglicher auf seinen Esel, und zogen wieder in die Stadt. Und Juda ging mit seinen Brüdern in Josephs Haus; denn er war noch daselbst, und sie fielen vor ihm nieder auf die Erde. Joseph aber sprach zu ihnen: Was ist das für eine That, die ihr gethan habt? Wisset ihr nicht, daß ein solcher Mann, wie ich bin, errathen könne? Juda sprach: Was sollen wir sagen meinem Herrn, oder wie sollen wir reden? Und was sollen wir vorwenden? Gott hat die Missethat deiner Knechte gefunden. Siehe da, wir und der, bei dem der Becher gefunden ist, sind meines Herrn Knechte. Er aber sprach: Das sei ferne von mir, solches zu thun. Der Mann, bei dem der Becher gefunden ist, soll mein Knecht sein, ihr aber ziehet hinauf mit Frieden zu eurem Vater. Da trat Juda zu ihm, und sprach: Mein Herr, laß deinen Knecht ein Wort reden vor deinen Ohren, mein Herr, und dein Zorn ergrimme nicht über deinen Knecht, denn du bist wie Pharao. Mein Herr fragte seine Knechte, und sprach: Habt ihr auch einen Vater oder Bruder? Da antworteten wir: Wir haben einen Vater, der ist alt, und einen jungen Knaben, in seinem Alter gezeuget, und sein Bruder ist todt, und er ist allein übergeblieben von seiner Mutter, und sein Vater hat ihn lieb. Da sprachest du [zu deinen Knechten]: Bringet ihn herab zu mir, so will ich meine Augen auf ihn werfen. Wir aber antworteten meinem Herrn: Der Knabe kann nicht von seinem Vater kommen; wo er von ihm käme, würde er sterben. Da sprachest du zu deinen Knechten: Wo euer jüngster Bruder nicht mit euch herkommt, sollt ihr mein Angesicht nicht mehr sehen. Da zogen wir hinauf zu deinem Knechte, meinem Vater, und sagten ihm an meines Herrn Rede. Da sprach unser Vater: Ziehet wieder hin und kaufet uns ein wenig Speise. Wir aber sprachen: Wir können nicht hinab ziehen, es sei denn unser jüngster Bruder mit uns, so wollen wir hinab ziehen; denn wir können des Mannes Angesicht nicht sehen, wo unser jüngster Bruder nicht mit uns ist. Da sprach dein Knecht, mein Vater, zu uns: Ihr wisset, daß mir mein Weib zwei Söhne geboren hat; einer ging hinaus von mir, und man sagt, er ist zerrissen, und habe ihn nicht gesehen bisher. Werdet ihr diesen auch von mir nehmen, und ihm ein Unfall widerfährt, so werdet ihr meine grauen Haare mit Jammer hinunter in die Hölle führen. Nun, so ich heim käme zu deinem Knechte, meinem Vater, und der Knabe wäre nicht mit mir, weil seine Seele an dieses Seele hanget, so wird es geschehen, wenn er sieht, daß der Knabe nicht da ist, daß er stirbt; so würden wir, deine Knechte,
\quellseite{306}
die grauen Haare deines Knechtes, unseres Vaters, mit Jammer hinunter in die Hölle bringen. Denn ich, dein Knecht, bin Bürge worden für den Knaben gegen meinem Vater, und sprach: Bringe ich ihn dir nicht wieder, so will ich mein Lebenlang die Schuld tragen. Darum laß deinen Knecht hier bleiben an des Knaben Statt, zum Knechte meines Herrn, und den Knaben mit seinen Brüdern hinauf ziehen. Denn wie soll ich hinauf ziehen zu meinem Vater, wenn der Knabe nicht mit mir ist? Ich würde den Jammer sehen müssen, der meinem Vater begegnen würde.

\section{Das fünfundvierzigste Capitel}

V. 1—28. Da konnte sich Joseph nicht enthalten vor allen, die umher standen, und er rief jedermann: Gehet von mir hinaus; und stand kein Mensch bei ihm, da sich Joseph mit seinen Brüdern bekannte. Und er weinete laut, daß es die Egypter und das Gesinde Pharao hörten, und sprach zu seinen Brüdern: Ich bin Joseph, lebt mein Vater noch? Und seine Brüder konnten ihm nicht antworten, so erschraken sie vor seinem Angesichte. Er sprach aber zu seinen Brüdern: Tretet doch her zu mir.\footnote{Randglosse: Zu mir. Das sind die süßen Worte des Evangelii; also redet Christus mit der Seele im Glauben, nachdem sie durch das Gewissen der Sünde wohl gedemüthigt und geängstet ist.} Und sie traten herzu. Und er sprach: Ich bin Joseph, euer Bruder, den ihr in Egypten verkauft habt. Und nun bekümmert euch nicht, und denket nicht, daß ich zorn sei,\footnote{In der Bibel: daß ich zürne.} daß ihr mich hieher verkauft habt, denn um eures Lebens willen hat mich GOtt vor euch her gesandt. Denn dies sind zwei Jahre, daß es theuer im Lande ist, und sind noch fünf Jahre, daß kein Pflügen noch kein Ernten sein wird. Aber GOtt hat mich vor euch her gesandt, daß er euch übrig bleiben lasse auf Erden, und euer Leben erhalte, durch eine große Errettung. Und nun, ihr habt mich nicht hergesandt, sondern GOtt, der hat mich Pharao zum Vater gesetzt, und zum Herrn über alle sein Haus, und einen Fürsten in ganz Egyptenland. Eilet nun und ziehet hinauf zu meinem Vater, und saget ihm: Das läßt dir Joseph, dein Sohn, sagen: GOtt hat mich zum Herrn in ganz Egypten gesetzt, komm herab zu mir, säume dich nicht. Du sollst im Lande Gosen wohnen, und nahe bei mir sein, du und deine Kinder und deine Kindeskinder, deine Schafe und Rinder, und alles, was dein ist, ich will dich daselbst versorgen. Denn es sind noch fünf Jahre der Theuerung, auf daß du nicht verderbest mit deinem Hause, und allem, das dein ist. Siehe, eure Augen sehen, und die Augen meines Bruders Benjamin, daß ich mündlich mit euch rede. Verkündiget meinem Vater alle meine Herrlichkeit in Egypten, und alles, was ihr gesehen habt; eilet und kommet hernieder mit meinem Vater hieher. Und er fiel seinem Bruder Benjamin um den Hals, und weinete, und Benjamin weinete auch an seinem Halse. Und küßte alle seine Brüder, und weinete über sie. Darnach redten seine Brüder mit ihm. Und da das Geschrei kam in Pharao's Haus, daß Josephs Brüder gekommen waren, gefiel es Pharao wohl, und allen seinen Knechten. Und Pharao sprach zu Joseph: Sage deinen Brüdern: Thut ihm also, beladet eure Thiere, ziehet hin, und wenn ihr kommt in das Land Canaan, so nehmet euren Vater, und euer Gesinde, und kommet zu mir, ich will euch Güter geben in Egyptenland, daß ihr essen sollet das Mark im Lande. Und gebet ihnen: Thut ihm also, nehmet zu euch aus Egyptenland Wagen zu euren Kindern und Weibern, und führet euren Vater, und kommet. Und schonet nicht eures Hausraths, denn die Güter des ganzen Landes Egypten sollen euer sein. Die Kinder Israel thaten also. Und Joseph gab ihnen Wagen, nach dem Befehle Pharao, und Zehrung auf den Weg, und gab ihnen allen, einem jeglichen ein Feierkleid; aber Benjamin gab er dreihundert Silberlinge und fünf Feierkleider. Aber seinem Vater sandte er dabei zehn Esel, mit Gut aus Egypten beladen, und zehn Eselinnen mit Getreide, und Brod und Speise seinem Vater auf den Weg. Also ließ er seine Brüder [und sie
\quellseite{307}
Predigt über 1 Mose 45, 1—28.


zogen hin], und sprach zu ihnen: Zanckt nicht auf dem Wege. Also zogen sie hinauf von Egypten, und kamen in das Land Canaan zu ihrem Vater Jakob. Und verkündigten ihm und sprachen: Dein Sohn Joseph lebt noch, und ist ein Herr im ganzen Egyptenland. Aber sein Herz schlug es in [den] Wind, denn er glaubte ihnen nicht. Da sagten sie ihm alle Worte Josephs, die er zu ihnen gesagt hatte. Und da er sahe die Wagen, die ihm Joseph gesandt hatte, ihn zu führen, ward sein Geist lebendig. Und [Israel] sprach: Ich habe genug, daß mein Sohn Joseph noch lebt; ich will hin und ihn sehen, ehe ich sterbe.

1. Diese vier Capitel hangen alle an einander, darum wollen wir sie auch zusammen überlaufen. Bisher haben wir gehört, wie der Joseph durch das Wort Gottes zu großen Ehren gekommen ist, daß er ein Landesherr worden, und Egypten durch ihn errettet ist; und nicht allein Egypten, sondern auch viel andere Leute; damit Gott zeugt und beweist, wie er der rechte Vater und Herr ist, der sich der ganzen Welt annehme. Also ist die Historie an ihr selbst lieblichen, denn niemand sagen und herausstreichen kann; darum mag es jedermann bei sich selbst fassen und bedenken.

2. Summa, wir sehen hier abermal, wie Gott seine Heiligen so wunderlich führt. Erstlich in dem Stücke, daß er den Altvater Jakob so lange läßt sitzen, nämlich, dreizehn ganze Jahre, in großer Trübniß und Leid, und zwar Joseph auch, daß sie nicht anders meinen, es sei aus mit ihnen; und darnach plötzlich so große Freude anrichtet, daß auch die, so es hören, müssen fröhlich werden, daß es den beiden so herzlich wohl geht.

3. Das ist der Punkt, den uns die Schrift überall vorhält: Was Gott will zu Ehren machen, das macht er vorhin zu Schanden; wen er auf das höchste erfreuen will, den macht er vor voll Traurens und Herzeleid. Also, daß er uns mit so viel Crempeln, damit er uns überschüttet, gerne dahin bringen wollte, daß wir auch einmal lernten seine Weise kennen, wie er handelt in allen Creaturen, sonderlich mit den Auserwählten. Je tiefer er sie demüthiget, drückt, und sinken läßt, je höher er sie erheben und empor setzen will.

4. Also sind aus dieser Historie viel Sprüche und Prophezeiungen der Schrift geflossen, daß die Väter daher genommen und gelernt haben. Darum sollten wir es doch auch einmal lernen; aber die Natur ist zu schwach, sie kann nicht anders richten, denn wie sie fühlt. Sollte Joseph dem gefolgt haben, wie er im Kerker lag, und seine Ehre und Gerücht ihm genommen war, daß er zweierlei Unglück auf einmal leiden müßte, Strafe und Schande, und beides unschuldig; denn es thut dennoch noch sanft, wenn einer so leidet, daß ihn die Leute wissen zu klagen, und Mitleiden haben, daß ihm Unrecht geschieht; aber wenn einer unschuldig mit allen Schanden steht, das ist erst bitter: sollte\footnote{Erlanger: soll. Wittenberger und Jenaer: „sollt“.} er nun (sage ich) darnach gerichtet haben, so hätte er hundertmal verzweifeln müssen. Wiederum, hätte er sich des versehen, daß er so zu großen Ehren kommen sollte, so hätte er gerne zehnmal so viel gelitten. Wer will es\footnote{Wittenberger: „muß“ statt: es.} ausstreichen und preisen, daß ihm widerfährt, daß ihn Gott ehrt zeitlich und ewig, geistlich und leiblich, daß er ein Herr wird über alle Güter Egyptenlands, und währet nicht zehn oder zwanzig, sondern siebzig, ja, fast achtzig Jahre!

5. Rechne nun gegen einander: Dreizehn Jahre leidet er Schmach, Schande und Schuld; die Zeit ist ihm gar lang worden, Fleisch und Blut könnte es nicht so lange erhalten; darnach wird er ein Herr, nicht allein über die Frau, die ihm muß zu Füße fallen und Gnade bitten, sondern über das ganze Land.

6. Darum kann es Gott nicht lassen: demüthiget er uns, und wirft uns herunter, so will er uns desto fröhlicher empor heben. Das will er uns einbilden, wenn wir es nur könnten erharren und Geduld haben. Das meint auch Paulus 2 Cor. 4, 17.; daß wir hier leiden, ist ein kurzer Augenblick, aber es schafft so viel desto reichere Ehre, die in uns soll offenbar werden, wenn wir nicht aufsehen auf das Sichtbare, sondern auf das Unsichtbare.

7. Also will uns Gott trösten; als wollte er sprechen: Sehet doch, wie ich gesinnt bin, nehmet vor euch den Joseph; wie ich mit ihm\footnote{„mit ihm“ fehlt in der Erlanger.} gehandelt habe, so will ich mit euch auch handeln. Dreizehn Jahre habe ich ihn lassen leiden, aber viel reichlicher vergolten mit leiblicher, und viel mehr mit geistlicher Ehre und Gütern, daß er
\quellseite{308}
der Heiland und Engel, von GOtt gesandt, wäre aller, die da sind, die Leute zu erleuchten, und zu GOtt [zu] bekehren.

8. Das ist das erste und beste Stück, die Summa der ganzen Historie, wie er zugleich ein weltlicher Herr, und geistlicher Regierer der Seelen wird, und ein rechter Bischof über alles, was in das Land schlägt; denn ohne Zweifel, seine Lehre ist sehr weit erschollen. So hat er nicht gefeiert, weil er alles in der Hand hatte und regierte, daß die rechte Lehre stark im Schwange ginge. Da kannst du denken, was er für Frucht geschafft habe mit Leuten unterweisen, daß sie GOtt erkenneten; wie der Psalter sagt [Psalm 32, 8.]. Aber es ist nicht dahin gekommen, er hat zuvor wohl müssen gebrochen werden.

9. Darüber ist noch ein groß Stück, wie ich gesagt habe, das dritte, daß er ewig ein Herr worden ist, also, daß er der Oberste ist worden im Volke Israel. Denn Israel hatte zwölf Söhne, von welchen GOtt den einen, nämlich Levi, sonderlich auszog zum Priesterthume; da blieb es bei. Darnach nimmt GOtt die zwei Söhne des Joseph, und macht die Zahl wieder voll, gibt ihm zwei Stämme vor den andern Brüdern, dazu auch das Königreich Israel, so lange es gestanden ist; also, daß der Mann nicht allein zeitlich und geistlich gesegnet ist, sondern auch sein Geschlecht und das Volk nach seinem Leben so lange geblieben ist, daß sie treffliche Thaten gethan, und oft Ehre eingelegt haben wider ihre Feinde, und große Propheten gehabt, daß sein Volk also gebenedeiet ist mit seinem Blut und Fleisch als dieses. Hätte er das alles zuvor gesehen, würde er gerne hundert Jahre den Tod gelitten haben, und wäre noch zu gering gewesen gegen das überschwängliche Gut, damit er gepreiset wird mit seinem Volke. Denn kein Volk auf Erden gewesen ist, das solche Leute, Propheten, Könige und Fürsten gehabt habe, als das Volk Israel, sonderlich Ephraim; ausgenommen den\footnote{Wittenberger: der.} Stamm Juda.

10. Das hat alles GOtt gethan um zukünftiges Trostes willen, wenn wir leiden; wo wir nur könnten stille stehen und halten, und gerne mehr leiden. Aber wir sind zu weich, wollen es nur zuvor sehen, verhindern damit keinen Rath. Wen nun das Exempel nicht bewegt, was soll ihn denn bewegen? Wie sind wir so gar Fleisch und Blut gegen diese! Darum auch, weil wir nicht können warten und harren, so werden wir nimmer gewahr, wie er die Seinen tröstet und segnet. So hat es ihm ja wehe gethan, daß er die Schmach mußte leiden; aber hätte er sollen sehen, was GOtt damit im Sinne hatte, hätte er gerne tausend Ehebrüche lassen über sich gehen; weil er aber ausharrt, so erfährt er, wie reichlich GOtt vergelten kann.

11. Das ist es, daß Paulus [2 Cor. 4, 17.] sagt: „Gering und leicht ist das Leiden gegen der Herrlichkeit, die uns GOtt geben will“, wenn wir glauben. Aber wenn das Leiden da ist, scheint es so groß, daß niemand meint, daß etwas Trostes da sei, also, daß es auch den Heiligen hier gefehlt hat; als auch dem heiligen Hiob, wie er [Cap. 6, 2. 3.] sagt: „Wenn man meinen Zorn wöge, und mein Leiden zusammen in eine Waage legte, so würde es schwerer sein denn Sand am Meere.“ Wer es aber GOtt heimstellen kann, der sieht das, wo er Eines kann leiden, daß da der ganze Himmel voll Trostes ist; als Christus sagt im Evangelio [Matth. 19, 29.]: „Ein jeglicher, der da verlässet Häuser, oder Brüder, oder Schwestern, oder Vater, oder Mutter, Weib oder Kind, oder Aecker, um meines Namens willen, der wird es hundertfältig nehmen, und das ewige Leben ererben.“

12. Da siehst du je, wie es wahr ist; er läßt Jakob mit elf Brüdern, und gewinnt das ganze Egyptenland, darnach auch das ewige Leben. Das können wir nicht thun. Das macht, wir wollen nur die Faust vorhin voll haben, und nicht glauben; indeß kommen wir nimmer dazu. Also fasse nun diese Historie, daß du wissest, was GOtt damit meint, und lernest, wie wir uns sollen halten, wenn er mit uns auch also fährt: daß man könne Trost daraus schöpfen, wenn es uns übel geht. Das sei das erste, vornehme Stück.

13. Darnach ist auch das, wie GOtt seine Heiligen so wunderlich regiert, daß er die Natur und Neigung nicht hinweg nimmt, so uns angeboren ist. Das predige ich nicht umsonst so oft, denn es wird gewöhnlich wieder kommen; wie man bereits etliche solche Heiligen findet, wie die gewesen sind, so in die Wüsten von Leuten gelaufen sind, und wie die Mönche und Nonnen vorgenommen haben, daß die Leute sollten Holz und Steine sein, meinten, so hoch in Tugend und Heiligkeit zu kommen, daß sie
\quellseite{309}
den Engeln gleich wären. Das sieht man in der ganzen Schrift von keinem Heiligen; denn dieselben natürlichen Neigungen sind Gottes Creaturen, die haben sie verleugnet und für Heiligkeit wollen machen, daß wir weder Blut noch Fleisch wären. Ich halte, Jakob sei ja so heilig, als irgend ein Papist, ja, St. Petrus und Paulus, und alle Apostel müßten ihm wohl die Ehre thun, und lieber Vater heißen, denn er überaus hoch und reich begnadet ist von GOtt: noch siehe, wie ihn GOtt läßt heulen und weinen, und wehklagen über den Sohn, als sei er\footnote{Jenaer und Erlanger: es.} gar fleischlich. Joseph, da er nun so viel erlitten hat, ist voll Heiligen Geistes und aller Güter; aber siehe, wie er so weich wird, und das Herz in ihm erwärmet über seinem leiblichen Bruder Benjamin, daß er muß von ihm gehen und weinen. Das ist eine natürliche Neigung, von GOtt eingepflanzt, die er nicht verdammen will. Also werden die Brüder auch froh, da sie ihn erkennen, und er, der Vater, wird darnach ganz wieder lebendig, als ob er aus einem tiefen Schlafe erwachte.

14. Das rede ich alles darum, daß man das rechte Werk GOttes erkenne; denn es ist seiner Meinung keine, daß man die Natur brechen will. Er hat die Augen geschaffen zu sehen, die Hände zu arbeiten, den Magen zu [ver]bauen; eine jegliche Creatur, wie er sie gemacht hat, so läßt er sein Werk bleiben; also hat er den Weibern ein Mutterherz gegeben, für die Kinder zu sorgen, welches niemand kann wehren noch hindern.

15. Darum soll man sein Reich also lassen bleiben, daß es ja die Natur nicht breche, daß man nicht wiederum solche tollen Heiligen mache, die ihnen das Herz selbst nehmen wollen. Das thut GOtt wohl, daß er die natürliche Bewegung versucht, und sich mit ihr zerrt, und wehrt ihr oder hält sie auf. Also war es natürlich, daß Jakob seinen Sohn gerne behalten hätte; aber GOtt reißt ihn hinweg, und nimmt doch dem Vater die Neigung nicht; thut nicht, was er gerne sähe, aber den Willen nimmt er ihm nicht. So verbeut er nicht, das oder jenes lieb zu haben, aber er hindert es oft, daß nicht geschieht, was wir gerne wollten; als, daß man einen recht schilt und verspricht\footnote{versprechen = übel von jemand reden.}, den ich gerne hörte loben; oder ein guter Freund stirbt, den ich nicht gerne verliere, und dergleichen.

16. Die Klüglinge aber meinen, es müsse alles hinweg, daß man des Dinges gar frei sei, und haben solche Meinung, [sie] wollen es mit Werken hinweg treiben. Er hat es geschaffen, daß es soll und muß sein; so wollen wir es besser machen, mit Geboten wehren und hinweg reißen. [Es] ist eben so viel, als wenn ich sagte: Das Auge soll nicht Auge sein, der Bauch soll nicht Bauch sein; doch hungert meinen Bauch gleichwohl, wenn es Zeit ist, und läßt ihn gebieten und verbieten, wie man will. Also thun alle natürlichen eingepflanzten Werke; die läßt er bleiben, aber martert sie wohl, daß er uns versuche. Darum siehe dich vor, vor denselben großen Heiligen, als die bald nach der Apostel Zeit aufgestanden sind; als man von etlichen liest in vitis patrum. Hat GOtt den Heiligen die Natur nicht genommen, die er in der Welt hat lassen ausschreien und beschreiben für die Besten, die er konnte finden, und so er sie besser gehabt hätte, auch hätte lassen schreiben, und sie vorgezogen als einen Ausbund und seine liebsten Kinder, so wollen wir es auch dabei bleiben lassen, und nicht heiliger werden.

17. Aber das müssen wir erwarten und uns erwegen, daß er sie wohl versucht und panzerfegt\footnote{panzerfegen = reinigen, poliren, reiben, martern. — sich verzeihen = sich einer Sache entschlagen.}, zu prüfen, wie wir in solchem Wesen wollen halten, ob wir um seinetwillen uns können alles Dinges verzeihen. Darum gibt er uns das Leben, welches das köstlichste Ding ist, das wir auf Erden haben, aber wirft uns oft in Fährlichkeit des Todes, zuletzt entweicht er es uns gar, daß man siehe, ob wir es um des willen, der es gegeben hat, können fahren lassen. Sollte es darum unrecht sein, das Leben lieb haben? Da sind seine Gaben nicht verworfen, sondern allein versucht, ob man sie lassen könne; wie er den Jakob versucht, ob er den Sohn könnte lassen um seinetwillen. Er that es auch, aber es that ihm doch wehe. Denn unmöglich ist’s, daß der Natur nicht sollte wehe thun, wenn man die eingepflanzte Bewegung brechen soll.

18. Das sind die vornehmsten Stücke in der Historie. Denn, daß Joseph mit den Brüdern also handelt und disputirt, und sie wohl martert und streckt\footnote{strecken = peinigen. Der Ausdruck ist hergenommen von der Folter.}, dient alles dahin, daß man immer GOttes Werke sehe. Er will sie fröhlich machen,
\quellseite{310}
und ein gut Mahl zurichten, darum tribulirt er sie vor recht wohl. Dabei man spüren könne und gewiß sein, wenn uns GOtt anhebt zu martern, daß er’s dann gut im Sinne hat; da ist er vor der Thüre, wie er auch in der Offenbarung Johannis [Cap. 3, 20.] spricht: „Ich stehe vor der Thüre“; aber das Klopfen thut uns wehe. Siehe aber, wie Joseph stille hält, und thut ihm auf, und läßt ihn hinein gehen. Also richtet er das Hochzeitsmahl zu. Das sei von der Historie genug geredet. Alle Punkte können wir nicht ausrichten, daß es nicht zu lang werde; wollen etliche sparen in die Mysteria, die wollen wir ausstreichen.

19. Dieser Joseph ist, wie wir oben [Cap. 37, § 10] gehört haben, eine Figur unsers Herrn Christi; darum hat er auch den Namen. Denn Joseph heißt ein Zunehmer, daß er immer wächst und zunimmt, und mehr und mehr wird, sammelt und häuft sich. Den Namen hat ihm GOtt auch zeitlich erfüllt, wie gehört, daß er ihn überhäuft und überschüttet hat mit zeitlichen,\footnote{Erlanger: Zeitlichem.} weltlichen Gütern.

20. Der Joseph hat nun einen bunten Rock gehabt, und war das rechte, liebe Kind des Vaters, dem auch das Regiment gebührte, wie auch seine Träume ausweisen; aber seine Brüder werden ihm feind, ziehen ihn aus, und verkaufen ihn in Egypten, daß er kommt an einen fremden Ort, und muß daselbst den Fremden dienen, und eine fremde Frau macht ihn zu Schanden, ergreift ihn einmal bei dem Mantel, den läßt er in ihrer Hand, und fleucht davon etc. Mit der Figur hat der Heilige Geist zuvor gespielt von dem Leiden und Auferstehung Christi. Denn daß er in Egypten kommt, ist, daß er Mensch worden ist auf Erden; da wird er ein Knecht und dient, hat nicht mehr denn Leibes Nothdurft davon.

21. Da wirft seine Frau die Augen auf ihn und will mit ihm buhlen; da er nicht will, ergreift sie ihn bei dem Mantel, und verklagt ihn vor ihrem Herrn. Also ist es ihm ergangen auf Erden. Die Synagoge oder jüdische Volk ist diese Hure, Mosis, seines Herrn Weib, welcher ist der Meister dieses Volkes, und regiert es wie ein Mann das Weib; unter dem hat Christus gedient, wie Paulus sagt, Gal. 3, 13. 4, 4., und sich unter das Gesetz gegeben und gehalten, ob ers wohl nicht bedurfte. Denn er ist fremd hinein gekommen, und hat ihm nichts zugehört, dient ihm aber also, daß er ihm alles unter die Hände thut, und selbst nicht weiß, was er hat. Das ist, daß Christus dem Gesetze genug gethan und alles erfüllt hat, daß [es] nicht mehr fordern kann.

22. Die Frau aber, weil sie sieht, daß er hübsch und schön von Angesicht war, gewann sie Lust und Liebe zu ihm, und reizte ihn täglich mit Worten, daß er bei ihr schliefe; er aber wollte nicht, so gar, daß er auch floh, um sie zu sein, so lange bis sie ihn unversehens heimlich erwischte. Das ist die Figur, da die Pharisäer und Schriftgelehrten haben seine Exempel und Wunderwerke, wie ihm das Volk anhing und groß Geschrei hatte, konnten sie seine Geberde und Wandel nicht strafen. Darum hätten sie gerne mit ihm gebuhlt, das ist, dahin gezogen, daß er es mit ihnen hielte, und predigte, wie sie wollten, und ließe ihr Ding recht sein; als im Evangelio Johannis sein angezeigt ist, wie sie ihn gerne zu sich gelockt hätten. Aber es ward nichts daraus, er sah ihre Büberei wohl. Wie oft sie es versuchten, so half es nicht; also blieb er unschuldig von ihnen.

23. Darnach erwischten sie ihn zuletzt und ergriffen ihn bei dem Mantel. Das geschah, da sie ihn gefangen nahmen und verklagten, er wäre ein Verführer, und hätte das Volk erreget und an sich gehangen, und wider das Gesetz gepredigt. Da verurtheilte man ihn als einen Ehebrecher mit allen Schanden. Daß er in [den] Kerker geworfen wird, das ist, abgethan, gestorben und begraben; da ließ er den Mantel in ihrer Hand, den behielten sie, das ist, sie haben die Schrift noch, die hat er hinter ihm gelassen; aber er ist davon, daß sie ihn nicht kriegen. Und wie die Hure den Mantel zum Zeugniß wider ihn heilegt und behält, also führen sie die Schrift, die wider sie selbst war, wider ihn zum Zeugniß, als, da sie sagten im Johanne [Cap. 19, 7.]: „Wir haben ein Gesetz, und nach dem Gesetze soll er sterben.“ Also muß er auf das allerschändlichste sterben, und liegt im Grabe.

24. Darnach gehen die Propheten her, das sind die Träumer, der Schenke und Bäcker; das ist, daß sich hieher reimen aller Propheten Sprüche, wie sein Reich ist, daß er die ganze Welt richtet, die Ungläubigen verdammt, und die Gläubigen selig macht. Indeß kommt Pharao
\quellseite{311}
mit seinem Traume; da kommt Christus hervor, und wird ein HErr seiner Feinde und aller Welt, der in Händen hat und regiert alles, was GOtt geschaffen hat im Himmel und auf Erden, daß alle Welt die Kniee müssen\footnote{„müssen“ fehlt in der Wittenberger.} vor ihm beugen.

25. Da ist die Auferstehung und das Königreich Christi auf das allerfeinste abgemalt. Denn wie Joseph im Zeitlichen regiert, also regieret er im Geistlichen, daß er alle Seelen speist mit seinem Worte, und die Gaben des Geistes austheilt, macht das ganze Land voll und satt, und wird sein Name und Geschlecht gemehret, daß unter seinem Volke aufstehen Könige und Propheten, also, daß seine Demuth und Elend über alle Maße reichlich bezahlt ist, wie Paulus Phil. 2, 8—11. sagt: „Er hat sich selbst erniedriget, und ist gehorsam worden bis zum Tode, ja, zum Tode des Kreuzes. Darum hat ihn GOtt einen Namen gegeben, der über alle Namen ist, daß in dem Namen Christi sich beugen sollen alle derer Kniee, die im Himmel und auf Erden, und unter der Erde sind, und alle Zungen bekennen sollen, daß Jesus Christus der HErr sei, zum Preise GOttes, des Vaters.“ Solches hat er alles durch sein Leiden erworben, doch alles uns zu Dienste, wie Joseph allein dazu in Egypten von GOtt geschickt ward, daß durch ihn vielen Leuten geholfen ist.

26. Die heimliche oder geistliche Deutung der andern Capitel lasse ich jetzt um der Kürze willen anstehen, denn sie sind leicht aus den vorigen Deutungen zu finden, auf daß andere auch etwas zu thun haben.

\section{Das sechsundvierzigste Capitel}

V. 1—34. Israel zog hin mit allem, das sein war. Und da er gen Bersaba kam, opferte er Opfer dem GOtt seines Vaters Isaac. Und GOtt sprach zu ihm des Nachts im Gesichte: Jakob, Jakob. Er sprach: Hier bin ich. Und er sprach: Ich bin der HErr, der GOtt deines Vaters, fürchte dich nicht in Egypten hinab zu ziehen, denn daselbst will ich dich zum großen Volke machen. Ich will mit dir hinab in Egypten ziehen, und will auch dich herauf führen, und Joseph soll seine Hände auf deine Augen legen. Da machte sich Jakob auf von Bersaba, und die Kinder Israel führten Jakob, ihren Vater, mit ihrem Kindlein und Weibern auf den Wagen, die Pharao gesandt hatte, ihn zu führen. Und nahmen ihr Vieh und Habe, die sie im Lande Canaan erworben hatten. Und kamen also in Egypten, Jakob und all sein Same mit ihm, seine Kinder und seine Kindeskinder mit ihm, seine Töchter, und seine Kindestöchter, und all sein Same [; die brachte er mit sich in Egypten]. Dies sind die Namen der Kinder Israel, die in Egypten kamen: Jakob und seine Söhne. Der erstgeborne Jakobs Sohn, Ruben. Die Kinder Ruben: Hanoch, Pallu, Hezron und Charmi. Die Kinder Simeon: Jemuel, Jamin, Ohad, Jachin, Zohar und Saul, der Sohn von dem cananäischen\footnote{Jenaer und Erlanger: cananäischen.} Weibe. Die Kinder Levi: Gerson, Kahath und Merari. Die Kinder Juda: Ger, Onan, Sela, Perez und Serah. Aber Ger und Onan starben im Lande Canaan. Die Kinder aber Perez: Hezron und Hamul. Die Kinder Issachar: Thola, Phua, Job und Semron. Die Kinder Sebulon: Sered, Elon und Jahleel. Das sind die Kinder von Lea, die sie Jakob gebar in Mesopotamia, mit seiner Tochter Dina, die machen allesammt, mit Söhnen und Töchtern, drei und dreißig Seelen. Die Kinder Gad: Ziphion, Haggi, Suni, Ezbon, Eri, Arodi und Areli. Die Kinder Asser: Jemna, Jesua, Jesui, Bria, und Serah, ihre Schwester. Aber die Kinder Bria: Heber und Malchiel. Das sind die Kinder von Silpa, die Laban gab Lea, seiner Tochter, und gebar Jakob diese sechzehn Seelen. Die Kinder Rahel, Jakobs Weibes: Joseph und Benjamin. Und Joseph wurden geboren in Egyptenland Manasse und Ephraim, die ihm gebar Asnath, die Tochter Potiphera, des Priesters zu On.
\quellseite{312}
Die Kinder Benjamin: Bela, Becher, Asbel, Gera, Naaman, Ehi, Ros, Muppim, Huppim und Ard. Das sind Kinder von Rahel, die Jakob geboren sind, allesamt vierzehn Seelen. Die Kinder Dan: Husim. Die Kinder Naphtali: Jahzeel, Guni, Jezer und Sillem. Das sind die Kinder Bilha, die Laban seiner Tochter Rahel gab, und gebar Jakob diese sieben Seelen. Alle Seelen, die mit Jakob in Egypten kamen, die von seinen Lenden ausgegangen waren (ausgenommen die Weiber seiner Kinder), sind alle zusammen sechs und sechzig Seelen. Und die Kinder Josephs, die in Egypten geboren sind, waren zwei Seelen. Also, daß alle Seelen des Hauses Jakob, die in Egypten kamen, waren siebzig. Und er sandte Juda vor ihm hin zu Joseph, daß er ihm den Weg weisete gen Gosen, und kamen in das Land Gosen. Da spannte Joseph seinen Wagen an, und zog hinauf seinem Vater Israel entgegen gen Gosen. Und da er ihn sahe, fiel er um seinen Hals, und weinete an seinem Halse sehr. Da sprach Israel zu Joseph: Ich will nun gerne sterben, nachdem ich dein Angesicht gesehen habe, daß du noch lebest. Joseph sprach zu seinen Brüdern, und seines Vaters Hause: Ich will hinauf ziehen, und Pharao ansagen, und zu ihm sprechen: Meine Brüder und meines Vaters Haus ist zu mir gekommen aus dem Lande Canaan, und sind Viehhirten; denn es sind Leute, die mit Vieh umgehen; ihre Schafe und Rinder, und alles, was sie haben, haben sie mitgebracht. Wenn euch nun Pharao wird rufen, und sagen: Was ist euer Handel? so sollet ihr sagen: Deine Knechte sind Leute, die mit Vieh umgehen, von unserer Jugend auf bisher, beide, wir und unsere Väter, auf daß ihr wohnen möget im Lande Gosen. Denn was Viehhirten sind, das ist den Egyptern ein Greuel.

\section{Das siebenundvierzigste Capitel}

V. 1—31. Da kam Joseph, und sagte es Pharao an, und sprach: Mein Vater und meine Brüder, ihre Schafe und Rinder, und [alles] was sie haben, sind kommen aus dem Lande Canaan, und siehe, sie sind im Lande Gosen. Und er nahm seiner [jüngsten] Brüder fünf, und stellte sie vor Pharao. Da sprach Pharao zu seinen Brüdern: Was ist euer Handel? Sie antworteten: Deine Knechte sind Viehhirten, wir und unsere Väter. Und sagten weiter zu Pharao: Wir sind kommen, bei euch zu wohnen im Lande; denn deine Knechte haben nicht Weide für ihr Vieh, so hart drückt die Theurung das Land Canaan. So laß doch nun deine Knechte im Lande Gosen wohnen. Pharao sprach zu Joseph: Es ist dein Vater, und sind deine Brüder, die sind zu dir gekommen; das Land Egypten steht dir offen, laß sie am besten Orte des Landes wohnen, laß sie im Lande Gosen wohnen. Und so du weißest, daß Leute unter ihnen sind, die tüchtig sind, so setze sie über mein Vieh. Joseph brachte auch seinen Vater [Jakob] hinein, und stellte ihn vor Pharao. Und Jakob segnete den Pharao. Pharao aber fragte Jakob: Wie alt bist du? Jakob sprach [zu Pharao]: Die Zeit meiner Wallfahrt ist hundert und dreißig Jahre, wenig und böse ist die Zeit meiner Wallfahrt, und langet nicht an die Zeit meiner Väter in ihrer Wallfahrt. Und Jakob segnete den Pharao, und ging heraus von ihm. Aber Joseph schaffte seinem Vater und seinen Brüdern Wohnung, und gab ihnen einen Sitz in Egyptenland, am besten Orte des Landes, nämlich, im Lande Raemses, wie Pharao geboten hatte. Und er versorgte [seinen Vater und] seine Brüder, und das ganze Haus seines Vaters, wie die jungen Kinder, mit Brod. (Es war aber kein Brod in allen Landen, denn die Theurung war fast schwer, daß das Land Egypten und Canaan verschmachteten vor der Theurung. Und Joseph brachte alles Geld zusammen, das in Egypten und Canaan gefunden ward, um das Getreide, das sie kauften, und er that alles Geld in das Haus Pharao. Da nun Geld gebrach im Lande Egypten und Canaan, kamen alle Egypter zu Joseph, und sprachen: Schaffe uns Brod, warum lässest du uns vor dir sterben, darum, daß wir ohne Geld sind? Joseph sprach: Schaffet euer Vieh her, so will ich euch um das Vieh geben, weil ihr ohne Geld seid. Da brachten
\quellseite{313}
sie Joseph ihr Vieh. Und er gab ihnen Brod um ihre Pferde, Schafe, Rinder und Esel. Also ernährte er sie mit Brod das Jahr, um alle ihr Vieh. Da das Jahr um war, kamen sie zu ihm im andern Jahre, und sprachen zu ihm: Wir wollen unserm Herrn nicht verbergen, daß nicht allein das Geld, sondern auch alles Vieh dahin ist zu unserm Herrn, und ist nichts mehr übrig vor unserm Herrn, denn nur unsere Leiber, und unser Feld. Warum lässest du beide, uns vor dir sterben, und unser Feld? Nimm uns an und unser Land um das Brod, daß wir und unser Land dienstbar seien dem Pharao, gib uns Samen, daß wir leben und nicht sterben, und das Feld nicht verwüste. Also nahm Joseph ein dem Pharao das ganze Egypten. Denn die Egypter verkauften ein jeglicher seinen Acker, denn die Theurung war zu stark über sie. Und ward also das Land Pharao eigen, sammt dem Volk, das zu seinen Städten aus- und einging, von einem Orte Egyptens bis ans andere. Ausgenommen der Priester Feld, das nahm er nicht ein; denn es war von Pharao für die Priester verordnet, daß sie essen sollten, was ihnen genannt war, das er ihnen gab. Darum durften sie ihr Feld nicht verkaufen. Da sprach Joseph zu dem Volke: Siehe, ich habe heute eingenommen euch und euer Feld dem Pharao; siehe, da habt ihr Samen und besäet das Feld, und von dem Getreide sollt ihr den Fünften Pharao geben, vier Theile sollen euer sein, zu besäen das Feld, zu eurer Speise, und für euer Haus und Kinder. Sie sprachen: Laß uns nur leben, und Gnade vor dir, unserm Herrn, finden, wir wollen gerne Pharao Knechte sein. Also machte Joseph ihnen ein Gesetz bis auf diesen Tag, über der Egypter Feld, den fünften Pharao zu geben, ausgenommen der Priester Feld, das ward nicht eigen Pharao. Also wohnete Israel in Egypten, im Lande Gosen, und hatten es inne, und wuchsen und mehreten sich sehr. Und Jakob lebte siebzehn Jahre in Egyptenland, daß sein ganzes Alter war hundert und sieben und vierzig Jahre. Da nun die Zeit herbei kam, daß Israel sterben sollte, rief er seinem Sohne Joseph, und sprach zu ihm: Habe ich Gnade vor dir funden, so lege deine Hand unter meine Hüfte, daß du Barmherzigkeit und Treue an mir thust, und begräbest mich nicht in Egypten; sondern ich will liegen bei meinen Vätern, und du sollst mich aus Egypten führen, und in ihrem Begräbnisse begraben. Er sprach: Ich will thun, wie du gesagt hast. Er aber sprach: So schwöre mir. Und er schwur ihm. Da neigte\footnote{Randglosse: knieete. Die alten Bücher im Griechischen und Latein haben hier also: Und er betete an die Spitze seines Scepters, wie die Epistel zu den Hebräern am 11. [V. 21.] meldet, als habe Jakob das Scepter Josephs angebetet; aber im Hebräischen ist es, wie hier steht.} Israel gegen dem Bette zum Häupten.

\section{Das achtundvierzigste Capitel}

V. 1—22. Darnach ward Joseph gesagt: Siehe, dein Vater ist krank. Und er nahm mit sich seine beiden Söhne, Manasse und Ephraim. Da ward es Jakob angesagt: Siehe, dein Sohn Joseph kommt zu dir. Und Israel machte sich stark, und setzte sich auf das Bette, und sprach zu Joseph: Der allmächtige Gott erschien mir zu Lus im Lande Canaan, und segnete mich, und sprach zu mir: Siehe, ich will dich wachsen lassen und mehren, und will dich zum Haufen Volks machen, und will dies Land zu eigen geben deinem Samen nach dir ewiglich. So sollen nun deine zwei Söhne [Ephraim und Manasse], die dir geboren sind in Egypten, ehe ich herein kommen bin zu dir, mein sein, gleich wie Ruben und Simeon. Welche du aber nach ihnen zeugest, sollen dein sein; aber diese sollen genannt sein mit ihrer Brüder Namen in ihrem Erbtheile. Und da ich aus Mesopotamia kam, starb bei mir Rahel, im Lande Canaan, auf dem Wege, da noch ein Feldweges war gen Ephrath, und ich begrub sie [daselbst] an dem Wege gen Ephrath, die nun Bethlehem heißt. Und Israel sahe die Söhne Josephs, und sprach: Wer sind die? Joseph antwortete seinem Vater: Es sind meine Söhne, die mir Gott hier gegeben hat. Er sprach: Bringe sie her zu mir, daß ich sie segne. Denn die Augen Israel waren schwer worden vor Alter, und konnten nicht wohl sehen.
\quellseite{314}
Und er brachte sie zu ihm. Er aber küssete sie, und herzte sie, und sprach zu Joseph: Siehe, ich habe dein Angesicht gesehen, daß ich nicht gedacht hätte, und siehe, GOtt hat mich auch deinen Samen sehen lassen. Und Joseph nahm sie von seinem Schooße, und sie fielen nieder auf die Erde auf ihr Angesicht. Da nahm sie Joseph beide, Ephraim in seine rechte Hand, gegen Israels linke Hand, und Manasse in seine linke Hand, gegen Israels rechte Hand, und brachte sie zu ihm. Aber Israel streckte seine rechte Hand aus, und legte sie auf Ephraims, des Jüngsten Haupt, und seine linke auf Manasse’s Haupt, und that wissend also mit seinen Händen, denn Manasse war der Erstgeborne. Und er segnete Joseph, und sprach: Der GOtt, vor dem meine Väter, Abraham und Isaak, gewandelt haben; der GOtt, der mich mein Lebelang ernähret hat, bis auf diesen Tag; der Engel, der mich erlöset hat von allem Uebel, der segne diese Knaben, daß sie nach meinem und nach meiner Väter, Abrahams und Isaaks, Namen genannt werden, daß sie wachsen, und viel werden auf Erden. Da aber Joseph sahe, daß sein Vater die rechte Hand auf Ephraims Haupt legte, gefiel es ihm übel, und fassete seines Vaters Hand, daß er sie von Ephraims Haupte auf Manasse’s Haupt wendete. Und sprach zu ihm: Nicht so, mein Vater, dieser ist der Erstgeborne, lege deine rechte Hand auf sein Haupt. Aber sein Vater weigerte das, und sprach: Ich weiß es wohl, mein Sohn, ich weiß es wohl, dieser soll auch ein Volk werden, und wird groß sein; aber sein jüngster Bruder wird größer denn er werden, und sein Same wird voll Volks werden. Also segnete er sie des Tages, und sprach: Nach deiner Weise werde Israel gesegnet, daß man sage: GOtt setze dich wie Ephraim und Manasse. Und setzte also Ephraim Manasse vor. Und Israel sprach zu Joseph: Siehe, ich sterbe, und GOtt wird mit euch sein, und wird euch wiederbringen in das Land eurer Väter. Ich habe dir ein Stück Landes gegeben, außer deinen Brüdern, das ich mit meinem Schwerdt und Bogen aus der Hand der Amoriter genommen habe.\footnote{In der Erlanger Ausgabe folgt hier der Text des 49. Capitels.}

Was hernach folgt, ist in der Summa eine Auslegung des 46. 47. 48. Capitels.

1. Da haben wir gehört, wie die schwere Theurung nun zwei Jahre gewähret hat in Egypten, ehe Jakob hinein gekommen ist, und noch fünf Jahre vorhanden war. Indes war nun verzehrt alles, was in Egypten war, auch das Geld, so im Lande und umher war, mit Haufen in des Königes Haus gekommen. Darnach, im sechsten Jahre, ging auch alles Vieh auf. Im siebenten ward Land und Leute des Königs eigen, wie die Knechte, daß er über sie Macht kriegte, wie über das Vieh. Darum hat ihm Joseph ein überschwänglich groß Reichthum zugebracht.

2. Also thut GOtt, wenn er ein Land reich und groß machen will, daß es muß groß werden, sollten auch alle anderen untergehen. Nun haben wir gehört, daß er ein feiner König gewesen sei, fromm und vernünftig, der GOtt erkannte, und GOttes Wort von Joseph gelernt hat. Wie geht es denn zu, daß Joseph, der heilige Mann, und Pharao, der fromme König, so greulich regieren, das Land beschweren und unterbringen, und lassen ihnen gar nichts Eigenes, machen sie zu armen, leibeigenen Leuten, so weit auch, daß er den Fünften ewig über das Land schlägt; ist das auch christlich oder billig?

3. Es ist jedermann wissend, daß ein gemeiner, ganghaftiger Zins in Landen ist, wiewohl es hier mit uns nicht ist, den man heißt den Zehuten, daß man jährlich den zehnten Theil der Obrigkeit gibt von allerlei Gütern im Lande; man hält es auch für ein zuträgliches, feines Recht. Ich wollte aber, daß man auch in diesem Lande den Zehuten, ja, ich dürfte noch rathen, daß man noch den Fünften möchte geben, also, daß es damit alles schlecht wäre; gleicher könnte man nicht handeln und zukommen. Denn wo man nicht den Zehuten gibt, sondern eine benannte Summe auf eine Stadt schlägt, daß eine so viel, die andere so viel gebe, so muß die Summe gefallen ein Jahr wie das andere, GOtt gebe, es gerathe, wie es wolle, ob sie allein kaum zwanzig oder vierzig Gulden einbringt. Geräth es einmal in sieben oder acht Jahren mit dem Einkommen und Vorrathe, so verdirbt es darnach wohl fünf oder sechs Jahre nach einander; also, daß es mit der Weile wohl die Hälfte hinweg kommt, und die Leute das mehrere Theil ihrer Güter zu Zins geben müssen. Was haben sie denn gewonnen?

4. Ich habe es überschlagen, was man bisher den geistlichen Bettlern gegeben hat, und gefunden
\quellseite{315}
, daß die Hälfte der armen Leute Güter dorthin gefallen sind, und sie sich von der andern Hälfte haben müssen nähren. Denn eines armen Mannes Gut, Acker, Vieh, und dergleichen, es trage etwas oder nichts, so müssen sie ein Jahr so viel geben, als das andere. Wenn aber ein Mann jährlich die zehnte Kuh und Kalb geben müßte, und so fortan, hätte er den Vortheil, trüge es ihm viel, so gäbe er viel, trüge es ihm wenig, so gäbe er wenig; so dürfte man die\footnote{Erlanger: der.} Leute nicht mit genannter Summe drücken, sondern müßte des Glücks mit ihnen warten. Aber jetzt muß ein armer Mann des Unglückes allein warten, die Oberherren sitzen frei und sicher; Gott gebe, man verderbe oder gedeihe, so müssen die Zinsen gefallen.

5. Nun sage, ob diese nicht vernünftigere Leute gewesen seien denn wir? haben natürlicher Weise Zins auf das Land gelegt, daß es hat mögen leiden und zukommen. Es wäre leidlicher, daß man das neunte, oder noch das achte Theil der Güter behielte, denn daß man oft die Hälfte hinweg gibt, zuweilen mehr, denn man einbringt. Daher ist es gekommen, daß die Städte haben müssen Zinse auf sich nehmen, und darüber verderben. Daher man klagt, es sei kein Geld und Gut schier in der Welt mehr, und bald jedermann verdorben ist.

6. Darum sind das weise Leute gewesen, ist auch nicht unbillig gethan, daß Joseph das Land und Güter unter den König gebracht hat. Das Korn und der Boden war sein; so war es eine Freundschaft, daß er es um einen gleichen Pfennig verkaufte, und Waare um Waare gab; welches in aller Welt gleich und billig ist, denn er war nicht schuldig, daß er sie alle ernährete; [es] wäre auch nicht gerathen. Denn man muß die Leute im Zaume und Zwange halten; sonst, wenn man Er\footnote{Er = Herr.} Omnes umsonst nährete, würde er zu muthwillig, und ginge auf das Eis tanzen.

7. Wenn eitel fromme Leute vorhanden wären, dürfte man sie nicht zwingen noch zähmen; aber der Bübel ist zu böse, so bald man ihm zu viel Willens läßt, so kann niemand mit ihm umgehen. Darum leidet es sich nicht, den gemeinen Mann ohne Zwang und Aufsätze gehen zu lassen. Es muß gefasset sein in ein Regiment und Ordnung, die muß man also lassen gehen, daß man ihnen nicht umsonst genug gebe, sondern Waare um Waare oder Geld. Wer wollte sonst arbeiten und still sitzen? Man sieht wohl, wie es ein Ding ist um den Haufen; wenn man ihm eine Hand breit erlaubt, so nimmt er vier und zwanzig Ellen.

8. Derhalben, wiewohl es ein hartes und greuliches Ding scheinen mag, daß Joseph solche Beschwerden macht über das Volk, ist es doch gnädig und nützlich gewesen, daß beide, [der] gemeine Mann und die Obrigkeit, dabei hat mögen gedeihen. Das ist aber kein Zweifel, daß Joseph so fromm gewesen ist, daß er arme Leute, die es nicht haben mögen bezahlen, nicht hat lassen Noth leiden, sondern umsonst gegeben; von andern aber, so es zu bezahlen hatten, hat er es auch mitgenommen. Weiter haben wir droben gehört, daß die heiligen Väter, beide Abraham und Isaac, und Jakob auch, viel eigene Leute oder Knechte haben gehabt, die man verkaufte, wie das Vieh; item, auch mehr Dinges gethan, das weltlich ist, als gekrieget und Blut vergößen; welches alles sich läßt ansehen, als nicht brüderlich oder christlich gehandelt. Aber christlich und brüderlich handeln, gehört nicht in das weltliche Regiment; Gott hat selbst gesetzt, wie man handeln soll mit Knechten, Mägdgen und ihren Kindern, und dergleichen.

9. Christlich und evangelisch Wesen gehört allein zu regieren die Gewissen; aber die Welt zu regieren, gehören scharfe und gestrenge Gesetze zu, daß man die Bosheit dämpfe, der sie alle voll ist; daß sie aber gedämpft werde, müssen die andern Frommen, die es nicht bedürfen, mithalten, um gemeinen Friedens willen. Also hat er nun das Land zu sich gebracht mit Gutt und Ehren. Wer es weiter rechnen will, der thue es, ob es schwer sei, den Fünften zu geben. Laß [es] kommen, daß man gebe den Sechsten, Siebenten, oder Achten. [Das] wäre alles zu rathen, wenn man andere Beschwerungen abthäte, Zölle, Zinstauf, und dergleichen, so jetzt täglich überhand nimmt, so behielten beide, die Obrigkeit und wir, genug. Und zwar im alten Testamente, bei den Juden, mußte man auch viel geben, den Zehnten und alle ersten Früchte, und Erstlinge von Vieh und Menschen, den Priestern, sonst auch den Zehnten des ersten Jahres, ohne andere Dinge [3 Mos. 27, 30. ff. Cap. 23, 10.], daß es wohl dem Fünften gleich worden ist, den Joseph aufgesetzt hat.
\quellseite{316}

10. Aber daneben schreibt Moses hier, wie die Priester ausgenommen sind, daß sie ihr Feld und Gut alles behielten, wie es für sie verordnet war, daß sie essen sollten, was ihnen bekannt war. Das ist auch der Texte einer, auf welchen sich des Papistes geistliches Recht will gründen, und spricht, daß seine Priester sollen ihre Güter frei haben, nicht unterworfen sein Zinsen oder Zehnten. Dazu sagt das Evangelium: Gehet hin und prediget, nehmet und esset, was man euch gibt, denn „ein Arbeiter ist seiner Speise werth“ [Matth. 10, 10.]. Und Paulus [1 Cor. 9, 14.]: „Der Herr hat befohlen, die das Evangelium predigen, sollen sich des Evangelii nähren.“ Also ist im alten Testamente geordnet gewesen, und hier auch unter den Heiden. Mich hat es oft gewundert von St. Paulo, daß er davon so viel schreibt, wie man den Bauch ernähre, und Christus selbst so sorgfältig dafür ist; wiewohl er es nicht bestimmt hat, wie viel man den Predigern soll geben, hat [er] es doch so weit befohlen, daß sie sich davon nähren und erhalten könnten.

11. Also ging es nun, weil Paulus lebte, daß man ihm schier nichts wollte geben, darum er so viel Worte davon macht 1 Cor. 9, 14., und Gal. 6, 6. 7. sagt er es ihnen hart: „Irret euch nicht, Gott lässet sich nicht höhnen; was der Mensch säet, das wird er ernten.“ Christus hätte auch selbst müssen nach Brode gehen und Hunger leiden, hätte er nicht können Steine zu Brode machen [Matth. 4, 3.], wenn es noth gewesen wäre, und die Weiber nicht an ihm gehangen wären, die ihm Nothdurft gegeben hätten. Wir können es auch nicht anders machen, es wird wohl bleiben, daß die rechten, frommen Prediger müssen der Welt halben nach Brod gehen. Einen rechten Prediger oder zwei können wir nicht ernähren, aber hundert und tausend Mönche und Pfaffen können wir wohl füllen, die uns zum Teufel führen.

12. Im alten Testamente konnten alle zwölf Stämme nicht den Einen Stamm Levi ernähren, daß sie zuletzt Gottes Wort mußten lassen liegen zu lehren, und sich mit den Händen nähren. Daneben war schier kein Flecken im Lande, es waren Altäre und Capellen da, da konnte man den Verführern reichlich geben: Als die Isabel [1 Kön. 18, 19.] allein vierhundert falsche Propheten von ihrem Tische nährte, und daneben sonst vierhundert Baals Propheten waren; dagegen waren hundert Propheten, die mußten Hunger und Noth leiden. Darum ist es wahr: Die da dienen im Worte, denen ist billig ihre Nahrung zu geben. Derhalben achte ich, daß Joseph und Pharao rechte Leute dazu gesetzt haben, und für das vornehmste Werk angesehen, daß die versorgt würden, als an denen mehr gelegen wäre, denn am ganzen Lande.

13. Ich habe auch gepredigt und geschrieben, daß vor allen Dingen sollte das Vornehmste sein, daß man die Pfarren und Schulen wohl versorgte. Aber da will niemand an. Anderen, Pfaffen und Mastbäuchen, die nichts thun, denn alles Unglück anrichten, gibt man übrig genug. Was Gott dazu sagen wird, werden wir auch einmal gewahr werden. Ich habe oft gesagt, Gott achtet seine Gaben theuer und werth; es hat ihn auch viel gekostet, er hat seines Sohnes Blut daran gewandt, daß das Wort möchte gepredigt werden, und seine Apostel ausgesandt würden. Darum denkt er also: Der Schatz ist zu edel und gut, sie sind nicht werth, daß sie ihn haben sollen; darum will ich meine werthsten Prediger selbst nähren. Darum gibt er der tollen Welt nicht in [den] Sinn, daß sie es thäte, denn seinen Christen, einem, oder zweien, oder dreien; die andern ließen wohl alle Prediger auf Einen Tag Hungers sterben. Es ist zu edel und zu groß, sie sind es nicht werth, dem Evangelio und seinen Predigern Ehre zu thun, möchten sonst sagen, sie hätten es verdient.

14. Also wird es bald wieder dahin kommen, daß die rechten Prediger ihr Amt müssen liegen lassen, und Handwerk lernen, allein um der Nahrung willen; so werden sie das Wort verlieren, durch Gottes Zorn und Strafe; und wiederum, denen mit Haufen zutragen, die ihnen Lügen und Verführung predigen. Weil der Schatz vorhanden ist, kehrt sich niemand daran; wenn er hinweg ist, so sieht man erst zurück. Darum muß es gehen, wie Christus [Joh. 5, 43.] sagt: „Ich bin gekommen in meines Vaters Namen, und ihr nehmet mich nicht an; so ein anderer wird in seinem Namen kommen, den werdet ihr annehmen.“ Wir mögen Gott danken, daß wir wissen, wie es müsse also gehen, und lassen die Welt zum Teufel laufen, und uns des trösten, daß wir einen Vater im Himmel haben, der uns ernähren wird; können wir nicht mehr, so arbeiten wir mit den Händen, so sind wir entschuldigt. Es ist uns nicht möglich
\quellseite{317}
zu predigen, und seine Nahrung davon haben, ob es wohl St. Paulus gethan hat; ist auch nicht vonnöthen. Wo es aber noth wäre, und ich mich es wollte Mühe lassen kosten, könnte ich mit GOttes Hülfe auch zugleich predigen, und ein Handwerk treiben.

15. Darum ist dies Exempel nicht dahin zu ziehen, da es der Pabst hinzeucht, daß sie in ihren Gütern sicher und fest sitzen, niemand nichts geben dürfen, welche sie nicht erworben haben, sondern von Fürsten gegeben sind, nicht um Predigens willen, sondern um Messens willen, und geachtet als ein gutes Werk, aber dem Teufel zu Dienste; darum sie auch müssen genug haben, und die rechten Prediger darben.

16. Aber dieser Pharao ist ein rechtschaffener Mann gewesen, und hat recht königlich gehandelt, daß er sich der Leute annimmt, die da predigen und lehren sollen. Man wird nicht viel heidnische Fürsten finden, die es ihm nachthun. Also hat er beide, weltliches und geistliches Regiment, auf das beste versorgt, daß sie keine Noth leiden an Leib und Seele.

17. Weiter ist nichts Sonderliches mehr in der Historie, denn im 48. Capitel, von den zwei Brüdern, Ephraim und Manasse, die er segnet und erwählt zu seinen Söhnen. Das zeucht er darum an, daß man hernach verstehe, so er oft meldet von den zwölf Stämmen des Volkes Israel; denn aus dem Joseph sind zwei Stämme gekommen, daß er zweimal so viel hat, als die andern. Denn es sollte ein Geschlecht ausgezogen werden zum Priesterthume, also, daß die Stätte wieder gefüllt, und die beiden Söhne eingerechnet mußten werden.

18. Da hat nun GOtt aber ein Wunder beweiset: Joseph bringt die zwei Kinder zum Vater, stellt den Aeltesten zu des Vaters rechten Seite, den andern zur linken, daß Jakob darnach die Hände auf sie legte, wie sie vor ihm standen; aber er kehrt es um, und thut Mutwillens wider seinen Sinn, daß es Joseph verdrießlich.

19. Nun ist es offenbar, wie im alten Testamente dem ersten Sohne viel Vorzug und Freiheit gegeben ward, daß er allezeit zwei Stücke Erbes vor andern nahm, dazu das Regiment und Priesterthum gewann, die andern mußten ihm unterthan sein, in geistlichem und weltlichem Regimente. Solches gebührt nun auch hier Rechts haben dem ersten Sohne Manasse; aber er entzeucht es ihm, und gibt es dem jüngsten, und spricht: „Dieser soll auch groß sein, aber der Jüngste wird größer, denn er, werden.“ Und aus dem Texte ward auch hernach der Stamm Ephraim stolz und hochmüthig, daß sein stolzer Volk unter den Juden war, das sich so brüstete auf den Segen; sie waren nur die Letzten, dazu die Jüngsten; noch waren sie die Größten. Denn daraus ist hernach bestanden das Königreich und Regiment in Israel. Wer nun die Propheten liest, der kann hieraus verstehen, warum sie den Namen Ephraim so anziehen, eben als wir bei uns nennen das römische Reich. Die Ehre hat verdient der Patriarch Joseph mit seinem Elende, so er in Egypten erlitten hatte; auf daß wir sehen sollen, wie GOtt wieder belohnt und zu großen Ehren setzt, die da können harren und erdulden seine Hand.

20. Daß er aber den jüngsten Sohn ehren vorzeucht, ist nach dem Evangelio gehandelt, daß GOtt also thut: Was groß vor der Welt ist, läßt er herunter; und wiederum, das Geringste hebt er empor. Welchen man nennt, der der Erste sein sollte, wird der Letzte; also, daß er nimmer seine Art läßt; alles, zu trösten die Christen, daß sie je mehr und mehr lernen, den GOtt erkennen, wie seine Art, Werk und Weise ist, nur herunter zu sehen, was klein ist; dagegen die Welt nur über sich sieht, was groß und hoch ist, und meint, der arme Haufe sei gar verlassen. Das müssen wir durch alle Exempel hören und sehen, dennoch bleiben wir immer im alten Sinne und Gedanken.

21. Dahin dient auch, daß wir oben [Cap. 43, 3.] gehört haben, wie Joseph seinen Brüdern befahl, sie sollten sein Angesicht nicht sehen, sie brächten denn den kleinsten Bruder mit sich. Das ist uns auch befohlen, daß wir immer Acht haben auf die Geringsten; dürfen vor GOtt nicht kommen, wir bringen denn die geringsten und schwächsten Christen mit uns. Wie Christus im Matthäo [Cap. 25, 40.] lehrt, und schließt also: „Was ihr gethan habt einem unter meinen geringsten Brüdern, das habt ihr mir gethan.“ Aber es will niemand herunter sehen, und sich unterwerfen zu Dienste derer, die da nichts sind. Darum möchten wir wohl dies Stück unter uns zum Sprichwort machen, und täglich deß ermahnen, weil uns Christus immerdar dahin weist.

22. Also haben wir diese Historie Josephs bis an das Ende gehört, damit nun erfüllt ist der Traum, der ihm zuvor geträumt hatte, wie
\quellseite{318}
auch der Text anzeigt, da er [Cap. 42, 9.] sagt: „Joseph gedachte an die Träume, die ihm von ihnen geträumet hatten“, darum sie ihm auch gram waren und von sich stießen, doch so wunderlich erfüllet sind worden, daß er eben dadurch zu solcher Ehre und Gewalt gekommen ist, damit sie meinten, ihn zu verderben, und nichts für so gewiß hielten, denn Joseph, wo er lebte, müßte alle seine Tage ein elender Mensch sein. Also fährt GOtt mit den\footnote{„den“ fehlt in der Erlanger.} Seinen, gibt ihnen Demuth, was geschehen soll, aber also, daß sich niemand versieht, wie es zugehen werde, sondern scheint, wie es solle zurück gehen, und wird doch so reichlich erfüllt. Es war ein Zeichen, daß ihn seine Brüder sollten gnädigen Herrn heißen, und vor die Füße fallen. Der Traum war von GOtt; darum mußte es so geschehen, daß [es] niemand glauben könnte; [es] mußte sich eben umkehren, daß sie ihn zum Knechte machten, und verkauften, und sicher waren, der Traum wäre lauter Lügen. Aber da sein allerdinge ungewiß ist, kommen sie, und fallen ihm unwissend zu Füßen.

23. Das sind alle Werke GOttes, die er thut und treibt. Er hat uns verheißen das ewige Leben, daß wir aus des Teufels Gewalt kommen, in ewige göttliche Gewalt. Das ist uns in der Taufe verheißen; aber wenn es soll angehen, so wendet es sich um; [wir] gehen dann mitten in [den] Tod, wie die lieben Märtyrer, die durch Schwert, Feuer und alles Unglück gehen mußten. Also, wenn er reich, hoch, fromm etc. will machen, so macht er arm, verachtet, Sünder und böse Gewissen, daß er es allenthalben so treibt nach dem äußerlichen Ansehen, als sollte es zunichte werden, und mit Lügen bestehen. Wohl dem, der es versteht, und sich nicht daran irrt noch ärgert! Joseph soll ein mächtiger Fürst werden, so wird er ein armer, gefangener Mann, dazu mit großen Schaden, als ein Ehebrecher; aber je größer Elend und Schande gewesen sind, je größere Wohlfahrt und Ehre ihm gegeben wird. Also ginge es uns auch, wenn wir nur einmal könnten trauen und harren, wir sollen keinen Mangel noch Noth an Leib und Seele haben, und GOtt will bei uns sein in aller Anfechtung und Trübsal, wie er sagt im 91. Psalm, V. 15., aber wenn er es angreift, so können wir nicht halten; da vergessen wir, was seine Art und Weise ist, sein Wort zu erfüllen, nämlich, daß er wegnimmt, was er geben will; spielt aber so freundlich mit uns, wie ein Vater mit seinen lieben Kindern, bis die Zeit kommt, daß er uns mit Freuden überschüttet.

\section{Das neunundvierzigste Capitel}

V. 1—33. Und Jakob berief seine Söhne, und sprach: Versammelt euch, daß ich euch verkündige, was euch begegnen wird in den letzten Zeiten. Kommet zu Haufen, und höret, ihr Kinder Jakob, höret euren Vater Israel. Ruben,\footnote{Handaufsatz: Ruben sollte der ersten Geburt Würde haben, nämlich das Priesterthum und Königreich. Nun aber wird beides von ihm genommen, und Levi das Priesterthum, und Juda das Königreich gegeben. Hier ist bedeutet die Synagoga, die das Bett Jakobs, das ist, die Schrift, befleckt mit falscher Lehre, darüber sie verloren hat Priesterthum und Königreich.} mein erster Sohn, du bist meine Kraft, und Anfang meiner Stärke, der Oberste im Opfer, und der Oberste in der Macht. Er fuhr leichtfertig dahin, wie Wasser. Du sollst nicht der Oberste sein, denn du bist auf deines Vaters Lager gestiegen, daselbst hast du mein Bette befleckt mit dem Aufsteigen. Die Brüder Simeon und Levi, ihre tödtliche Wehre sind schädlich Zeug. Meine Seele komme nicht in ihr Geheimniß, und meine Herrlichkeit sei nicht eins mit ihrer Sammlung. Denn in ihrem Zorn haben sie den Mann erwürget, und in ihrem Muthwillen haben sie den Ochsen verlähmet. Verflucht sei ihr Zorn, daß er so heftig ist, und ihr Grimm, daß er so störrig ist. Ich will sie zertheilen in Jakob, und zerstreuen in Israel. Juda, du bist es, dich werden deine Brüder loben, deine Hand wird deinen Feinden auf dem Halse sein. Vor dir werden deines Vaters Kinder sich bücken. Juda ist ein junger Löwe, du bist hoch gekommen, mein Sohn, sint dem Rauben; er hat nieder gekniet, und sich gelagert wie ein Löwe,
\quellseite{319}
und wie eine Löwin, wer will sich wider ihn auflehnen? Es wird das Scepter\footnote{Randglosse: Scepter. Hier sähet an der Segen von Christo, der von Juda geboren sollte werden, und heißt ihn Schilo, daß ist, der glückselig sei, und frisch durchdringen sollte mit Geist und Glauben, daß zuvor durch Werke sauer und unseelig Ding war. Darum nennen wir Schilo einen Helden. Denn das vorige Theil des Segens betrifft den König David, und ist sonst in allen Segen nichts mehr von Christo, sondern alles andere ist von zeitlichem Heil, das den Kindern Israel gegeben ist. Alles, daß Sebulun sollte am Meere wohnen bis gen Sidon, und Isaschar mitten im Lande vom Meere wohnen, und doch zinsbar gewesen ist den Königen von Assyrien. Den Segen Dan hat Simson erfüllt, Richt. 13, 24. 25. Gad hat seinen Segen ausgerichtet, wie sie vor Israel herzogen, Jos. 1, 14. Asser hat gut Getreideland inne gehabt. Naphtali Segen ist erfüllt durch Debora und Barak, Richt. 4, 6. 5, 1. Der Segen Josephs geht auf das Königreich Israel, und ist ganz vom leiblichen Regimente gesagt, daß die Töchter, das ist, die Städte im Lande, wohl regiert worden zeitlich, und viel Propheten und große Leute zu Hofeinen hatten. Und wiewohl sie oft angefochten wurden, gewannen sie doch; und das Königreich war im Geschlechte Ephraim. Also blieb der geistliche Segen und Reich auf Juda, und das leibliche Reich auf Ephraim. Benjamins Segen hat St. Paulus erfüllt, Phil. 3, 5., der der König Saul, und die Bürger zu Gaba, Richt. 20, 25.} von Juda nicht entwendet werden, noch ein Meister von seinen Füßen, bis daß der Held komme, und demselben werden die Völker zufallen. Er wird sein Füllen an den Weinstock binden, und seiner Eselin Sohn an den edlen Reben. Er wird sein Kleid in Wein waschen, und seinen Mantel in Weinerblute. Seine Augen sind röthlicher denn Wein, und seine Zähne weißer denn Milch. Sebulon wird am Anfurt des Meeres wohnen, und am Anfurt der Schiffe, und stoßen an Sidon. Isaschar wird ein beinerner Esel sein, und sich lagern zwischen die Grenzen. Und er sahe die Ruhe, daß sie gut ist, und das Land, daß es lustig ist, und hat seine Schultern geneigt zu tragen, und ist ein zinsbarer Knecht worden. Dan wird Richter sein in seinem Volke, wie ein ander Geschlecht in Israel. Dan wird eine Schlange werden auf dem Wege, und eine Otter auf dem Steige, und das Pferd in die Ferse beißen, daß sein Reiter zurück falle. HErr, ich warte auf dein Heil. Gad, gerüstet, wird das Heer führen, und wieder herum führen. Von Asser kommt sein fett Brod, und er wird den Königen niedliche Speise geben. Naphtali ist ein schneller Hirsch, und gibt hübsche Rede. Das holdselige Kind Joseph, das holdselige Kind anzugehen. Die Töchter treten einher im Regimente. Und wiewohl ihn die Schützen erbittert, mit ihm gezankt, und ihn gehaßt haben, so ist doch sein Bogen fest geblieben, und die Arme seiner Hände sind ermauert, durch die Hände des Mächtigen in Jakob. Aus ihnen sind gekommen Hirten und Steine in Israel. Von deines Vaters Gott ist dir geholfen, und von dem Allmächtigen bist du gesegnet, mit Segen oben vom Himmel herab, mit Segen von der Tiefe, die unten liegt, mit Segen an Brüsten und Bäuchen. Der Segen, deinem Vater und meinen Voreltern verheißen, gehet stark, nach Wunsch der Hohen in der Welt. Aus Joseph sollen Häupter werden, und oberste Nasiräer unter seinen Brüdern. Benjamin ist ein reißender Wolf; des Morgens wird er Raub fressen, aber des Abends wird er den Raub austheilen. Das sind die zwölf Geschlechte Israel alle, und das ist es, das ihr Vater mit ihnen geredet hat, da er sie segnete, einen jeglichen mit einem besonderen Segen. Und er gebot ihnen, und sprach zu ihnen: Ich werde versammelt zu meinem Volke, begrabet mich bei meine Väter in der Höhle auf dem Acker Ephron, des Hethiters, in der zwiefachen Höhle, die gegen Mamre liegt, im Lande Canaan, die Abraham kaufte sammt dem Acker von Ephron, dem Hethiter, zum Erbbegräbnisse. Daselbst haben sie Abraham begraben, und Sara, sein Weib. Daselbst haben sie auch Isaak begraben, und Rebekka, sein Weib. Daselbst habe ich auch Lea begraben in dem Gut des Ackers und der Höhle drinnen. Und da Jakob vollendet hatte die Gebote an seine Kinder, that er seine Füße zusammen auf das Bette, nahm ab und ward versammelt zu seinem Volke.

Cap. 50, 1. Da fiel Joseph auf seines Vaters Angesicht, und weinete [über ihn] und küssete ihn.

B. 1—4. Und Jakob berief seine Söhne, und sprach: Versammelt euch, daß ich euch verkündige, was euch begegnen wird in den letzten Zeiten. Kommet zu Haufen, und höret, ihr Kinder Jakob, höret euren Vater Israel etc.

1. Das ist nun die Letze, die Jakob seinen Söhnen gibt, jeglichem seine Prophezei, wie es ihm gehen soll, und nimmt sie nach einander vor, wie er sie gezeugt hat, nach ihrem Alter.

2. Der erste ist Ruben. „Meine Kraft (spricht er) und Anfang meiner Stärke“, oder Vermögens; das ist, das erste Kind, als das erste Gut, das ich überkommen habe. Daher die
\quellseite{320}
Schrift auch die Früchte auf dem Felde Kraft nennt, als droben [Cap. 4, 12.]: „Das Feld soll dir seine Kraft oder Vermögen nicht geben.“ Und Anfang meines Vermögens, das ist, ich habe zum ersten nichts gehabt, denn dich. Der Oberste im Opfer und Macht, das ist,\footnote{„das ist“ fehlt in der Erlanger.} du solltest wohl der Oberste sein, als das erste Kind, welchem von Recht das Priesterthum gebührt, das ist, das geistliche Regiment, zu predigen und die Gewissen zu regieren, dazu auch das weltliche Regiment, das er die Macht nennt. Dazu bist du geboren; aber du fährst zu leichtfertig dahin, das ist, du hast zu geringe Acht darauf, hast dich zu sicher darauf verlassen, und meinst, du wärest Junker, man müßte es alles lassen gut sein, was du thätest, hast es leichtfertig dahin geschlagen und nicht zu Sinne genommen: um der Leichtfertigkeit willen sollst du nicht der Oberste sein.

3. Also muß er hier deß entgelten, das er dem Vater zu Leide gethan hatte, und bei seiner Stiefmutter geschlafen. Denn Gott läßt nichts ungerochen, wie langsam er auch straft; wie er den Seinen nichts unvergolten läßt. Jakob hat es müssen gehen lassen und nicht gestraft, wiewohl es ihm sehr verdrossen und wehe gethan hat; aber hier wird es bezahlt in der Austheilung väterlicher Güter und Segens. Darum ist auch dieser Stamm und Geschlecht nie etwas Sonderliches gewesen, sondern fast der geringste und hinterste. Also straft Gott langsam, aber hart, fragt nicht nach der Person, und nach dem Geblüte, daß er die erste Geburt hat, und Jakobs Sohn ist. So wird diesem die Ehre genommen, die ihm sonst von Recht gebührt, und das Priesterthum auf Levi, das Königreich auf Juda und Ephraim gewendet.

\section[B. 5—7. Die Brüder Simeon und Levi, ihre]{B. 5—7. Die Brüder Simeon und Levi, ihre tödtliche Wehre sind schädliches Zeug. Meine Seele komme nicht in die Geheimniß, und meine Herrlichkeit sei nicht eins mit ihrer Versammlung etc}

4. Die zween Brüder nimmt er zuhause, und gibt ihnen auch nichts Gutes, bezahlt immer, wie ein jeglicher verdient hat. Was sie ausgerichtet haben, haben wir droben gehört, im 34. Capitel, B. 25. ff., wie sie die Sichemiter alle erwürgten, wiewohl sie ihnen Sicherheit gelobt hatten. Das war auch bisher ungestraft geblieben, aber nicht vergessen; also, daß es ihre Kinder müssen entgelten, daß sie zertheilt und zerstreuet werden in Jakob oder Israel, das ist, daß sie nicht eigene Wohnung haben im Lande. Denn Levi hat nie keine sonderliche eigene Wohnung gehabt, sondern ward hin und her geworfen, und unter den andern zertheilt. Simeon mußte sich auch behelfen im Stamme Juda, und in einem Winkel bleiben.

5. Das machen nun etliche geistlich, daß Judas, der Verräther Christi, von Simeon gekommen sei; die Priester aber, die sich über Christum versammelt haben und im Zorn erwürget, sind des Geschlechtes Levi gewesen, darum sie nun in aller Welt zerstreut sind. Aber ich lasse es bleiben im schlechten, leiblichen Verstande, nach der Historie; denn um Christi willen sind alle dreizehn\footnote{Jenaer: zwölf.} Geschlechte zerstreut; aber diese zwei haben im Volke Israel zertheilt müssen werden.

\section[B. 8. 9. Juda, du bist es, dich]{B. 8. 9. Juda, du bist es, dich werden deine Brüder loben etc}

6. Das ist ein Stück des Segens über Juda, da wird Christus mit einkommen, um welches willen es alles geschrieben ist. Hier wird nun Juda die Ehre gegeben, und allen Brüdern vorgezogen, daß sie sich vor ihm bücken sollen, und ist mit ihm so zugegangen, wie die andern Historien geben, daß er der Königsstamm ward, der unter sich hatte das ganze Israel, unter David und Salomon, schier bei hundert Jahren lang. Darnach ward es zertrennt und getheilt, daß nur zwei Stämme bei einander blieben, Juda und Benjamin, die andern kriegte Ephraim oder Joseph zu sich. Dennoch blieb das Königreich bei Juda, wie ihm zugesagt war, bis auf Christum. Aber Ephraim ist nicht geblieben, also, daß es mit einander nicht viel über zweihundert Jahre gestanden ist. Darum sagt Jakob hier zuvor, daß Juda groß werden sollte, und das Scepter behalten, bis der rechte König, Christus, käme. Also, daß der Text gewaltig zwinget, daß Christus von Juda kommen mußte, wie wir hören werden.

7. Also hat diese Prophezei zwei Theile in sich; erstlich das leibliche Königreich, welches David und Salomo und ihre Nachkommen gehabt haben bis auf Christum, in welchem das geistliche angefangen hat, und ohne Ende besteht.
\quellseite{321}

8. So sagt er nun: „Juda, du bist der Mann“, das rechte Kind von Lea, der armen Frau, „dich werden deine Brüder“, alle zwölf Geschlechter, „in Ehren haben; dazu soll deine Hand den Feinden auf dem Hals sein“, das ist, Gott wird dir Gnade geben, daß, wohin du dich wendest, wirst du deine Feinde schlagen; welches alles in David erfüllt ist. Denn wiewohl er viel Unglück daheim in seinem Hause hatte, doch wenn er auszog, schlug er alles zu Boden, daß er nie keine Schlacht verlor. Also beschreibt der Heilige Geist so lange zuvor, wie er es treiben würde; denn es auf ihn sonderlich geht, daß ihn alle Juden sollten gnädigen Herrn heißen.

9. Darum sagt er weiter: „Du bist ein junger Löwe“, unerschrocken und muthig, dazu so stark, den niemand überwältigen kann; wie die ganze Historie Davids ausweist. Item: „Du bist hoch gekommen sint dem Raube“, das ist, du hast einen Raub gethan, der hat dich hoch erhoben, nämlich an allen Königen umher; denn er alles, was sich wider ihn legte, unter sich brachte. Item: „Er hat nieder gekniet wie ein Löwe“, das ist, weil er so geraubt hat, und alles unter sich gebrochen, setzte er sich nieder, und saß stille im Regimente, war ein König, wie ein Löwe, der so gewaltig lag und saß, daß sich niemand durfte wider ihn regen. So weit redet er vom Königreiche Juda, darinne er sonderlich den David meint. Folgt nun das geistliche [Königreich].

\section[B. 10. 11. Es wird das Scepter von]{B. 10. 11. Es wird das Scepter von Juda nicht entwenden werden, noch ein Meister von seinen Füßen, bis daß der Held komme, und demselben werden die Völker zufallen. Er wird sein Füllen an den Weinstock binden, und seiner Eselin Sohn an den edlen Reben. Er wird sein Kleid etc}

10. Das geht nun alles auf Christum. Das „Scepter“ heißt er das Königreich und Regiment, das soll bei Juda bleiben. Als sollte er sagen: das ganze Volk Israel soll nicht bleiben, sondern abfallen; aber das Geschlecht Juda mußte das Regiment behalten. Und ob es nicht immerdar Könige haben würde, doch wird es zum wenigsten Fürsten haben, daß je das Regiment bleibe, nicht allein weltlich, sondern auch geistlich.

11. Denn, daß er sagt: „Noch ein Meister von seinen Füßen“, meint er das Priesteramt, Propheten und Lehrer, denen man zuhöre, bis auf den Helden, Christum, deß sollten sie sicher sein. Christum aber nennt er auf Hebräisch Schilo, welches die Juden noch nicht wissen, was es eigentlich heißt, wird auch sonst nirgend in der Schrift gebraucht, denn hier. Wir haben es gedeutet „ein Held“. Sonst heißt es einen solchen Mann, dem es glücklich von statten geht, der es frei hinaus führt, und geräth alles wohl, was er anfähet. Denn es kommt von dem Worte \unsicher{שלו}, welches heißt Glück und Wohlfahrt, als im 30. Psalm, V. 7.: „Ich sprach in meiner Wohlfahrt, oder da mir es wohl ging: Ich werde ewiglich bleiben.“ So deutet er es auf Christum mit verdeckten Worten, und bestimmt das Königreich Juda.

12. Daraus können wir nun mit den Juden handeln und ihr Ding verlegen. Denn weil sie weder Königreich noch Regiment, weder geistliches noch weltliches haben, sind dazu auch\footnote{„auch“ fehlt in der Erlanger.} aus dem Lande verstoßen, und durch die Welt zerstreuet: kurz, da ist weder Scepter noch Meister: so muß der Held oder Messias gewiß gekommen sein. Sie können auch nichts dawider aufbringen, denn es ist so ein gewaltiger Spruch, als keiner im Neuen Testamente. So ist nun unleugbar, daß, als Christus kam, war ihr Regiment aus, und das Volk ging immer unter, bis sie jetzt gar vertrieben sind, daß sie nirgend keine gewisse Stätte haben, und sind nun anderthalbhundert Jahre vergangen, daß sie müssen ohne ihrem Dank bekennen, er müsse längst gekommen sein. Also, daß auch unser Evangelium und Neues Testament gegründet ist aus diesem Texte, daß es muß rechtschaffen sein.

13. Weiter will Jakob anzeigen, was er für ein König sein solle. Da hört das jüdische Königreich auf, und geht an das geistliche. [Er] spricht erstlich also: „Demselben werden die Völker zufallen“, das ist, es soll ein solcher König, und ein solches Reich werden, das die Leute nicht mit dem Schwerte zwingen und unter sich bringen wird, wie alle andern Königreiche, sondern dem nicht allein die Juden, sondern Völker zulaufen und sich unter ihn geben. Der Königreiche ist nie keines gewesen, dem ein Volk muthwillig zugefallen sei, ohne dieses Königes: das muß so köstlich sein, daß allerlei Volk mit Lust und Liebe dazu laufe. Sagen aber die Juden, er rede von einem schlechten Herrn, dem viel Leute
\quellseite{322}
sollen zufallen und anhangen, so lasse sie es bedeuten, wie sie wollen, wir haben das unwidersprechlich, daß er soll vom Stamme Juda sein, des königlichen Geblütes, nämlich des Königs Davids, und allererst dann kommen, wenn das Scepter, das ist, ihr Königreich aufhört, und soll ein solcher Herr sein, dem nicht allein das Judenthum, sondern auch die Heiden zufallen. Was wäre es sonst geredet, wenn er sollte meinen, wie ihm die Juden allein würden zufallen, und an ihm halten, sintemal das Königreich Juda nicht soll bleiben, und schon die Juden alle zu ihm gefallen sind, und doch nicht bei ihm blieben? Was wäre es auch für Wunder, daß ihm die Juden, sein eigen Volk, zufielen?

14. Nun lasse sie sagen, wo sie irgend einen Herrn gehabt, nachdem das Scepter und leibliche Königreich aufgehört haben, dem die Völker, beide sie und die Heiden, zugefallen sind. Wir aber haben den Christum aus dem Stamm Juda, zu dem nicht allein viel Juden, sondern alle Heiden sich geschlagen haben, und willig für ihren Herrn angenommen. Denn Jacob will mit den Worten anzeigen, daß der Schilo viel größer soll werden denn das vorige Scepter und Königreich Juda. Nun konnte es nicht größer werden, wenn ihm allein die Juden zufielen; ist auch noch bei ihnen kein größeres gewesen, ohne dieses Christi; also, daß er es für einen sonderlichen Preis dieses Königes angeucht, wie es auch ein groß Wunder ist, daß alle Welt einem jüdischen Manne sollte zufallen, so doch die Heiden nie keinem Volke sind feinder gewesen denn dem jüdischen.

15. Womit wird er aber solches ausrichten, oder was wird der König für Regiment führen? „Er wird (spricht Jacob) sein\footnote{Füllen: „sein Füllen“; Wittenberger und Erlanger: eine Füllen.} Füllen an den Weinstock binden, und seiner“ 2c. Da steht sein Reich beschrieben, und scheint ja ein närrisch Beschreiben [zu] sein. Was ist das für köstlich Ding, daß des Königs, der ohne Schwert regieren und größer denn David und Salomon sein soll, höchster Preis sei, daß er seinen Esel an Reben oder Weinstöcke bindet, und sein Kleid im Weine wäscht, rothe Augen und weiße Zähne hat? Wer hat je mehr gehört einen König so loben und beschreiben? Sollte er ihn nicht billig so preisen: Er wird Land und Leute freßen, Städte und Schlösser umschlagen mit Heereskraft und Blutvergießen, daß man darinne möchte baden, und herein reiten in großem, herrlichem Pracht, und mit eitel Golde umgehen, wie der König Salomo, welcher wohl anders beschrieben ist.

16. Wenn ich einen schlechten Bürger loben wollte, und sollte ihn so loben, wie er den Esel an [den] Weinstock bände, und sein Kleid im Weine wüsche; item, von rothen Augen und weißen Zähnen, so würde jedermann meiner lachen und spotten; ja, wenn es einer von einem Narren redete, wie er nicht mehr könnte, denn den Esel an Reben binden, und ließe ihn freßen, und verderbete beide, Kleid und Wein, mit Waschen, so ginge es hin. Aber von einem solchen gewaltigen Könige, dem alle Welt zufällt, also zu reden, lautet doch gar nichts, sonderlich, wenn es so großes Ding soll geredet sein, als Gott redet. Wenn er von einem Königreiche redet, so meint er es mit Ernst, daß es etwas Großes sei.

17. Derhalben muß hier der Glaube regieren, und der Verstand gesehen werden, den Gott meint, welchen [die] Vernunft nicht erlangen kann. Die Juden machen ihm wohl eine Nase, und sagen, er meine so ein reich, überflüssig Königreich, daß der Wein so wohl gerathen solle, daß man die Esel wird lassen Trauben freßen, und Kleider im Weine waschen, daß jedermann darinne baden möchte. Aber wie\footnote{Erlanger: was.} reimt sich der Esel zum Reben, oder das Kleid im Weine? und wo bleiben die rothen Augen und weißen Zähne? oder wann ist solches je leiblich erfüllt? Darum kann solches nicht anders denn geistlich verstanden werden, wie er auch nicht ein leiblicher, sondern geistlicher König ist. Also, daß es durch und durch von dem Predigtamte geredet sei; wie wir oft gesagt haben, daß Christi Königreich nicht anders, denn durch den Heiligen Geist mit dem Evangelio regiert wird. Alle sein Schwert, Harnisch, Kraft und Macht ist das Wort Gottes, wie St. Paulus zum Römern Cap. 1, 16. sagt: „Ich schäme mich des Evangelii nicht; denn es ist eine Gotteskraft, die da selig macht alle, so daran glauben.“

18. Das ist alle Pracht und Herrlichkeit dieses Königs, daß er ein ohnmächtig, verachtet Wort führet, das vor aller Welt verdammt ist und mit Füßen getreten; aber niemand solle selig werden, er glaube denn\footnote{„denn“ fehlt in der Erlanger.} daran, und halte es für
\quellseite{323}
mächtiger denn aller Welt Kraft. Es ist wohl ein Spott und Thorheit vor ihr, sonderlich vor den heiligen Leuten; aber nichtsdestoweniger richtet er das damit an, daß die Herzen, so es fassen, gerührt vom Heiligen Geiste, fallen ab von allen andern Dingen, so die Welt für köstlich achtet, von aller\footnote{So von uns herübergenommen aus der alten Augsburger Bibel. In den andern Ausgaben: alle Abgötterei.} Abgötterei und unseren Werken, und verlassen sich allein auf Gottes Gnade. Damit wird alles zerstört auf Erden, [es] reißt uns von allen Dingen, das nicht Gott ist, welches keine Menschenkraft vermag zu thun. Denn nie keine weltliche Gewalt so groß gewesen ist, daß sie könnte die Herzen regieren, daß sie sich williglich ließen berücken und führen.

19. Also regiert er durch das Wort und mit dem Worte in den Herzen, daß sein königlich Regiment nichts anders ist, denn das Predigtamt, welches hier so lange zuvor bedeutet und verkündigt hat müssen werden von Jakob. Darum es auch St. Paulus zum Römern am Anfange [B. 2.] rühmt, daß es verkündigt sei in der heiligen Schrift. Man weiß von Christi Reiche nichts mehr denn das Evangelium; [es] regiert durch kein Schwert noch Gewalt, das Wort richtet allein\footnote{Erlanger: allen.} den Rumor an, daß sich darüber scheiden Mutter und Tochter, Vater und Sohn, wie Christus [Matth. 10, 35.] sagt.

20. Das ist es nun, das Jakob sagt: „Er wird den Esel an [den] Weinstock binden“; nämlich nichts anders, denn seine Christen binden an das heilige Wort Gottes, daß sie davon trinken werden. Was Gott geredet und geschrieben hat, das ist der edle Weinstock und Reben, wir aber sind der Esel oder Füllen. So redet er nicht von großen Leuten, sondern geringen, armen Eseln, zu deuten, daß sein Reich ein Reich der Gnaden ist, wie er auch am Palmtage anzeigte, da er einritt, und nahm das dumme Thier, den Esel dazu [Matth. 21, 7.]. Uns, die wir im Fleisch und Blut leben, heißt die Schrift Esel, wie wir auch grobe Esel sind, darum, daß der Esel gerüstet ist, zu tragen, geht aber langsam fort, daß man ihn immer treiben muß. Das sind eben seine lieben Christen auf Erden, dazu geordnet, daß wir das heilige Kreuz tragen, haben am Halse den faulen Schelm, den alten Adam, der nicht zum Guten geneigt, und nur widerstrebt; dem muß gewehrt werden, und geplagt durch das Kreuz, so lange bis er stirbt. Er muß aber am Stocke stehen, und Wein trinken, das ist, gebunden an das Wort, das uns stärkt, trunken, fröhlich und fett macht, zu leiden und sterben, daß man hindurch gehe ohne alle Furcht.

21. Eben dasselbige ist auch, das da folgt: „Er wird sein Kleid waschen im Weine, und seinen Mantel in Weinbeerblute“, das ist, in rothen Trauben; wie man gemeiniglich in den Landen rothen Wein hat, so roth wie Blut. Das ist aber ein seltsam Waschen: wenn einer wollte Kleider waschen, und nähme Blut dazu, den hielte man für unsinnig, als der das Kleid mehr besudeln wollte, denn rein machen. Welches, weil es sich nicht reimt, wie gesagt, zwingt es der Text von geistlichen Sachen zu verstehen, und ist nicht anders geredet, denn von dem Königreiche Christi, was es für Art hat. Das Kleid oder der Mantel sind wir. So steht sein Reich darinne, daß es nur ein Bad ist; denn wir sind nicht rein, werden aber rein; es ist nicht gethan, soll aber gethan werden; es ist nicht derer, die da vollkommen sind, sondern derer, die da täglich müssen zunehmen und fortfahren. Darum muß man auch immer predigen und das Wort lassen im Schwange gehen, damit immer zu fegen, sonst dürfte man keines Predigens. Denn dies Waschen und Baden geht nicht zu, denn durch das Wort; das ist der schöne rothe Wein, dadurch unsere Gewissen ohne Unterlaß gereinigt werden, wie die Apostel schreiben, man solle immer zunehmen im Glauben und Liebe, [wir] bleiben doch nimmer anders denn halbe Menschen, Fleisch und Blut steckt zu tief in uns.

22. Darum hat der König sein\footnote{Erlanger: nicht ein.} ander Amt, denn sein Kleid und Mantel, das ist, seine liebe Christenheit, die er anzieht, und sich darein schmücket; wie er sich rühmt im Propheten [Jes. 61, 10.]: Ich will dich schmücken, wie eine Braut den Mantel und das\footnote{„das“ fehlt in der Erlanger.} Band anthut, und aufsetzt. Das sind eitel starke Sprüche wider Menschengefeße. Denn wir können uns nicht mit Menschenlehre waschen, sondern man muß sie waschen im Worte; das ist, man wird ihnen predigen, und durch das Predigen werden sie glauben, dadurch werden sie rein werden, und immer je mehr und mehr gefegt. Weiter:
\quellseite{324}

\section[B. 12. Seine Augen sind röthlicher denn Wein,]{B. 12. Seine Augen sind röthlicher denn Wein, und seine Zähne weißer denn Milch}

23. Das wäre billiger eines Weibes Lob, denn eines Königes; der steht es wohl an, daß sie schön sei. Eines Mannes Schmuck wäre besser Schwert und Harnisch. Aber weil er so durch Reden in Narrheit kommt, zwingt er uns, von der Narrheit in seine Weisheit zu kommen. „Seine Augen sind röthlich“, spricht er, wie starker rother Wein, und setzt wieder einander „rothe Augen“ und „weiße Zähne“.

24. Die Augen sind, die da vorstehen dem Worte GOttes; denn er beschreibt noch immer das Predigtamt und Wort, damit er die Christenheit nährt und führt, und will, daß sie allein damit umgehe, und keine andere Lehre noch Regiment wisse. Die Prediger aber, die es treiben, führen, und in der Hand haben, sind die Augen und Zähne, das ist, daß ein Prediger zwei Stücke soll haben, Gesetz und Gnade predigen. Rothe oder braune Augen sind zornig und ernst; das sind Christi Augen, welcher ist ein Bischof und Aufseher, der keine andere Predigt noch Regiment zuläßt.

25. [Das] ist sonderlich geredet auf die Lehre, denn es liegt nicht so viel am Leben. Wenn man die Lehre rein behielte, könnte man allen Sachen rathen. GOtt gebe nur Gnade, daß man recht predige; ob wir nicht alle recht leben, müßte man Geduld tragen. Denn es geht also: Wo man recht predigt, da geht das Leben schwach hernach, faul und langsam; wiederum, wo man die Leute verführt mit falscher Lehre, geht es gar nicht hernach. Denn wo die Lehre nicht recht ist, kann das Leben, welches sich nach der Lehre richtet, auch nicht recht sein.

26. Darum sind „die Augen“ Prediger und Apostel; die sind roth wie starker Wein, das ist, scharf, daß sie strafen alles, was nicht reine Lehre ist, das ist, was nicht Glauben und Liebe lehrt. „Die Zähne“ aber sind, die da beißen, und sind doch weiß, wie Milch, die da süß ist. So wunderliche Rede führt er vor großer Lust an dem Reiche Christi. Das sind nun auch dieselben Prediger, die da beißen und strafen, wenn man schon im Glauben ist, daß man die reine Lehre erhalte, tadeln durch das Gesetz aller Menschen Ding, und lassen dann aufkommen den lautern Glauben. Danach, wenn Christen worden sind im Glauben, thun sie nichts, denn strafen immerdar, und sind doch süß, rein und lauter wie Milch; nicht giftig noch falsch, die nicht das Ihre suchen, begehren nichts denn der andern Nutzen und Besserung; sind wohl scharf, sie meinen es aber gut; heucheln nicht, sondern handeln aufrichtig, recht und schlecht. Das heißen Zähne wie Milch, candidi dentes, dagegen heißt man lividos dentes, faule und giftige Zähne, die nicht rein sind. Die Kleinen strafen und schelten auch, meinen es aber herzlich und brüderlich; wie St. Paulus die Corinther, item, die Galater auch straft.

27. Siehe, durch solche wunderliche Rede, so verblümt und geschmückt, zeigt er so meisterlich an das ganze Königreich Christi, wie es gar nicht stehe in leiblichen Dingen, sondern nur darum zu thun ist, daß die Lehre allein recht sei, und führe auf den lautern Glauben, dazu auch die Prediger, so darüber halten. Wenn nur solches recht geht, kann GOtt wohl Geduld haben am Leben. Es ist nicht zu thun so viel um die Zuhörer als um die Lehrer; denn die Zuhörer sind bald geurtheilt, ob sie recht, oder nicht, leben.

28. Darum kann keine größere Gnade und theurere Gabe widerfahren, denn daß GOtt rechte Prediger schickt, die es mit Treuen meinen, und den falschen Lehrern wehren. Wiederum, keine größere Plage noch Herzeleid, denn wenn er die rechten hinweg nimmt, läßt falsche Lehrer mit blauen, giftigen Zähnen aufkommen. Da lasset uns zusehen und davor hüten. Sie sind schon vorhanden, und reißen ein; ja, auch unter uns sind ihrer, die es bereits sind, oder werden es werden. Wir haben das Evangelium, wissen, was wir thun, wie wir glauben und lieben sollen; aber wir sind bald gewöhnt und gestorben; so ist es gar bald geschehen, daß alles wieder voll falscher Prediger wird. Darum lasset uns GOtt danken, und bitten, daß es rein bleibe, und wehre den falschen, verführerischen Geistern.

29. Das ist der edle Text von dem Reiche Christi, wie es gehen soll mit dem rechten Verstande vom heiligen Evangelio, welches nimmer genung zu preisen ist; [es] ist so groß, daß sich ein christlich Herz davor wundern muß; [es] thut und schafft so viel, daß [es] niemand genung ermessen kann, davon wir immerdar gepredigt haben, auch nichts anders zu predigen haben. Darum hat auch der heilige Patriarch
\quellseite{325}
so lustig davon geredet, auf das allerbeste er konnte und gewußt hat.

30. So schließe nun aus diesem Texte, daß Christus gekommen ist; denn die Predigt ist schon längst angegangen, die die Leute so kühn und fröhlich macht, daß sie sich nicht fürchten vor Tod noch Teufel, und geschickt sind, das heilige Kreuz zu tragen, dazu seine, reine und scharfe Predigt haben; wie es erstlich durch die Apostel und Märtyrer erfüllt ist. Weil nun die Predigt angegangen ist, können wir kein beweisen, daß wir den König haben, weil ihm so viel Völker zugefallen sind, die um seinetwillen fröhlich das Blut vergossen haben und solche rechtschaffene Lehre bewiesen.

31. Siehe, solches hat Gott dem Juda zuvor kund gethan, und verheißen. Das ist auch das Hauptstück des ganzen Capitels. Denn wiewohl er dem ganzen Geschlechte, allen zwölf Stämmen, jeglichem seinen Segen gibt, sind sie doch alle um des Hauptsegens willen geschehen und geschrieben, daß Christus kommen sollte vom Stamm Juda. Darum, was nun hernach folgt, ist eitel zeitlich Ding, so den andern verkündigt ist, wie es ihnen gehen solle, wie auch den vorigen zweien, Simeon und Levi.

32. Doch lasset das dahin, wie Gott alle Dinge aus lauter Gnaden gibt, und niemand nichts überkommt aus seiner Arbeit, Mühe und Vorsichtigkeit, sintemal er diesen so lange zuvor sagt, wie es ihnen gehen solle, ehe sie geboren werden, daß man lerne, was jedermann widerfährt, sei also beschlossen von Gott von Ewigkeit, beide, wie ein jeglicher leben und sterben solle.

33. Es ist nicht Menschen Vernunft noch That, sondern Gottes Wille und Ordnung, es sei gut oder böse. Wir lassen uns dünken, es gehe nach unserm Träumen, wenn etwas geschieht, meinen, wir haben weislich gethan, oder genarrt. [Wir] sollten uns billig schämen, wenn wir diesen Text recht ansehen, daß wir meinen, es gehe, wie wir es machen, nicht, wie er es führt. Wenn ich denn das weiß, daß alles geht, wie er es ordnet, und nichts begegnet, es ist von ihm so versehen, und geschickt: was wollen wir denn daraus machen? [Wir] müssen sagen, du wollst es so haben. Weil wir das nicht thun, so murren wir darum, und wollen Gott meistern. So wollen wir nun die andern Segen vollends überlaufen, wiewohl sie nicht schwer sind; die Sophisten aber sind darein gefallen, und haben es mit ihrem Deuten verfinstert.

\section[B. 13. Sebulon wird am Anfurt des Meeres]{B. 13. Sebulon wird am Anfurt des Meeres wohnen, und am Anfurt der Schiffe, und stoßen an Sidon}

34. Was darf Jakob davon unnütze Worte machen? „Anfurt des Meeres und grenzen an Sidon“, ist nicht mehr denn so viel verheißen dem Sebulon, er solle die Grenze inne halten, die da ist am Meere, und langet bis gen Sidon. Das wissen die Landfahrer, wie und wo das große Meer, darauf man schifft von Venedig gen Jerusalem, an das Land stößt, da man anfährt. Wie, und wann sie nun das Land eingenommen haben, liest man im Buche Josua hin und wieder.

35. Das ist nun gar ein zeitlich Ding, doch ist das geistlich darinnen, wie gesagt ist, daß Gott nicht allein gibt, was der Seele gehört, sondern auch ordnet und schickt leiblich gut Regiment und alle Dinge, daß nichts ohne seine Ordnung sich begeben kann. Daß die Fürsten von Sachsen dies Land haben, ist nicht ihre Gewalt, sondern Gottes Ordnung und Wille, der da sagt: So weit sei Herr, und nicht weiter. Also, wie Paulus sagt in den Geschichten [Apost. 17, 26.], daß er einer\footnote{Jenaer und Erlanger: einem.} jeglichen Herrschaft abzirket, eben und dürr abmißt, wie weit sie regieren soll; Trotz aber, daß sie mehr ein Haarbreit könnte einnehmen,\footnote{In den alten Ausgaben: „trotz aber daß sie mehr ein Haarbreit kund einnehmen“. — „über“ scheint uns ein Druckfehler zu sein statt: „aber“, was wir gesetzt haben.} daß man bekenne, was jedermann hat, er von Gottes Hand habe und nehmen müsse, der aller Dinge Gewalt hat, ob es gleich die Welt nicht erkennt; wie es auch die Juden nicht erkannt haben, meinten, sie thäten es durch ihr Schwert. So spricht er droben: Ja, wenn ich es nicht thäte, könntest du nicht ein Haarbreit haben.

36. Wir Christen müssen das wissen, so wir glauben, daß er der Schöpfer ist Himmels und der Erde, daß ohne sein Wissen und Willen niemand ein Haar vom Haupte fällt [Luc. 21, 18.] und nicht einen Bissen Brodes haben könne, er gebe es denn. Es ist nimmer genugsam zu lernen, ob es gleich leicht zu sagen ist. Wenn wir es glaubten und für wahr hielten, würden wir für nichts sorgen, noch uns fürchten, sondern
\quellseite{326}
ihn machen lassen, und alles von seiner Hand warten. Weil wir aber nicht glauben, denken wir, er sei weit von uns, wolle, daß wir selbst sorgen und denken, wie wir kriegen, was uns noth ist. Das sei nun dem Sebulon geschenkt, weiter soll er nicht greifen noch langen.

B. 14. 15. Issachar wird ein beinern Esel sein, und sich lagern zwischen die Grenzen. Und er sähe die Ruhe, daß sie gut ist, und das Land, daß es lustig ist, und hat seine Schultern geneiget zu tragen, und ist ein zinsbarer Knecht geworden.

37. Das ist nicht gepredigt für den gemeinen Mann, dient für die Tiefsinnigen, die in der Schrift studiren. Diesem Stamme gibt er den Segen: Er soll wohl ein fein Land haben, aber ein Knecht, und zinsbar sein. Das Land ist gelegen mitten in Canaan, auf einem fruchtbaren Boden, wie Sebulon und Naphtali auch; aber es liegt in der Grenze, da die Heiden an stoßen; darum sollte es immerdar angefochten werden, und dem Assyrer Könige Schoß und Zinse geben. [Er] will also sagen: Du siehst wohl, und wirst es inne, daß [es] ein köstlich Land ist; aber du wirst nicht Herr darin sein. Als sollte er sprechen: Du wirst zu fressen und saufen und alle Fülle haben, aber nicht Gott dankbar sein, darum wirst du es müssen andern Leuten geben, und fressen lassen. Das sieht man nun in der Historie, wie sie geschlagen sind von Königen, haben das Geschlecht nie errettet, [es] ist immer geblieben unter den Königen von Syrien und Assyrien. So ist es wahr worden, daß er ein beinern Esel wäre, das ist, der große Last tragen und unterdrückt werden müßte, mehr Bürden tragen, denn er vermögen würde.

B. 16—18. Dan wird Richter sein in seinem Volke, wie ein Geschlecht in Israel. Dan wird eine Schlange werden auf dem Wege, und eine Otter auf dem Steige, und das Pferd in die Fersen beißen, daß sein Reiter zurück falle. Herr, ich warte auf dein Heil.

38. Diesen Text müssen wir darum handeln, daß wir begegnen dem losen Geschwätze und Fabeln. Denn man trägt ein Buch vom Antichrist um, darin gesagt wird, er solle vom Stamm Dan geboren werden, darum, daß geschrieben ist: „Er soll sein eine Schlange und Otter auf dem Wege.“ Es ist aber lauter Narrenwerk. Der Antichrist soll nicht von Juden kommen; wir haben schon den rechten vorhanden, davon Paulus [2 Thess. 2, 3.] sagt, den Pabst zu Rom. So ist aber wohl gesagt, daß die zwölf Geschlechte zerstreuet sollten werden in alle Welt, und nimmer aufkommen, noch [es] dahin bringen alle auf einen Haufen, daß sie so stark sollten werden, daß [es] möchte ein Fürstenthum heißen. Die Schrift steht dürr da, daß sie sollen zunichte werden, und kein Regiment kriegen; wie der 110. Psalm, V. 6., prophezeiet: „Gott wird den Kopf der Feinde zerschmeißen.“

39. Des Antichrist Reich aber, weil es eine solche Macht ist, die alle Welt überwindet, kann es davon nicht kommen, sondern muß einen andern Ursprung haben, wie St. Paulus sagt. Wenn das römische Reich hinweg ist, so sehe man sich vor, dann wird es aufkommen.

40. So sagt er nun: „Dan wird auch ein Richter werden, wie die andern.“ Welches ist erfüllt in dem starken Helden Simson, in der Richter Buche; der ist eben der, davon er hier redet: Die Schlange auf dem Wege, und Otter auf dem Steige, so das Pferd in die Fersen beißt, daß der Reiter zurück fällt; das ist, daß sie hergeht, und macht Reiter und Roß zu Schanden; wie er gethan hat mit den Philistern. [Er] war ein einiger Mann, ohne allen Zeug, er schlich und griff sie an, und trieb sie alle zurück. So ist dies die Allegorie dieser Historie, so Simson gethan hat, wie oft er mit eigener Person eingerissen, und die Philister geschlagen hat, daß ihm niemand konnte widerstehen. Davon besehe das 15. und 16. Capitel im Buche der Richter.

41. Daß nun Jakob dazwischen mengt, und spricht: „Herr, ich warte auf dein Heil“, redet er auf Christum; als sollte er sagen: So wird es zugehen im Volke, so lange wird die Schlange regieren, wollte Gott, dein Christus wäre bald hier!

B. 19. Gad, gerüstet, wird das Heer führen, und wieder herum führen.

42. Davon lies die Historie des Buches Josua, wie sie vor Israel hergezogen sind, und an der Spitze gewesen.

B. 20. Von Asser kommt sein fett Brod, und er wird den Königen niedliche Speise geben.

43. Asser soll den Segen haben, daß er in einem köstlichen Kornlande soll sitzen, und reich werden.
\quellseite{327}

\section[V. 21. Naphtali ist ein schneller Hirsch, und]{V. 21. Naphtali ist ein schneller Hirsch, und gibt hübsche Rede}

44. Davon lies auch die Historie im Buche der Richter [Cap. 4, 4. ff.] von der Prophetin Debora und Barak aus diesem Geschlechte, wie sie den großen königlichen Hauptmann\footnote{So die Jenaer. Wittenberger und Erlanger: König.} Sisera schlugen mit einem geringen Heere, kamen kurz hinter ihm her, so behende und schnell, wie ein Hirsch. Davon macht sie einen hübschen geistlichen Gesang, daselbst [Cap. 5.] beschrieben. Das sind die zehn Söhne von Lea und den zwei Mägden.

\section[V. 22—26. Das holdselige Kind Joseph, das holdselige]{V. 22—26. Das holdselige Kind Joseph, das holdselige Kind anzusehen, die Töchter treten einher im Regimente. Und wiewohl ihn die Schützen erbittert, mit ihm gezankt, und ihn gehaßt haben etc}

45. Daß man dies verstehe, habe ich vorgesagt, daß Gott das Land und Volk in zwei Theile getheilt hat nach den zwölf Brüdern, und jeglichem sein Theil ward gegeben. Aber das Geschlecht Levi zog er aus. So macht er nun, daß das König-Geschlecht kam auf Juda, bis auf Christum, welcher ein geistlich Reich daraus machte; doch ward das Königreich getheilt in zwei Theile; ein Stück führte Juda, das andere Joseph oder Ephraim. Juda behielt nun zwei Geschlechte, die andern zehn fielen alle dorthin; jenes war das geringste, aber doch das beste; dieses war das gewaltigste, aber nicht so gut. Denn Juda hatte viel fromme Könige, aber Ephraim hatte nie seinen frommen, darum sie auch Gott strafte und aus tilgte, daß sich zuletzt die zehn Geschlechte gar zerstreuten, und nie wieder zusammen kamen.

46. So redet nun hier Jakob von dem großen Königreiche Ephraim, aber nichts, denn vom zeitlichen Regimente, daß es ein fein Volk und reich sei, und ihm wohl gehen solle um deswillen, daß unter dem Volke viele waren, die zu Gottes Reiche gehörten. Denn es hat auch viel treffliche Propheten gehabt, darum mußte es auch Glück haben, die Könige wären fromm oder nicht; der Herr ist aber ein wenig finster.

47. „Joseph (sagt er) ist das liebe Kind“, bindet die zwei Söhne, Ephraim und Manasse, an einander an Eine Person; „und ist ein holdselig Kind“, das ist, das Volk von Joseph wird ein fein Volk sein, nach dem zeitlichen Regimente und Wesen; nicht seiner Person halben, daß er ein schöner Knabe war, sondern seines Stammes, Blutes und Fleisches. Denn es ist ein solch schön Reich gewesen, daß sich andere Könige verwundert haben, von geschickten, weisen Leuten und seinen Propheten. Die Töchter aber heißen in der Schrift: Städte und Flecken des Landes; als im Matthäo [Cap. 21, 5.] steht: „Saget der Tochter Zion“, das ist, Jerusalem. So will er nun sagen: Die Städte in dem Reiche sind wohl geordnet im Regimente, [es] ist fein im Schwange und gutem Frieden gegangen, und wohl gediehen, welches ein großer Schatz und Gnade ist in einem Lande. Denn was Gott lobt als ein fein Regiment, muß nicht etwas Geringes sein.

48. Solche Gnade soll nun das Reich vor andern haben, daß [es] ein köstlich, wohlgeordnet Regiment habe. Und nicht allein das, sondern auch, wenn sich gleich andere Könige an sie machten, als, von Syrien und Assyrien, mit ihm zu streiten und unter sich zu bringen, soll es nicht Noth haben, sein Bogen und Arm soll feste und stark bleiben, und immer obliegen; wie die Historien von [den] Königen zeigen. Wiewohl sie oft angefochten sind, haben sie doch gewonnen und das Regiment erhalten, sind aber fest geblieben und ermannet; nicht durch ihre Kräfte (denn sie waren oft kaum das zehnte Theil so stark und viel, als die Feinde), sondern „durch die Hände (sagt er) des Mächtigen in Jakob“. Also liest man vom Könige Ahab, der doch ein Bösewicht war vor Gott, aber vor der Welt ein feiner Mann, wie er auf einen Tag hunderttausend Mann mit einem kleinen Häuflein schlug [1 Kön. 20, 27—29.].

49. Weiter „sollen aus ihnen kommen Hirten und Steine in Israel“. Das sind die großen, feinen Propheten, Elias, Elisäus, Micha, und andere; das waren die Edelsteine, darauf sie sich sollten verlassen.

50. Weiter will er sagen, dieser Segen widerfährt dir nicht um deines Verdienstes willen, oder um deiner Väter willen; [es] ist ein lauter Segen Gottes, was du hast. „Mit Segen vom Himmel, und von der Tiefe“; das ist, die Sonne wird dir scheinen, regnen,\footnote{In den alten Ausgaben: regen.} thauen, kalt und warm sein, wie es der Erde nöthig ist: [es] soll dir alles dienen, auch die Tiefe oder Wasser soll
\quellseite{328}
das Land feuchten, und Dampf geben, daß der Regen komme vom Himmel. Dazu auch „mit Segen an Brüsten und Bäuchen“, also, daß alles schwanger sein, Frucht bringen, und säugen soll, was nur Bäuche und Brüste hat. So ist mit bunten Worten alles verfasset; was fruchtbar sein kann von Menschen und Vieh, soll sich alles mehren und genug haben. [Es] ist so viel geredet: Das Kinder-Tragen, Vieh-Gerathen und Gedeihen, dazu das Getreide auf dem Felde, ist alles nicht nur Ding, sondern GOttes Gnade und Segen ist es.

51. Item: „Der Segen, deinem Vater (spricht er) und meinen Voreltern verheißen, gehet stark.“ Mir und meinen Vätern ist verheißen der Segen, daß wir das Land sollen einnehmen und besitzen, mit aller Fülle des Segens, daß [es] ein köstlich Königreich werde, alles genug und wohlgerathen, und ordentlich gehe mit Streit und Frieden. „Nach Wünsche der Hohen in der Welt“, das ist, wie alle Könige und Fürsten, und was da hoch ist, möchten wünschen. Also streckte sich dieser Segen nicht weiter, denn auf das leibliche Königreich. Und ist alles darum so geschrieben, daß wir wissen sollen, wie es von Gott zuvor alles so geordnet ist, daß er es gibt, und wir nicht, und sich niemand rühme, es sei seine Arbeit, Geschicklichkeit und Vorsichtigkeit. Es muß alles sein Segen sein. Ist es nicht Segen, so gibt es und schafft nichts; ist es aber gegeben, so ist es lauter Gottes Segen. Also zeucht Gott alle Güter zu sich, was wir haben an Leib und Seele.

52. Auf das letzte sagt er: „Aus Joseph sollen Häupter werden, und oberste Nasiräer unter seinen Brüdern.“ Das ist gewesen Elias und Elisa, [die] haben regiert wie die Häupter; dazu Nasiräer, davon Moses schreibt [4. Buch, Cap. 6, 2. ff.]; das sind leibliche Heilige gewesen, heilige Propheten, geordnet das Volk zu regieren und lehren, als auch der Prophet Samuel war. Also ist das Geschlecht gesegnet, daß es soll wohl stehen von Regimente, seinen Leuten, und allerlei Fülle. Wiewohl es daran gefehlt hat, daß sie nicht recht Gott dienten; dennoch den leiblichen Segen gekriegt haben, um dieses Patriarchen willen.

\begin{quote}
V. 27. Benjamin ein reißender Wolf; des Morgens wird er Raub fressen, aber des Abends wird er den Raub austheilen.
\end{quote}

53. Was der Text will, weiß ich nicht allzu gewiß. Es möchte verstanden werden von dem Könige Saul, oder auch St. Paulo, denn sie beide dieses Stammes gewesen sind. Aber weil der Segen leiblich lautet, meine ich auf das einfältigste, er rede von dem Könige Saul, welcher der erste König gewesen ist, und gestritten hat wider Amalek, und ihn nieder gelegt, und den Raub ausgetheilt hat [1 Sam. 14, 47. 48.]. Wer es geistlich will machen, und auf St. Paulum ziehen, mag es thun, also, daß er ein reißender Wolf ist gewesen; erstlich, als er die Christenheit verfolgt hat, Apost. 8, 3. 9, 1. ff., darnach aber bekehrt worden, den Raub gerissen und ausgetheilt, das ist, die Seelen dem Teufel und der Welt geraubt, und zu Christo gebracht hat, und viel Gutes geschafft in der Christenheit, ja, fast das größte Theil der Welt allein bepredigt\footnote{So die Jenaer. In den andern Ausgaben: gepredigt.} und Christen gemacht. Darum mir es auch nicht übel gefällt, den Text von ihm zu verstehen. Denn Gott hat insgemein einen Stamme Einen Mann, oder zwei oder drei gegeben, durch welche das ganze Geschlecht gepriesen würde. Darum soll man diese Segen nicht auf allerlei Personen, sondern auf Einen oder zwei ziehen. Also liest man von niemand aus dem Stamme Benjamin, der etwas Sonderliches gewesen oder ausgerichtet habe, denn von dem Könige Saul, und von dem Apostel Paulo. Item, von den Bürgern zu Gibea, wie im Buch der Richter [Cap. 20, 15.] steht.

54. So sei nun dies die Summa dieses Capitels: Wie alles, das da geschieht, durch Gottes Verordnung geschieht, daß es so gehen muß und nicht anders, man schlage es vor, wie man wolle, und niemand durch seine Vernunft oder Arbeit erlangen, auch niemand hindern kann, sondern alle Welt das Maul halte, und rühme sich nichts von ihrem Gute. Hast du etwas, so wisse, daß es allein sein Segen ist; wird dir was genommen, so ist es alles, bei einem Heller, durch ihn geordnet und versehen, daß [es] also gehen muß.
\quellseite{329}

\section{Das fünfzigste Capitel}

V. 2—14. Und Joseph befahl seinen Knechten, den Aerzten, daß sie seinen Vater salbeten. Und die Aerzte salbeten Israel, bis daß vierzig Tage um waren; denn so lange währen die Salbetage, und die Egypter beweineten ihn siebenzig Tage. Da nun die Leidetage aus waren, redete Joseph mit Pharao Gesinde, und sprach: Habe ich Gnade vor euch funden, so redet mit Pharao, und sprechet: Mein Vater hat einen Eid von mir genommen, und gesagt: Siehe, ich sterbe, begrabe mich in meinem Grabe, das ich mir im Lande Canaan gegraben habe. So will ich nun hinauf ziehen, und meinen Vater begraben und wieder kommen. Pharao sprach: Zeuch hinauf und begrabe deinen Vater, wie du ihm geschworen hast. Also zog Joseph hinauf, seinen Vater zu begraben. Und es zogen mit ihm alle Knechte Pharao, die Aeltesten seines Gesindes, und alle Aeltesten des Landes Egypten. Dazu das ganze Gesinde Josephs, und seine Brüder, und das Gesinde seines Vaters; allein ihre Kinder, Schafe und Ochsen ließen sie im Lande Gosen. Und zogen auch mit ihm hinauf Wagen und Reitige, und war ein fast großes Heer. Da sie nun an den Fluten Atad kamen, der jenseit dem Jordan liegt, da hielten sie eine sehr große und bittere Klage, und er trug über seinen Vater Leide sieben Tage. Und da die Leute im Lande, die Canaaniter, sahen das Leidetragen im Flecken Atad, sprachen sie: Die Egypter tragen da große Leide. Daher heißt man den Ort, der Egypter Leide, welcher liegt jenseit dem Jordan. Und seine Kinder thaten, wie er ihnen befohlen hatte, und führten ihn in das Land Canaan, und begruben ihn in der zwiefachen Höhle des Ackers, die Abraham erkauft hatte mit dem Acker, zum Erbbegräbniße, von Ephron, dem Hethiter, gegen Mamre. Also zog Joseph wieder in Egypten mit seinen Brüdern, und mit allen, die mit ihm hinauf gezogen waren, seinen Vater zu begraben, als sie ihn begraben hatten.

1. Das ist ein ehrlich Begräbniß gewesen der heiligen sechs Patriarchen, Abraham, Isaak, Jakob mit ihren Weibern, die doch nie eines Fußes breit gehabt haben in dem Lande, so ihnen von GOtt verheißen war, wie St. Stephanus [Apost. 7, 16.] sagt. Warum haben sie denn eben an dem Orte wollen begraben werden, daß auch Jakob so theuer befiehlt, man soll je seinen Leib hinauf führen und begraben in seiner Väter Grabe? Es ist allein geschehen, die Verheißung zu bestätigen, auf daß die Juden hernach desto freudiger wären, die Heiden auszurotten, auf daß sie das Land, da ihre Väter lagen, einnehmen.

2. Denn wir sind so geschickt, was GOtt zusagt und gelobt, kann man uns nicht genug einbilden und vorbläuen, daß wir es fest glauben. Das macht, daß GOtt so wunderlich regiert. Denn sie waren nun in Egypten, in fremdem Lande, und bald nach Josephs Tode wurden sie so untergedrückt, daß nicht möglich war vor der Welt, daß sie das Land sollten einnehmen. Und waren über dreißig Könige in dem Lande Canaan, große, gewaltige Leute, die waren also gesessen, daß nicht möglich war, sie auszutreiben, sonderlich von so geringem, schwachem Volke; wie man spricht: Der Hahn ist nicht gut von seinem Miste zu treiben, und Ein Mann vermag mehr in seinem Hause, denn vier draußen. Derhalben wollten die Väter allda begraben werden, ihren Glauben zu bezeugen, und die Kinder zu stärken. Sonst hätte man gesagt: Ja, wie sollen wir es glauben, haben sie doch selbst nicht geglaubt. Aber solches zu wehren, hat sich Jakob mit solchem Gepränge in das Land wollen führen lassen, als der da gewißlich glaubte, daß er von GOtt in das Land geführt würde.

3. Also ist Gottes Werk und Wort allezeit gerichtet, daß [es] Vernunft nicht glauben noch begreifen kann. Denn, wie sollte das die Vernunft können ermessen? Dieser stirbt im fremden Lande, daß er auch sein Begräbniß muß kaufen; wie fein ist er denn ein Herr des Landes? Es scheint so gar unwahrhaftig, daß die Vernunft spotten und lachen muß. Aber ist es Gottes Wort, so muß es geschehen, es müßte eher Himmel und Erde untergehen.

4. Darum braucht Moses so viel unnützer Worte, wie uns dünkt, daß er ihren Glauben
\quellseite{330}
anzeige, wie sie nichts Eigenes im Lande gehabt haben, daß sie sollten besitzen, und gestorben sind auf dasselbe Wort Gottes, auch nach dem Tode so fest geblieben, daß man sie darinne begraben mußte; also, daß nicht mehr denn das Wort da ist, und sie in dem bloßen Worte das Land einnehmen. Der Leib läßt sich in das Land führen und begraben; aber die Seele hält sich an das Wort, und bleibt darinne in ewigem Leben; welches die Epistel an die Hebräer [Cap. 11, 20. f.] fein heraus gestrichen hat.

5. Also haben wir den Jakob auch ausgerichtet; in welcher Legende wir nichts gesehen haben, denn eitel Uebung des Glaubens im Kreuz und Leiden, von Anfang an, daß er durch so viel Jammer und Herzeleid wohl gekocht und mürbe ist worden, daß er des Lebens müßte müde werden, und aufs letzte hin verdrossen ist, wie ein Licht. Denn die, so versucht sind und gebrochen auf Erden, denen wird der Tod nicht bitter, aber die viel guter Tage haben, und ohne Anfechtung hingehen, sterben gar schwerlich und ungerne. Denn es gar hart ist, so lange zu harren, bis es an die Züge geht, und dann allererst Lust zu sterben gewinnen, wenn man zuvor nicht wohl gewohnt ist, allerlei zu leiden.

6. So sagt nun die Schrift: „Er ist versammelt zu seinem Volke“, wie er auch droben [Cap. 25, 8.] von Abraham geredet hat. [Es] ist eine hebräische Weise zu reden, und so viel gesagt, nicht, wo er hingefahren sei, sondern er ist gekommen zu den Todten, da sie begraben sind, wie wir auf unsere Weise sagen: Er ist auf den Kirchhof gekommen. Aber da ist heimlich die Auferstehung von [den] Todten angezeigt, daß sie Gott wiederum wollen auferwecken; daher die Schrift den Tod auch einen Schlaf heißt [Joh. 11, 11. Marc. 5, 39.]. Denn die da schlafen, haben die Hoffnung, daß sie wieder sollen auferstehen. Also auch sind die Väter versammelt zu ihrem Volke, als auf die Gott ein Auge hat, und will sie zu seiner Zeit auferwecken.

7. Daß nun Moses schreibt, wie Jakob gesalbt und begraben ist, scheinet, als seien es eitel unnütze Theidinge; [es] ist aber darum geschrieben, daß es rüchtig wäre bei ihren Kindern und Nachkommen, ihren Glauben zu stärken. In dem hat nun Joseph des Landes Weise gehalten, daß er ihn salben läßt vierzig Tage lang; wie sich den Christen gebührt immer zu stellen nach anderer Leute Sitten und Brauch, und nicht ihrem eigenen Sinne nach zu leben. Nun haben die Egypter die Weise gehabt, ihre Todten zu salben mit edlen, köstlichen Salben, oder Würze und Specerei, so dasselbe Land trägt, welche auch hernach die Juden von ihnen genommen haben; wie auch Christus gesalbt wird mit Myrrhen und Aloe [Joh. 19, 39. 40.], daß auch etliche ihre Eltern pflegten zu behalten in einem Crystall, wie den Kaiser Titum zu Rom. Sonst könnte der todte Leib nicht gewähret haben, sonderlich in dem heißen Lande; denn sie haben ihn einen Tag oder hundert müssen führen, ehe sie zum Begräbnisse gekommen sind.

8. Item, daß Joseph Urlaub nimmt von dem Könige Pharao, den Vater zu begraben, will nicht ohne seinen Willen von seinem Amte ziehen, ist darum geschrieben, daß ein jeglicher der Obrigkeit gehorchen soll, und nichts ohne ihren Willen thun, sonderlich die im Amte sitzen. Denn er sich hier als ein Knecht herunter läßt, und bittet des Königes Gesinde, alle, die unter ihm waren, daß sie für ihn sollten bitten. Also hoch will Gott die weltliche Obrigkeit geehrt haben. Dawider sich der Papst mit seiner Rotte aufgeworfen, gelehrt und gethan hat, und sich von aller weltlichen Oberhand ausgezogen.

\begin{quote}
V. 15—21. Die Brüder aber Josephs fürchteten sich, da ihr Vater gestorben war, und sprachen: Joseph möchte uns gram sein und vergelten alle Bosheit, die wir an ihm gethan haben. Darum ließen sie ihm sagen: Dein Vater befahl vor seinem Tode und sprach: Also sollt ihr Joseph sagen: Lieber, vergib deinen Brüdern die Missethat und ihre Sünde, daß sie so übel an dir gethan haben. Lieber, so vergib nun diese Missethat uns, den Dienern des Gottes deines Vaters. Aber Joseph weinete, da sie solches mit ihm redeten. Und seine Brüder gingen hin, und fielen vor ihm nieder und sprachen: Siehe, hier sind wir, deine Knechte. Joseph sprach zu ihnen: Fürchtet euch nicht, denn ich bin unter Gott. Ihr gedachtet Böses über mich; aber Gott hat es zum Guten gewandt, daß er thäte, wie es jetzt am Tage ist, zu erhalten viel Volks. So fürchtet euch nun nicht, ich will euch versorgen und eure Kinder. Und er tröstete sie und redete freundlich mit ihnen.
\end{quote}

9. Das ist abermal ein seiner, lieblicher Text, und groß Exempel, uns zu trösten. Da stehen die Brüder, bekennen und beichten ihre Sünde,
\quellseite{331}
und ziehen an GOtt für einen Vater, wie die rechten Christen. Joseph aber ist so voll Liebe, daß ihm das Herz erwarmt, und sich nicht enthalten kann, er muß weinen; so gar hat er der Missethat vergessen, so die Brüder an ihm gethan haben. Wer den Spiegel könnte in sich bilden, der würde seinen Feinden leichtlich vergeben.

10. Ja, noch mehr sagt er: „GOtt hat es zum Besten gewandt, euch und mir, und viel Volks\footnote{Oben im Text und in der Bibel: zu erhalten viel Volks.} zugut gethan.“ Er lobt ihre That nicht, sondern bekennt, sie haben es böse gemeint, ihn zu verderben; aber er lehrt uns, wie GOtt aller unserer Feinde Herzen und Gedanken in seiner Hand hat, welche, ob sie gleich böse sind, kann er sie doch so umkehren, daß es gerade wider ihre Meinung hinaus geht. Wer das glaubt, der kann seine Feinde leichtlich überwinden, darf sich vor niemand fürchten, denn er weiß, daß GOtt alle Dinge regiert und meistert; darnach, daß er einen solchen GOtt hat, der ihm alles Unglück wolle zu Glück wenden.

11. Darum, bist du ein Christ, so laß die Welt feindlich zürnen, und alles dir nehmen, so sie kann, aber tröste dich deß, je ärger sie es meinen und gedenken, je näher GOtt bei dir steht, und meint es auf das allerbeste. Also, daß der Christen Leiden und Verfolgung allein dazu dient, daß wir darauf trotzen, und wissen, GOtt wende es zum Besten. Das ist es, daß Christus sagt im Johanne [Cap. 12, 24.]: „Wo nicht das Weizenkorn in die Erde fällt und erstirbt, bleibt es allein, wo es aber erstirbt, so bringt es viel Frucht“; das ist, wenn man euer Ding verdammt und zunichte macht, so wird es alles besser werden. Daher die Schrift [Ps. 33, 10.] sagt: „Der HErr macht zunichte den Rath der Heiden, und hindert die Gedanken der Völker“, wie es mit Christo gegangen ist; denn sie eben damit, daß sie ihn kreuzigten, das ausgerichtet haben, daß er ewig ein König ward.

12. So will Joseph auch sagen: Ihr wollet mich umbringen und meine Träume zunichte machen, aber sie sind eben dadurch erfüllt worden; wenn ihr nicht solches Uebel wider mich gedacht und gehandelt hättet, so wäre ich zu diesen Ehren nicht gekommen. Das sollte uns je ein großer Trost sein; aber denen, die nicht Christen sind, ist es umsonst gesagt, sie glauben es nicht. Vernunft kann solches nicht fassen, sie will sich nur rächen, und nichts leiden; der Glaube aber sagt: Ich will gerne leiden, was ich soll, mein GOtt, du wirst es wohl zum Besten wenden.

\begin{quote}
V. 22, 23. Also wohnte Joseph in Egypten mit seines Vaters Hause, und lebte hundert und zehn Jahre, und sahe Ephraims Kinder bis in das dritte Glied. Derselben gleichen die Kinder Machir, Manasse’s Sohnes, zeugeten auch Kinder auf Josephs Schooß.
\end{quote}

13. Das ist ja reichlich vergolten, aus solchem Elende zu solchen Ehren zu kommen; daß dieser Joseph ein Mann von trefflichem Glauben muß gewesen sein, weil ihm so überflüssig vergolten wird. Er hat bei achtzig Jahren im Lande regiert, und ist der Oberste gewesen nach Pharao, vom dreißigsten Jahre an bis in das hundert und zehnte, [hat] dazu bei seinem Leben das fünfte Glied seines Geschlechtes gesehen.

\begin{quote}
V. 24—26. Und Joseph sprach zu seinen Brüdern: Ich sterbe, und GOtt wird euch heimsuchen, und aus diesem Lande führen in das Land, das er Abraham, Isaak und Jakob geschworen hat. Darum nahm er einen Eid von den Kindern Israel, und sprach: Wenn euch GOtt heimsuchen wird, so führet meine Gebeine von dannen. Also starb Joseph, da er war hundert und zehn Jahre alt, und sie salbeten ihn, und legten ihn in eine Lade in Egypten.
\end{quote}

14. Das ist abermal der Glaube auf die Verheißung des Landes Canaan; deß zum Zeichen und Sicherung, dadurch ihrer Kinder Glaube gestärkt würde, befiehlt er auch, seinen Leib mit ihnen zu führen, wenn sie ausziehen würden. [Er] ist fast nach seinem Tode zweihundert Jahre lang in einer Lade behalten.

15. Also haben wir das erste und fast das edelste Buch des Alten Testamentes, das an allen Orten, durch und durch voll merklicher Exempel des Glaubens ist, dazu der Liebe, und vornehmlich des heiligen Kreuzes, und so reich von Figuren unseres HErrn Christi und seines Reiches als kein Buch der Schrift, also, daß beide, Worte und Exempel, nichts anders zeigen und lehren denn den einigen Christum.

GOtt sei ewig Lob!
\quellseite{332}

\chapter{C. Predigt über den Text 1 Mosis 3, 15}

Ipsum semen conteret caput serpentis, des Weibes Same wird der Schlange den Kopf zertreten\footnote{Luther hielt diese und die folgende Predigt am Ostermontage, den 2. April 1526 (nicht 1525, was die Ausgaben bieten, irregeleitet durch die Eißlebensche Ausgabe). Herr D. Enders hat in dem 17. Bande der Erlanger Ausgabe, S. 153 ff., schon dargethan, daß die Predigt im Jahre 1525 nicht gehalten sein könne, weil Luther zu der Zeit nicht in Wittenberg anwesend war; er vermochte aber nicht das rechte Jahr zu bestimmen. Aus dem Index Noachs zu Predigten Luthers, der sich in der Zwickauer Rathsschulbibliothek befindet, ersehen wir, daß Luther an dem von uns angegebenen Tage diese Predigt gehalten habe. (Buchwald, Andr. Noachs handschriftliche Sammlung etc., Bd. I, erste Hälfte, S. XXIV.) Diese Predigten wurden zuerst im Jahre 1564 von Georg Buchholzer, Probst in Berlin, herausgegeben unter dem Titel: „Drei schöne herrliche Sermones des ehrwürdigen Herrn und Vaters D. Martini Lutheri, seligen und heiliger Gedächtniß, zu Wittenberg in der Pfarrkirche in den Osterfeiertagen gepredigt. Die 1. Predigt über den Text Gen. 3.: Ipsum semen conteret caput serpentis. Die 2. Predigt über den Text Gen. 22.: In semine tuo benedicentur omnes gentes. Die 3. Predigt über den Text Gen. 3.: Rubum, quem viderat Moises etc. Am Ende der letzten Predigt, daß Joseph in Egypten auch eine Figur Christi gewesen ist etc. Zu Ehren der römischen königlichen Majestät etc. Jetzt neu und zuvor niemals dergleichen ausgegangen, und durch Georgium Buchholzer, Probst zu Berlin, in Druck verfertiget, Anno 1564.“ Gedruckt wurde diese Schrift „durch Johan Eichorn“ zu Frankfurt an der Oder. In den Sammlungen: Eißlebensche, Bd. I, Bl. 198b; Altenburger, Bd. II, S. 878; Leipziger, Bd. XII, S. 140; Erlanger (1.), Bd. 19, S. 1, und (2.), Bd. 17, S. 153. Wir geben den Text nach der Erlanger Ausgabe (2. Aufl.). Die Predigt, welche in der vorher genannten Ausgabe Buchholzers als die dritte bezeichnet ist, findet sich in den „Predigten über etliche Capitel des andern Buches Mosis“ in diesem Bande nach der Auslegung des dritten Capitels.}

Gehalten zu Wittenberg, Montags in den Osterfeiertagen, Vormittags, den 2. April 1526.

1 Mos. 3, 15. Des Weibes Same wird der Schlange den Kopf zertreten.

1. Lieben Freunde, dieweil wir noch das Osterfest feiern und halten, und begehen den Artikel unseres christlichen Glaubens, da wir sprechen: Am dritten Tage auferstanden von den Todten, dadurch uns Christus, Gottes Sohn, gebracht hat Gerechtigkeit, das ist, Vergebung der Sünden, Ueberwindung des Todes, des Teufels und der Hölle, ewiges Leben und Seligkeit; wie ihr denn das, Gott Lob und Dank, gestern genugsam gehört habt. So wollen wir heute vor uns nehmen die erste Verheißung Gottes, oder das erste Evangelium, so auf Erden im Paradies, Adam und Eva gepredigt worden, da Gott zur Schlange, die Evam verführt, spricht: „Und ich will Feindschaft setzen zwischen dir und dem Weibe, und zwischen deinem Samen und ihrem Samen. Derselbe soll dir den Kopf zertreten, und du wirst ihn in die Ferse stechen.“

2. Das sind die tröstlichen Worte, die Adam und Evam wieder erquicket und erfreuet haben; sonst hätten sie in der Angst des Todes, darein sie durch die Schlange, den Teufel, gebracht, müssen verzagen und verzweifeln, ewiglich sterben und verderben. Und das ist ihr Evangelium gewesen, daß sie ihren Kindern geprediget. Darum hat es auch am ersten Adams Schooß geheißen, daß sie gleich wie in einem Schooß alle Gläubigen versammelt, durch das Wort erhalten und selig worden. Die Predigt hat gewährt bis auf Abraham, als 2080 Jahr; da Abraham Isaak will opfern, da verneuert Gott dem Abraham diese Verheißung, 1 Mos. 22, 18., und spricht, daß in seinem Samen sollen gebenedeiet werden alle Heiden; da ist's Abrahams Schooß genannt worden, wie Christus spricht Luc. 16, 22. Jetzund aber heißt’s Christi Schooß; denn der ist nun kommen, welchen Gott Adam, Abraham verheißen hat.

3. Darum wollen wir auch nicht viel disputiren von der Schlange, als gehen nur die Worte und die Strafe auf die leibliche Schlange allein, wie die gottlosen Juden narren; sondern gehen vornehmlich auf den Teufel, der die Schlange besessen, und durch sie redet, denn keiner Creatur auf Erden von Anfang der Welt sonst gegeben ist die Gabe zu reden, denn allein dem Menschen. Dieweil aber hier die Schlange redet mit Eva, als mit einer leiblichen Stimme und menschlicher Sprache, daß Adam und Eva verstehen können: dazu so redet sie nicht allein schlechte Worte, sondern redet von Gottes Wort
\quellseite{333}
und Gebot, davon die Schlange nichts weiß; ja, die Schlange hat nicht verstanden, was der Teufel durch ihre Zunge und Mund geredet hat, wie auch Augustinus sagt, und verfälscht Gottes Wort mit Lügen, da er spricht: „Ihr werdet nicht des Todes sterben“; das kann keine schlechte Schlange, sondern muß ein hoher Geist sein. Wie denn der Teufel im Anfang ein hoher Geist, und ein guter Engel von Gott geschaffen gewesen ist, der die Schlange besessen hat, und durch sie redet. Es hilft auch nicht die Einrede der Juden: Hat doch Bileams Eselin geredet, warum denn nicht hier die Schlange auch reden sollte? Antwort: Aber dort redet der Engel des Herrn durch die Eselin, wie im 4. Buch Mosis am 22., V. 28. ff., geschrieben; aber hier durch die Schlange thut es der Teufel.

4. Warum gehen diese Worte Gottes zu Adam und Eua wider den Teufel und sein Reich, als, Sünde, Tod, Gottes Zorn und Hölle, und sind gesagt von unserm Herrn Jesu Christo, Gottes Sohn, wie der Mensch werden soll; und darnach, wie sich der Teufel an ihn legen wird, ihn tödten und morden. Aber der Herr Christus soll wieder erstehen, den Teufel überwinden, ihm seine Kraft nehmen, Adam und Eua und alle Gläubigen erlösen, ihnen Gerechtigkeit, Leben und Seligkeit wieder bringen; wie wir denn das durch Gottes Gnade hören wollen.

5. Erstlich ist es eine greuliche Abgötterei gewesen in der Kirche, daß man diesen schönen, herrlichen Text von Christo auf die zarte, keusche, reine Jungfrau Maria, die Mutter Christi, gedeutet und gezogen hat, als sei sie die, die der Schlange den Kopf zertreten soll, das ist, dem Teufel seine Kraft nehmen, die Sünde vertilgen, den Tod überwinden, Leben und Seligkeit wieder bringen, und haben das Wort Ipsum, das ist, der, verkehret, und daraus gemacht Ipsa, das ist, die, als Maria; wie man in allen Kirchen von ihr gesungen; welches doch ihr Amt nicht ist. Man hat sie auch damit nicht geehrt, wie man gemeint, sondern auf das höchste geschändet, daß ihr das zugelegt wird, das ihrem Sohne gebührt und gehört. Darum wollen wir den teuflischen Irrthum fahren lassen, und bleiben bei unserm Herrn Jesu Christo, von dem Moses und alle Propheten geschrieben haben; wie Christus gestern am Ostertag zu seinen Jüngern sagt, wie Luca am letzten [B. 44.] zu lesen.

6. Darum ist nun dieser Same des Weibes Christus unser Herr, und die hebräische Sprache heißt die Kinder, so von Weibern geboren werden, Weibes-Same, das ist, wie wir auf deutsch sagen, ein recht natürlich Kind vom Weibe, oder eines Weibes. Und ist hier nun das erste Stück, daß unser Herr Christus soll werden ein rechter, natürlicher, pur lauter Mensch, geboren vom Weibe, oder nicht vom Manne; gleich als sagte Gott: Nun, du Teufel, der du durch das Weib, Eua (die vom Manne, dieweil er schlief, aber ohne des Mannes Zuthun und Werk, geschaffen und genommen), die Sünde, den Tod und ewige Verdammniß in die Welt gebracht hast: also will ich dir wieder durch ein Kind, so vom Weibe, und nicht vom Manne, kommen soll und geboren werden, ein Glück anrichten; dasselbe Kind oder Weibes Same soll alles an dir rächen; was du nur durch sie zuwege gebracht, das soll über deinen Kopf kommen und ausgehen. Soll nun aber dies Kind von einem Weibe, und nicht von einem Manne, herkommen und geboren werden, so muß das Weib oder die Mutter eine Jungfrau sein.

7. Das hat aber Adam und Eua die Zeit noch nicht verstanden, daß die Mutter soll eine Jungfrau sein, und ohne des Mannes Zuthun gebären. Wie denn Moses im 4. Capitel, V. 1., hernach schreibt, da Eua ihren ersten Sohn Kain geboren, da meint sie nicht anders, der sei der Mann, den ihr Gott verheißen und zugesagt, der sie an der Schlange rächen werde. Darum spricht sie auch: „Ich habe bekommen den Mann, den Herrn“, den Gott Jehova, des Weibes Samen. Also auch, da Noah geboren wird, nach 1057 Jahren, meint der Patriarch oder Altvater Lamech auch nicht anders, sein Noah würde des Weibes Same, oder das Kind sein. Derhalben sagt er: „Der wird uns trösten in unserer Mühe und Arbeit auf Erden, die der Herr verflucht hat“, wie 1 Mos. 5, 29. geschrieben.

8. Aber der Mann oder das Kind muß von Adam oder Lamech nicht gezeuget werden, sondern allein von einem Weibe, das ist, von einer Jungfrau kommen. Wie denn Gott das hernach 1244 Jahr durch den Propheten Jesaiam Cap. 7, 14. öffentlich sagen lassen: „Siehe, eine Jungfrau ist schwanger, und wird einen Sohn gebären, den wird sie heißen Immanuel.“ Da ist es recht offenbar und laut worden, daß dieses
\quellseite{334}
Sohns Mutter muß eine Jungfrau sein, und soll nicht eines Mannes, sondern einer Jungfrau Sohn sein. Denn alles, was vom Mann und Weib geboren wird, das wird in Untugend und Sünden empfangen und geboren, wie der 51. Psalm, V. 7., sagt. Soll nun dieser Sohn ohne Sünde geboren werden, so muß er nicht vom Mann und Weib herkommen, sondern von einem Weibe allein, das ist von einer Jungfrau.

9. Denn wenn seine Empfängniß und Geburt wäre unrein gewesen, wie aller Menschen, so hätte er uns nicht von Sünden und Tod können helfen. Denn ein Sünder oder Dieb kann dem andern nicht vom Galgen helfen, es muß es ein Frommer, Unschuldiger thun. Soll nun uns Christus helfen, so muß seine Empfängniß und Geburt rein und heilig sein, der Heilige Geist muß solches durch eine Jungfrau ausrichten: wie wir in unserm Glauben bekennen, daß es also geschehen ist, und im Neuen Testament durch St. Matthäum [Cap. 1, 18.] und St. Lucam [Cap. 1, 35.], die Evangelisten, beschrieben ist. Darum hat Adam und Eva geglaubt, Christus, Gottes Sohn, würde kommen, und Mensch werden, sie wiederum erlösen und erretten von Sünde, Tod, Teufel und Hölle. Ob er aber von einer Jungfrau sollte geboren werden, das haben sie nicht gewußt, denn es ist ihnen auch nicht so öffentlich und so klar angezeigt. Verhalten haben sie sich auch nicht damit so weit bekümmert, sind einfältig bei der Verheißung Gottes blieben, daß des Weibes Same, unser Herr Christus, Gott und Gottes Sohn, würde kommen und geboren werden, ihnen Gerechtigkeit, Vergebung der Sünden, Leben und Seligkeit wieder bringen.

10. Das andere Stück hier ist: Wie wird es nun des Weibes Samen, dem Herrn Christo, gehen? So spricht Gott zur Schlange: „Und du wirst ihn in die Ferse stechen“, das ist, du Teufel wirst ihn tödten und morden. Denn in die Ferse stechen, heißt, daß der Teufel dies Kind morden und tödten wird. Aber merket eben auf die Worte, wie Gott sagt; spricht nicht: Du Schlange, oder Teufel, wirst des Weibes Samen den Kopf zertreten oder zermalmen: Nein, nein, den soll er ihm ganz lassen, sondern schlecht „in die Ferse stechen“. Wenn einer in die Ferse wird gestochen, der wird versehrt und verwundet, der Kopf, Herz und der Leib bleibt ganz unversehrt, da hat es denn keine Noth. Also auch hier mit des Weibes Samen, da soll der Teufel wohl so viel zuwege bringen, daß Christus ermordet und verwürget wird, und ihm einen Stich in die Ferse thun, aber er soll nichts damit ausrichten, sondern Christus soll den Kopf ganz behalten, den aufheben und hervorgehen, das ist, wieder vom Tod erstehen, als einer, dem des Teufels Stich in die Ferse nicht geschadet hat, sondern der da triumphirt, und den Sieg behalten, und vom Tod erstanden. Das ist vom andern Stück gesagt, wie es des Weibes Samen, dem Herrn Christo, gehen soll.

11. Nun folgt das dritte Stück, was nun des Weibes Same, das ist, Christus, wenn er von des Teufels Stich, das ist, vom Tod wird wieder erstehen, ausrichten soll, als nämlich so sagt Gott: „Und derselbe Same des Weibes soll dir den Kopf zertreten.“ Das ist der rechte Trost, der Adam und Eva wieder zurecht gebracht, erfreut und erquicket hat, daß sie in so großer Furcht und Angst ihres Gewissens hören, wie des Weibes Same, Christus, dem Teufel oder der Schlange, nicht ein Stück vom Schwanz abhauen soll, sondern eben den Kopf zertreten. Denn das sagt man, wenn gleich eine Schlange ein groß Stück vom Schwanz abgehauen, oder am Leibe hart verwundet wird, und ihr nur der Kopf ganz bleibt, so hat es keine Noth mit ihr, heilt sich wieder aus, und bleibt lebendig. Derhalben, wenn man eine Schlange schlagen will, so windet sie sich in einander, und thut allen Fleiß, daß sie nur den Kopf bewahre: wenn sie den unzerhauen davon bringt, so ist sie genesen. Hauet oder zerschlägt man ihr den Kopf, so ist's mit ihr aus, und hat ihre Kraft verloren; und ob wohl noch der Leib und Schwanz grausam und scheußlich anzusehen, so hat er doch keine Macht mehr; denn der Kopf ist weg. Darum spricht hier Gott auch: Christus, des Weibes Same, soll es unterthun, dem Teufel, der alten Schlange, nicht den Schwanz zerhauen, oder den Leib verwunden, sondern soll ihm den Kopf zertreten und zermalmen.

12. Was ist nun des Teufels Kopf? Erstmals ist es die Sünde, die Adam und Eva durch ihren Fall auf alle Menschen, die von ihrem Fleisch und Blut herkommen und geboren, geerbt und gebracht, also, daß sie alle in Sünden empfangen und geboren werden. Durch die Sünde kommt über sie Gottes Zorn und der Tod, die fressen alle Menschen, wie heilig sie
\quellseite{335}
auch sind; darnach kommt die Hölle und der Teufel, die verdammen ewiglich. Da soll nun Christus, wie Johannes der Täufer sagt, als das rechte Osterlämmlein Gottes durch sein Sterben und Leiden die Sünde der ganzen Welt tragen [Joh. 1, 29.] und von Gottes Zorn, dem Tod, Hölle und Teufel, gnädiglich erlösen, und durch sein Auferstehen Gerechtigkeit, Vergebung der Sünden, Leben und Seligkeit, und alles, was Adam und Eva im Paradies verloren haben, wieder bringen, und soll der Christus sagen: Siehe, Sünde, Tod, Hölle und Teufel, ihr habt mich in die Ferse gestochen, getödtet und gemordet, habt mir aber Unrecht gethan, denn ich bin nicht der Mann, der unter euer Reich gehört; derhalben sollt ihr meine Gefangenen sein. Und wirft die unter sich, und tritt sie mit Füßen, und läßt ausdreyen und ausrufen durch sein Evangelium in aller Welt, daß alle, die unter den Feinden und unter der Last des Teufels liegen, sollen sich an ihn halten und an ihn glauben, die sollen gerecht sein, das ist, von Sünde, Gottes Zorn, Tod, Hölle und Teufel erlöst und gefreiet sein, wie St. Paulus sagt Röm. 4, 25.: „Christus ist gestorben um unserer Sünde willen, und erstanden um unserer Gerechtigkeit willen.“

13. Das ist’s nun, was Gott hier sagt: Des Weibes Same soll dem Teufel, der alten, bösen, giftigen Schlange, den Kopf zertreten. Denn also spricht David im 68. Psalm, V. 21., von Christo: „Wir haben einen Gott,\footnote{„Gott“ von uns hinzugefügt.} der da hilft, und den Herrn Herrn, der vom Tod errettet.“ Also sagt auch St. Paulus 2 Tim. 1, 10.: „Christus hat dem Tod die Macht genommen, und das Leben, und ein unvergänglich Wesen ans Licht bracht durch das Evangelium.“ Wie fein lieblich, klar und hell reden aus Einem Munde David und Paulus, und haben doch von einander gelebt 1100 Jahr. Das haben nun Adam und Eva, wie ihnen Gott verheißen, geglaubt; dadurch sind sie auch erhalten und selig worden. Also müssen wir auch selig werden. Das ist nun das dritte Stück, das Christus dem Teufel, als seinem und der ganzen Christenheit abgesagtem Feinde, thun soll. Ach! wer wollte sich doch nun vor dem Feinde, dem Teufel, fürchten, und wollte nicht fröhlich mit dem alten Simeon begehren zu sterben [Luc. 2, 29. ff.]? „Gott aber sei Dank, der uns den Sieg gegeben hat durch unsern Herrn Jesum Christum“ [1 Cor. 15, 57.].

14. Das vierte Stück: Soll aber nun Christus von dem Tode wieder erstehen, dem Teufel seinen Kopf zertreten, von Sünde, Gottes Zorn, Tod, Hölle erlösen, Gerechtigkeit, Leben und Seligkeit wieder bringen, so muß er und kann nicht allein sein des Weibes Same, das ist, schlechter, natürlicher Mensch. Denn so wäre er dem Teufel und dem Tode viel zu schwach gewesen. Denn Sünde, Gottes Zorn, Tod, Hölle und Teufel, sind so eine schwere Last, daß sie alle Menschen unterdrücken in den Tod, wie heilig sie auch sind. Ja, wenn gleich alle Engel, alle Creaturen, Sonne, Mond und Sterne, wie im Hiob [Cap. 31, 23.] geschrieben, diese Last vor Gott tragen sollten, könnten sie nicht bestehen. Derhalben, soll nun des Weibes Same diese Last alle tragen, so muß er auch ewiger, allmächtiger Gott sein. Denn solche Bürde zu tragen ist allein Gottes Werk, wie Jesaias, Cap. 43, 24. 25., klarlich sagt. Darum ist nun dieses Weibes Same, Christus, ewiger, allmächtiger Gott und Schöpfer, dadurch, wie Johannes [Cap. 1, 3.] sagt, alle Creaturen gemacht und geschaffen sind, was gemacht und geschaffen ist.

15. Das haben nun Adam und Eva hier betrachtet, und mit den Augen des Herzens im Glauben gesehen, daß dieser Same des Weibes würde auch rechter natürlicher Gott sein, gleich dem Vater, doch eine andere Person. Darum spricht auch Eva, da sie Kain gebiert, 1 Mos. 4, 1.: „Ich habe bekommen den Mann, den HErrn“, den Gott Jehovah; spricht nicht schlecht, des Weibes Samen, sondern „den Mann, den Gott“, der sich nun als ein rechter Mann gegen der bösen Teufelschlange, die uns in Noth gebracht, wird erzeigen, welcher der rechte Gott ist.

16. So sagt David auch Ps. 68, 21.: „Wir haben einen Gott, der da hilft, und den Jehovah, den HErrn HErrn, der vom Tode erlöset.“ Ja, lieber David, wer ist denn der Gott? So antwortet er im 68. Psalm, V. 8. 9.: „Gott, der du vor deinem Volke hergezogen, da du einher gingest in der Wüste, da bebte die Erde, und die Himmel troffen vor diesem Gott in Sinai, vor dem Gott, der Israels Gott ist.“ Das ist klar und hell genug, daß der Gott, der vom Himmel herunter gestiegen auf Erden, wie Paulus sagt Eph. 4, 8. 9., gestorben, auferstanden, und wieder
\quellseite{336}
aufgefahren in die Höhe, und nimmt das Gefängniß gefangen, und errettet vom Tode, ist der ewige, allmächtige GOtt, der die zehn Gebote auf dem Berge Sinai den Kindern von Israel gegeben hat. Das ist ein gewaltig Zeugniß, daß dieser Christus ewiger GOtt ist. So spricht auch Jeremias am 23, V. 6., und Cap. 33, 15. 16.: „Siehe, es kommt die Zeit, spricht der HErr, daß ich dem David ein Gewächs der Gerechtigkeit erwecken will, und soll ein König sein, der wohl regieren wird, Recht und Gerechtigkeit auf Erden anrichten. Zu derselbigen Zeit soll Juda geholfen werden, und Israel sicher wohnen. Und das wird sein Name sein, daß man ihn nennen wird: HErr, der uns gerecht macht.“

17. Da steht es auch klarlich, daß Davids Gewächs in Einer Person vereinigt soll zwo Naturen haben: eine menschliche, von Davids Fleisch und Blut, die andere göttlich, die er von Ewigkeit vom Vater gehabt. Und steht hier der große Name GOttes Jehovah, den alle Juden bekennen müssen, daß er allein der göttlichen Majestät gehört. Ich habe auch diesen Text den Juden oftmals vor die Nase gehalten, daß Davids Sohn, Christus, wäre allmächtiger, ewiger GOtt; denn hier stände Jehovah; haben mir nichts können drauf antworten. Darum ist dieser Davids oder Marien Sohn nicht ein schlechter Mensch, sondern zugleich wahrer, allmächtiger GOtt, der ewige Gewalt hat, Himmel und Erde und alles erschaffen und erhält; derhalben, daß Menschheit und Gottheit Eine Person ist, und auch rechter GOtt ist. Ist auch dazu der GOtt, der uns gerecht macht, der darum „Davids Gewächs“, das ist, ein rechter Mensch worden, daß er uns durch sein Leiden, Sterben, Blutvergießen und Auferstehen Gerechtigkeit, die Adam und Eva verloren haben, wiederbringen und erwerben, Vergebung der Sünde, GOttes Gnade, Leben und Seligkeit schenken und geben will; denn das heißt „gerecht machen“.

18. Das haben nun Adam und Eva geglaubt. Denn daran liegt alle unsere Seligkeit, daß wir diese beiden Naturen nicht voneinander scheiden, sondern in Einer Person zusammen vereinigt bleiben lassen. Wie denn auch St. Paulus sagt 1 Cor. 2, 8.: „Sie haben den HErrn der Herrlichkeit gekreuziget“; nicht den schlechten Menschen Christum, sondern den gerechten GOtt; nicht nach der Gottheit, sondern nach der menschlichen Natur. Wie die Kirche gesungen am guten Freitage: Dum fabricator mundi mortis supplicium pateretur in cruce etc. Und am heiligen Osterzeit: Crucifixum in cruce laudate etc. Wie auch Petrus spricht, 1. Ep. 4, 1. f.

19. Aber das hat GOtt dem Teufel verborgen. Denn wenn er’s gewußt hätte, daß dieser Christus, der Jungfrau Marien Sohn, des Weibes Same, wäre ewiger, allmächtiger GOtt gewesen, so hätte er sich an ihn nicht gelegt, und die Hände an ihm verbrannt. Denn er hat wohl gesehen im Himmel, wie er GOtt des Vaters einig Herz gewesen, in dem GOtt der Vater sein ganz göttliches Wesen und Substanz gegossen. Darum hat er seine göttliche Kraft wohl erfahren, da er von wegen des Sohns aus dem Himmel in Abgrund der Hölle gestoßen ward, Jes. 14, 12. Luc. 10, 18., darum, daß der Satan dem Sohne nach seiner Ehre stand, und ihm wollen gleich sein. Das könnte GOtt nicht leiden. Erwegen verbirgt es GOtt dem Teufel, daß er ihn in dem Sohne, der Mensch worden, so fahen will, und thut wie ein Fischer, der Fische fangen will: der bindet eine Schnur an einen Stecken, und unten an die Schnur eine scharfe Angel; daran hängt er ein Würmlein, und wirft das ins Wasser. Da kommt denn der Fisch, sieht das arme Würmlein, sieht aber nicht die scharfe Angel in dem Würmlein verborgen, und beißt drein, meint, er bekomme ein gut niedlich Bißlein: aber die Angel bleibt ihm im Munde oder Halse stecken, wird also gefangen und ergriffen. Also thut GOtt der Vater auch: da nimmt er seinen eingebornen, geliebten Sohn, den hänget er an die Linee oder Schnur der Patriarchen und Propheten, muß Adams, Abrahams, Davids Fleisch und Blut annehmen, und läßt ihn aus dem hohen Himmel in die Welt kommen; da der Teufel sieht, wie Christus als ein armer, elender Wurm, wie er sich im 22. Psalm, V. 7., selbst so nennt, daß ist, als ein armer, geplagter Mensch einher geht, leidet Hunger, Durst, Kälte, Frost und Hitze, weint und hat eitel Elend und Jammer in dieser Welt, geboren wie ein anderer Mensch.

20. Es weiß aber der Satan nicht, daß dies in diesem Christo verborgen, daß er ewiger, allmächtiger GOtt ist, gleich dem Vater; doch eine andere Person. Darum denkt er: Ob wohl Christus Wunderzeichen thut, haben doch die Propheten auch Wunderzeichen gethan; ich habe sie aber alle aufgesressen, wie hoch und heilig sie gewesen, will den Christum auch wohl verschlingen
\quellseite{337}
, und meint, er wolle an dem Manne ein niedlich Büßlein haben; hebt an, frißt den Christum, und verschlingt ihn. Aber es bekommt ihm, wie dem Hunde das Gras. Denn der Christus bleibt ihm im Halse stecken, und muß ihn wieder speien, wie der Walfisch den Propheten Jonas, Jona 2, 11., und eben mit dem Freßen ermordet und erwürgt sich der Teufel selbst, und wird von Christo gefangen. Denn dieser Mensch, Christus, ist rechter, wahrer Gott, der nicht im Tode kann bleiben, sondern muß wieder hervor, auferstehen und leben, und den Teufel also fangen. Und hier wird wahr das Wort des Propheten Hoseä, der vor Christo gelebt 760 Jahr, der da spricht am 13. Capitel, V. 14., von Christo: „Ich will sie erlösen aus der Hölle, und vom Tode erretten. Tod, ich will dir ein Gift sein, Hölle, ich will dir eine Pestilenz sein.“ Darum sagt auch der Prophet Hoseas zuvor, Cap. 6, 1—3., daß die Christen zur Zeit Christi sagen werden: „Kommt, wir wollen wieder zum Herrn gehen; denn er hat uns zerrissen, er wird uns heilen; er hat uns geschlagen, er wird uns auch verbinden; er macht uns lebendig nach zweien Tagen, er wird uns am dritten Tage erquicken, daß wir vor ihm leben werden. Dann werden wir die Art drauf haben, und fleißig sein, daß wir den Herrn erkennen.“ Das ist gestern alles durch unsern Herrn Jesum Christum, Gottes Sohn, da er am dritten Tage vom Tode auferstanden, erfüllt und vollbracht worden; da sind wir von unsern Sünden erlöst, gerechtfertigt, geheilt, verbunden und gesund gemacht; da sind wir erquickt und lebendig worden vom ewigen Tode, und mangelt nirgend an, denn daß wir solches von Herzen glaubten, und den Herrn Christum erkennen lernten, und möchten fröhlich singen: „Deß sollen wir alle froh sein, Christ will unser Trost sein, Kyrieleis.“

21. Das ist nun die hohe Weisheit Gottes, die allein vom Heiligen Geist muß herkommen. Denn kein Mensch aus menschlicher Vernunft verstehen kann, daß dieses Weibes Same, unser Herr Christus, nicht allein Mensch ist; sondern allmächtiger Gott, gleich dem Vater und Heiligen Geiste, in Einer Person, nicht in zweien Personen, wie die beiden Ketzer, Nestorius und Eutyches, genarrt haben, sondern ist nur Eine Person, wie ein Mensch mit Leib und Seele nur Eine Person ist.

22. Aber gleichwohl ist diese einige Person, Christus, rechter, wahrer, allmächtiger Gott und Mensch; ist doch nicht der Vater, oder der Heilige Geist, sondern der Sohn, die andere Person in der Gottheit; und doch gleicher Gott mit dem Vater und dem Heiligen Geiste in dem göttlichen Wesen oder Substanz. Das haben die Juden zu Christi Zeiten nicht verstanden, verstehen es auch noch nicht; desgleichen die Heiden und Türken, der Papst mit seinen Genossen auch nicht; müssen toll und thöricht darüber werden, denn sie wollen es mit ihrer Vernunft auspeculieren und erforschen. Aber da wird nichts aus, es muß allein vom Himmel kommen, wie Christus Matth. 16, 17. zu Petro sagt: „Selig bist du, Simon, Jonas Sohn, denn Fleisch und Blut hat dir das nicht offenbaret, sondern mein Vater im Himmel.“ Darum muß man’s durchs Wort im Glauben fassen, wie Adam und Eva mußten thun; sonst ist’s verloren.

23. Derhalben ist zwischen Adam und Eva ihrem, und unserm Glauben,\footnote{Erlanger: Adam und Eva, ihrem und unserm Glauben 2c.} kein Unterschied; allein dieser,\footnote{Erlanger: diese.} daß sie geglaubt, daß Christus Gottes Sohn, als wahrer, allmächtiger Gott, soll des Weibes Same, das ist, rechter Mensch werden, und in die Welt kommen: den würde der Teufel in die Ferse stechen, das ist, tödten und morden; aber er würde erstehen vom Tode, und dem Teufel seinen Kopf zertreten, von Sünden, Gottes Zorn, Tod und Hölle erlösen, Gerechtigkeit, Leben, Seligkeit wiederbringen. Wir aber glauben nun, daß es nicht noch zukünftig, und geschehen soll, sondern daß es schon geschehen, und alles vollbracht und erfüllt ist, wie Gott Adam und Eva zugesagt und verheißen hat, und mangelt nirgend an, denn daß wir warten der Auferstehung des Leibes oder unseres Fleisches. Darum lasset uns ja Gott getreulich bitten, daß wir in diesem Artikel, daß Christus, als wahrer Gott, Mensch worden, gestorben, und am dritten Tage vom Tode erstanden, Gerechtigkeit, Vergebung der Sünden, Leben und Seligkeit wiedergebracht hat, wie Adam und Eva gethan, mit festem Glauben mögen bleiben. Thun wir das, so wohl uns; wo nicht, so wehe uns allen! Daß wir aber darinne mögen festiglich bleiben bis an unser Ende, dazu helfe uns Gott der Vater, Gott der Sohn, und Gott der Heilige Geist, Amen.
\quellseite{338}

\chapter{D. Predigt über den Text 1 Mosis 22, 18}

\begin{quote}
In semine tuo benedicentur omnes gentes, in deinem Samen sollen alle Völker auf Erden gesegnet werden.
\end{quote}

Gehalten zu Wittenberg in der Pfarrkirche, Montags in den Osterfeiertagen zur Vesper, den 2. April 1526.\footnote{*) Wegen der Zeitbestimmung 2c. vergleiche die erste Anmerkung zur vorigen Predigt.}

1 Mos. 22, 18. In deinem Samen werden gesegnet werden alle Geschlechter auf Erden.

1. Lieben Freunde, dieweil ihr heut vor Mittage gehört habt aus Mose das erste Evangelium, so Gott selber Adam und Even im Paradies gepredigt hat, dadurch sie auch getröstet, gerecht, erhalten und selig worden; sonst hätten sie in Gottes Zorn und ihren Sünden müssen sterben und verderben. Jetztund aber wollen wir vor uns nehmen die andere Verheißung, oder das andere Evangelium Gottes, zu Abraham geschehen, als er seinen Sohn Isaac wollte opfern, nach der Welt Anfang 2080) Jahr, da Gott Abraham mit einem Eide verheißt und spricht: „In deinem Samen sollen Col Goym“, das ist, „alle Heiden auf Erden gesegnet werden“, wie im ersten Buch Mosis am 22. Capitel, V. 18., geschrieben.

2. Nun wollen wir hier nicht viel disputiren, ob dieser Text von unserm Herrn Jesu Christo gesagt ist oder nicht; denn da haben wir unsern Doctor und Lehrer der Heiden, der von Gott selber, und unserm Herrn Jesu Christo, zum Doctor Theologiä promovirt und gemacht ist, und steht das testimonium seines Doctorats in der Apostelgeschichte, Cap. 9, 3. ff., reichlich und herrlich beschrieben, als Sanct Paulus; der sagt klärlich mit hellen Worten zum Galatern am dritten, V. 8., daß diese Worte Mosis allein von Christo gesagt sind, und er sei auch der Same allein, und sonst niemand.

3. Ueber das so spricht unser Herr Jesus Christus selber zu den Juden, Joh. 8, 56.: „Abraham, euer Vater, ward froh, daß er meinen Tag sehen sollte, und er sahe ihn und freuete sich.“ Ja, wo sahe er ihn? In diesem Spruche, da er hörte, wie sein Same soll Gott und Mensch sein, der alle Heiden segnen, das ist, von Sünde, Tod und Teufel erlösen, ewiglich gerecht, lebendig, heilig und selig machen soll.

4. Derhalben wollen wir nun zum Text greifen, und am ersten hören, was das für ein „Fluch“ ist, davon dieser Christus, Abrahams Same, solle alle Heiden erlösen, und was „segnen“ ist. Zum andern, wie der Abrahams Same muß ein rechter Mensch sein und sterben. Zum dritten, wie der Abrahams Same muß wieder vom Tode erstehen, soll er anders alle Heiden segnen. Zum vierten, wie er auch muß rechter Gott sein, doch eine andere Person, denn Gott der Vater und der Heilige Geist. Das wollen wir auf diesmal vor uns nehmen zu handeln, so viel als Gott seine Gnade dazu verleihen und geben wird.

5. Da hören wir nun im Text klärlich, daß Gott sagt zu Abraham: „In deinem Samen sollen alle Heiden auf Erden gesegnet werden.“ Sollen nun alle Heiden, das sind alle Menschen, in diesem Christo gesegnet werden, so müssen vormals alle Heiden unter dem Fluch sein; sonst könnte sie der Same nicht segnen.

6. Nun, was ist es denn für ein Fluch? So spricht unser Doctor, Sanct Paulus, zum Römern am fünften Capitel, V. 12.: „Durch Einen Menschen ist kommen die Sünde in die Welt, und der Tod durch die Sünde, und ist also der Tod zu allen Menschen durchgedrungen, dieweil sie alle gesündiget haben.“ Das ist ein rechter Meister, der kann Adams Fall im Paradies recht abmalen, was er mit dem Apfelbiß zuwege gebracht, als nämlich, daß er über alle Menschen, niemand ausgenommen, ohne unsern Herrn Jesum Christum, einen solchen greulichen Fluch gebracht und geerbt, daß sie in Sünden empfangen und geboren werden, und dadurch des Todes, des Teufels, ewiglich verdammt und verloren sind.

7. Dieser Fluch heißt peccatum originale, die Erbsünde, oder Natursünde, die wir von Natur von unsern Eltern empfangen haben in
\quellseite{339}
Mutterleibe; wie der ein und fünfzigste Psalm, V. 7., sagt: „Siehe, ich bin aus sündlichem Samen gezeuget, und meine Mutter hat mich in Sünden empfangen.“ Und zum Ephesern am andern, V. 3.: „Wir waren von Natur Kinder des Zorns“; das ist, wir bringen von Natur, wie wir empfangen und geboren werden, die Sünde mit in die Welt, und durch die Sünde kommt Gottes Zorn und Tod, daß wir alle verloren und verdammt sind. Und diese Erbsünde ist der rechte Brunnquell, daraus andere wirkliche Sünden der Menschen entspringen und herkommen.

8. Von diesem Stücke weiß kein Jude oder Türke, ja, der Pabst auch nicht; aber das hat Adam, Abraham ihren Kindern mündlich gepredigt, desgleichen die Propheten, bis daß es durch den König David klärlich beschrieben worden. Aber wir Christen müssen nun aus dem Neuen Testament von diesem Fluch, wie greulich und groß der ist, daß er auch alle Menschen wegreißt und hinnimmt in den Tod; wie heilig sie auch sind, so hilft es nichts; es kann keiner den Hals aus der Schlinge ziehen [Ps. 89, 49.], so schwer ist die Sünde, Gottes Zorn, Tod, Hölle und Verdammniß, daß wir von unsern Eltern, Adam und Eva, ererbt haben. Das ist nun von dem Fluche gesagt, der über alle Heiden kommen ist; wollen nun vom Segen hören, den dieser Same, Jesus Christus, bringen soll.

9. Nun, dieser Segen ist nicht eines schlechten Menschen Segen, da ein Mensch den andern mit grüßt, einen guten Morgen oder Tag bietet oder wünscht, wie denn die Menschen einander nicht anders segnen können; sondern dieser Segen ist Gottes Segen, den er auch allein geben kann und will, durch [das] Evangelium, so er verheißen und zugesagt. Solcher Segen ist nicht ein schlecht, ledig Wort, das uns viel Gutes wünscht, und nichts daraus wird, sondern gibt und bringt alles, was das Wort verspricht und zusagt. Derhalben ist dieser göttliche Segen, in Abrahams Samen verheißen und versprochen, ein gewaltiger, kräftiger und lebendiger Segen, der da bringt und ausrichtet, was er zusagt oder segnet.

10. Denn er ist von Gott verheißen und gegeben wider den Fluch, darein uns die teuflische alte Schlange gebracht hat durch Adams Ungehorsam und Sünde. Und hiermit hat Gott erneuet die Verheißung von des Weibes Samen, und soll nun Abrahams Same heißen, und hernach weiter, Davids Same genannt [2 Sam. 7, 12. 13.], und zuletzt der reinen Jungfrau Marien Sohn worden ist.

11. Darum heißt hier „Segen in Abrahams Samen“, wie ihr heute vor Mittage gehört habt, „der Same des Weibes soll der Schlange den Kopf zertreten“, das ist, er soll die Sünde und Hölle wegnehmen, und Gerechtigkeit, Unschuld, Leben und Seligkeit wiederbringen. Denn Sünde und ewiger Tod ist der Fluch, wie gesagt ist, darunter wir müßten ewiglich (wo wir nicht durch diesen Samen wiederum geheilet, das ist, von Sünden frei und gerecht, vom Tode lebendig, heilig und selig gemacht würden, und solches in aller Welt durch das Evangelium allen Völkern und Heiden gepredigt und vorgetragen würde) verdammt werden.

12. Das hat der liebe Patriarch oder Erzvater Abraham wohl verstanden; darum hat er sich auch gefreut, wie Christus sagt. Und ich halte, das wird auch ein Stück von dieser Freude gewesen sein, da ihm Gott den Isaac verheißt, er auf sein Angesicht nieder fällt, und im Herzen vor Freuden lacht, 1 Mos. 17, 17.

13. Das soll auch sein unsere Freude, die wir haben sollen, daß wir in diesem Samen Abrahams, Davids und Marien, der Jungfrau, gesegnet werden; ja, wir Col Goym, alle Heiden, wollen uns deß annehmen, und rühmen den Glauben, wie Abraham gethan. Und wollen mit dem Segen trotzen wider Sünde, Gottes Zorn, Tod, Teufel und Hölle, singen und sprechen: Durch Christum, Abrahams Samen, haben wir Gerechtigkeit, das ist, Vergebung, Abwaschung der Sünden, Erlösung vom Tode, Teufel und allem Uebel, und Gott dem Vater sei Lob und Preis, daß er uns seinen eingebornen Sohn, Jesum Christum, geschenkt und gegeben hat, daß er ist worden unsere Weisheit, Gerechtigkeit, Heiligung und Erlösung; deß wollen wir uns rühmen, wie St. Paulus 1 Cor. 1, 31. sagt: „Wer sich rühmen will, der soll sich des Herrn rühmen“ [Jer. 9, 23.] und keines andern. Das ist nun von Abrahams Segen gesagt, der unser Trost, Freude, Leben und Seligkeit in Ewigkeit sein soll. Wollen nun hören, was Abrahams Same für eine Person sein muß, der jenen solle alle Heiden.

14. Da steht es nun klar und hell, Christus soll Abrahams Same sein, das ist, ein recht
\quellseite{340}
natürlich Kind und lauter Mensch, der Abrahams Fleisch und Blut an sich habe.

15. Soll er nun ein rechter Mensch sein, so muß er sterben; denn was ein Mensch ist, das muß sterben [Hebr. 9, 27.]. Das Urtheil ist über Adam und Even, die von dem Baum gessen, und über alle ihre Nachkommen, die von ihnen geboren werden, auch gegangen, daß sie, als Menschen, sterben müssen. Dieweil aber Christus Abrahams Saame, das ist, ein rechter Mensch sein soll, so muß er auch sterben, denn er hat alle Eigenschaft der menschlichen Natur an sich genommen, ausgenommen die Sünde. Das ist nun das andere Stück, daß dieser Christus, als ein rechter, wahrer, lauterer Mensch, sterben soll.

16. Zum dritten: Soll nun dieser Christus segnen, das ist, gerecht machen, Vergebung der Sünden, Erlösung vom Tode, Teufel und Hölle, Leben und Seligkeit wiederbringen, so muß er nicht im Tode bleiben; sonst könnte er nicht segnen, oder, wie Jeremias, Cap. 33, 15. 16., sagt, gerecht machen; sondern muß wieder erstehen vom Tode, und in seinem Namen allen Heiden und Völkern solchen Segen, Gerechtigkeit, Buße und Vergebung der Sünden predigen und verkündigen lassen; wie der Herr Christus zu seinen Jüngern am Ostertage selbst sagt, und Lucä am letzten, V. 47., geschrieben steht. Wie wir Christen auch mit gläubigen Herzen singen: „Wär er nicht erstanden, so wär die Welt vergangen; sint daß er erstanden ist, so loben wir den Vater Jesu Christ, Kyrieleis.“ Das ist recht gesungen. Das hat Abraham gesehen in dem Worte, daß der Mann, sein Sohn Christus, muß sterben und wieder auferstehen vom Tode, Gerechtigkeit, Vergebung der Sünde, Leben und Seligkeit bringen allen Heiden.

17. Zum vierten: Soll nun dieser Saame Abrahams solchen kräftigen und mächtigen Segen geben und unter den Heiden ausrichten, so muß er nicht allein Abrahams Saame, das ist, ein bloßer, lauterer Mensch sein, der nicht mehr kann, denn guten Tag oder Abend zu uns sagen, welches alle Menschen wohl thun können, sondern muß der rechte, natürliche, wahre, allmächtige, einige Gott sein, der solchen starken, thätlichen Segen gewaltiglich in der Hand habe.

18. Denn Sünde, Tod, Hölle und Teufel aufheben und überwinden, Gerechtigkeit, Leben und Seligkeit wiederbringen, sind nicht Menschen- oder Engelwerk, sondern allein der einigen, ewigen, göttlichen Majestät, Schöpfers Himmels und der Erde, Jes. 63, 3. ff. Darum muß dieser Abrahams Saame rechter, wahrer, ewiger, allmächtiger Gott sein, gleich dem Vater von Ewigkeit; doch eine andere Person.

19. Und wir Christen müssen uns wohl vorsehen, wie mir auch heute vor Mittage gesagt haben, und jezuund noch mehr davon sagen wollen, daß wir die Person Christi nicht trennen, noch die zwo Naturen, als das göttliche und menschliche Wesen, nicht in Eine Natur oder Wesen mengen, sondern die Natur oder Wesen hier unterscheiden, und die Person einig behalten.

20. Denn hier hat sich der Teufel durch seine Gesellen wunderlich verdreht und kraus gemacht, daß sie entweder die Gottheit und Menschheit Eine Natur und Wesen, oder aber zwo Personen haben machen wollen, wie etliche Ketzer gethan haben.

21. Nestor, ein Patriarch oder Bischof in der kaiserlichen Stadt Constantinopel, zur Zeit des Kaisers Theodosii des Jüngern (wider den das Concilium zu Ephesio, eben das Jahr, da Augustinus gestorben, nach Christi Geburt 431\footnote{. Eislebensche Ausg. 435. — Augustinus starb den 28. August 430.} Jahr, gehalten worden), hat öffentlich gelehrt, Gott, oder Jesus Christus, GOttes Sohn, sei nicht geboren aus Maria der Jungfrau, gestorben, und auferstanden. Denn GOtt sei ein Geist [Joh. 4, 24.], der nicht kann geboren werden, leiden oder sterben, sondern allein der Mensch, Marien Sohn.

22. Das heißt gemacht aus Christo zwo Personen, und hier wird zertrennt die einige Person Christi, daß eine andere Person ist, die geboren, gestorben und auferstanden, eine andere Person, die nicht geboren, gestorben und auferstanden ist. Und wird also eine jegliche Natur für sich selbst eine Person, und zween unterschiedliche Söhne, der Christi, einer GOttes, der andere Marien Sohn; welches eben so viel ist, als daß GOtt nicht ist Mensch worden, gestorben und auferstanden.

23. Und ist öffentlich wider die Schrift, die da sagt Johannes am ersten, B. 14.: „Das Wort ist Fleisch worden.“ Der Engel Gabriel spricht Lucä am ersten, B. 35.: „Das Heilige, das aus dir geboren wird, wird GOttes Sohn genannt werden.“ St. Paulus Galater am vierten,
\quellseite{341}
B. 4.: „Gott hat seinen Sohn gesandt, von einem Weibesbilde geboren.“ Item, 1 Cor. 2, 8.: „Die Fürsten dieser Welt haben den Herrn der Majestät gekreuziget.“ Und Apost. 20, 28. spricht St. Paulus: „Gott hat die Kirche mit seinem Blute erworben“, so doch, nach der Vernunft zu reden, Gott kein Blut nicht hat. Wie wir denn auch im Kinderglauben sprechen: Ich glaube an Jesum Christum, Gottes Sohn, unsern Herrn, der empfangen ist vom Heiligen Geist, geboren, gestorben und auferstanden; sagen nicht, daß Gottes Sohn ein anderer ist, denn der von Marien geboren, gestorben und auferstanden. Darum ist es ein einiger Sohn, und nicht zween Söhne. Das sei gesagt von Nestorio, dem Ketzer.

24. Der andere Ketzer, Eutyches (welcher ist gewesen ein Abt zu Constantinopel, wider welchen das vierte Hauptconcilium zu Chalcedon in Ponto oder Asia, im vierten Jahr des Kaisers Marcian, gehalten worden, nach Christi Geburt im 451.)\footnote{Eislebensche Ausgabe: 455.} der lehret auch öffentlich, daß der Mensch, Jesus Christus, von Marien geboren, ist nicht der Schöpfer Himmels und der Erde, den man anbeten soll. Dieser Ketzer trennt abermal die Person Christi, und macht aus Einer Person zwo Personen.

25. Der Nestorius, der Ketzer, trennt die Person damit, daß er die Menschheit von der Gottheit reißt, und macht aus einer jeglichen Natur eine besondere Person, daß allein der Mensch Christus geboren, gestorben und auferstanden ist. Eutyches aber, der Ketzer, trennt und reißt die Gottheit von der Menschheit damit, daß er aus einer jeden Natur eine besondere Person macht: eine, die Himmel und Erde geschaffen hat, und die man anbeten soll, als die göttliche Natur allein; trennt und sondert die göttliche Natur von der menschlichen, als, daß Marien Sohn ist nicht der Schöpfer Himmels und der Erde, den man anbeten soll.

26. Aber das ist auch wider die heilige Schrift und unsern christlichen Glauben. Denn wenn wir sagen: Christus, der Mensch, ist Schöpfer Himmels und der Erde, an den man glauben und vertrauen, oder den man anbeten soll, so reden wir’s nicht von dem bloßen und abgesonderten Menschen, von, außer und ohne Gott, sondern reden von dem Menschen, oder von der Person, die zugleich Gott und Mensch ist in Einer Person, ungesondert und unzertrennt, nämlich de Deo incarnato, von dem ich nicht in abstracto oder absolute, wie es die Alten genannt, sondern in concreto also sagen muß: Christus, Gottes und Marien Sohn, ist Schöpfer Himmels und der Erde, den man ehren und anbeten soll; welcher hat den Tod überwunden, Sünde vertilgt, Hölle zerbrochen, und durch sein Leiden, Sterben und Auferstehen Gerechtigkeit, Vergebung der Sünden, ewiges Leben und Seligkeit wiederbracht; der mit Gott dem Vater und Heiligem Geiste ein einiger Gott ist, und mit seiner Gottheit und Menschheit eine einige und unzertrennte Person ist. Das müssen wir wohl lernen. Ich habe es euch nach der Länge desto weitläufiger gesagt, daß ihr es desto besser begreifen, verstehen und behalten könnet. Denn ich sehe und vermerke, daß der Teufel durch die neue Secte der Sacramentirer diese alte Ketzerei wieder hervorzubringen will, und diese Person Christi trennen und theilen.

27. Darum warne ich und bitte, lernet diesen Artikel wohl, und lasset euch nicht die Schwärmer irre machen und verführen. Denn in diesem Stücke stoßen und ärgern sich Juden, Türken, ja, der Papst, und unsere neue Secte, die Sacramentirer. Aber das soll unser Christen-Ruhm und Freude sein, daß Gottes Sohn, Jesus Christus, der rechte einige Gott ist, mit dem Vater und Heiligem Geist, der ist um unser willen Mensch worden.

28. Derwegen sagen wir recht: Gott hat gelitten, Gott ist gestorben und auferstanden, aber nach dem Fleisch, wie St. Petrus spricht 1 Petr. 3, 18. Denn Gott und Mensch ist Eine Person. Denn was der Mensch Christus thut, redet und leidet, das thut, redet und leidet Gott; und was Gott thut, redet und leidet der Mensch Christus. Denn der ist zugleich Gottes und Marien Sohn, und hat doch zwo unterschiedliche Naturen, als, die göttliche und menschliche. Das alles hat der liebe Abraham in dieser Verheißung Gottes erkannt, und in seinem Herzen gewiß geglaubt und gesehen.

29. Aber eines müssen wir auch noch melden und anzeigen. Denn dieweil unser Herr Jesus Christus nicht ist die Person, die zu Abraham sagt: „In deinem Samen sollen alle Heiden gesegnet werden“, so muß er eine andere, unterschiedliche Person sein; denn der Gott, der zu Abraham spricht: „In deinem Samen“,
\quellseite{342}
der ist nicht Abrahams Same, sondern redet als von einem andern, der Abrahams Same noch werden soll. Daraus folgt denn gewaltiglich der Unterschied der zwo Personen; und bleibt doch der einige, ewige, unzertrennte GOtt in seinem einigen göttlichen Wesen. Wo nun diese zwo Personen sind, da ist auch die dritte, als, der Heilige Geist, der nun solches durch den Engel zu Abraham mündlich gesagt, aber durch Mosen hat schreiben lassen, wie der heilige Petrus [2. Ep. 1, 21.] spricht: „Die heiligen Menschen Gottes haben geredet, getrieben von dem Heiligen Geist.“

30. Da haben wir unsern Artikel des Glaubens, daß drei unterschiedliche Personen in der Gottheit sind, als, GOtt der Vater, Sohn und der Heilige Geist; und ist doch nur Ein einiger, ewiger, allmächtiger GOtt.

31. Wie denn auch Moses, 1 Mos. 18, 1. ff., schreibt, daß GOtt Abraham im Hain Mamre, da er in seiner Thüre saß, erschienen, und „da er seine Augen aufthat, da standen drei Männer gegen ihm, und er lief ihnen entgegen, und fiel vor ihnen nieder, und sprach: Herr, habe ich Gnade funden vor deinen Augen, so gehe nicht vor deinem Knecht vorüber. Man soll euch ein wenig Wasser geben“. Hier redet Abraham mit GOtt in singulari, als mit Einem, da er spricht: „Herr, habe ich vor deinen Augen Gnade funden.“ Darnach in plurali, als mit vielen, da er aber spricht: „Man soll euch ein wenig Wasser geben.“ Und der Text sagt deutlich, daß dies Gesicht oder Erscheinung ist GOtt selbst gewesen, der ihm vor seiner Thüre erschienen ist.

32. Wiewohl die Juden hiewider sagen, daß man diesen Text soll verstehen von den Engeln Gottes, die Abraham erschienen sind, und wollen das beweisen aus dem 19. Capitel hernach, mit den zween Engeln, die gegen Sodom gehen. Aber das ist gar ein Anderes mit diesen dreien, die mit Abraham reden und essen, als ein einiger GOtt, wie das ganze Capitel zeuget. Denn da steht der Text klar und hell: der HErr ist es, der ihm erschienen in drei Personen, hat sie auch alle drei als Einen angebetet. Darum Abraham die heilige Dreifaltigkeit recht erkannt und geglaubt.

33. Und wir Christen müssen auch das lernen erkennen und glauben, denn es liegt alle unsere Seligkeit daran. Dieser Artikel sondert uns ab von Juden, Türken und Heiden, daß in dem ewigen, einigen, göttlichen Wesen sind drei unterschiedliche Personen, und müssen uns ja wohl vorsehen, daß wir die Personen nicht mengen, noch das Wesen trennen, wie viel Ketzer gethan haben.

34. Darum müssen wir allein nach der heiligen Schrift richten, denn unsere Vernunft gilt hier nichts, und bei dem Worte Credo bleiben.

35. Der Unterschied aber der drei Personen in dem göttlichen Wesen ist dieser. Erstlich, des Vaters, als der ersten Person, welcher ist der Brunn oder Quelle der Gottheit, daß der ewiger, allmächtiger GOtt ist von ihm selbst, und hat von Ewigkeit geboren in seinem väterlichen Herzen den Sohn, und durch die ewige Geburt ihm gegeben seine Gottheit.

36. Die andere Person aber, als der Sohn, ist ewiglich vom Vater geboren, der hat seine Gottheit nicht von ihm selbst, allein vom Vater.

37. Die dritte Person, als der Heilige Geist, der ist nicht geboren, wie der Sohn, sondern geht zugleich aus vom Vater und Sohn, von denen er auch von Ewigkeit seine Gottheit hat.

38. Also ist es nun kurz gesagt: Der Vater, der gebiert; der Sohn, der geboren ist; der Heilige Geist, der von beiden, als Vater und Sohn, ausgeht. Mit diesem Glauben können wir dem Teufel, Juden und allen Ketzern Widerstand thun.

39. Darum lasset uns die drei Personen nicht in Eine Person mengen, sondern bleiben einfältig dabei, daß drei unterschiedliche Personen in dem einigen, ewigen, göttlichen Wesen sind, die doch alle drei gegen uns und allen Creaturen Ein einiger GOtt, Schöpfer, Erhalter und Wirker aller Dinge ist.

40. Das hat Abraham, wie Moses schreibt, und [ihr] jetzt gehört habt, an den dreien Männern, die ihm erschienen, erkannt und gesehen. Und sind derhalben Abraham und Moses zween gute Christen gewesen. Denn Moses verkündigt uns mit seinem Schreiben, daß GOtt Abraham einen solchen reichen Segen verheißen, der mit dem ganzen Neuen Testament übereinstimmt, daß GOtt seinen Sohn (der gleicher GOtt mit ihm ist) soll senden, Abrahams Fleisch und Blut anzunehmen. Ja, wenn er jetzt hier wäre, würden ihn seine eigenen jetzigen Juden (wo sie anders seines Geblüts sind, wie ich denn nicht glaube) als einen Erzsetzer, wie manchmal schier in der
\quellseite{343}
Wüste geschehen wäre [2 Mos. 17, 4. 4 Mos. 14, 10.], tödten und steinigen.

41. Das wollen wir auf diesmal gesagt haben; denn ich habe es fast lange gemacht, wie unser Herr Jesus Christus, Gottes Sohn, rechter Abrahams Same und Mensch worden ist; und durch sein Leiden und Sterben uns hat vom Fluche errettet, und als ein ewiger, allmächtiger Gott vom Tode erstanden; doch eine andere Person, denn Gott der Vater und der Heilige Geist; und wiederum, alle Heiden reichlich gesegnet, das ist, Gerechtigkeit, Vergebung der Sünde, Leben und Seligkeit, Erlösung vom ewigen Tod, Teufel und Hölle wiedergebracht.

42. Das hat der liebe Abraham alles geglaubt, und durch den Glauben ist er, wie St. Paulus spricht [Röm. 4, 3.], mit allen Gläubigen gerecht und selig worden. Darum heißt auch das selige Wort Abrahams Schooß, wie Luc. 16, 22. geschrieben. Jetzt heißt es Christi Schooß; denn der ist kommen, und hat alles erfüllt, wie Abraham verheißen und zugesagt ist.

43. Und mangelt nun nirgends an, denn daß wir es auch (wie Abraham gethan) von Herzen glauben, und mit dem Munde bekennen, so werden wir gerecht und selig, wie St. Paulus sagt zu den Römern am 10., V. 10. Zu solchem rechten, wahrhaftigen und lebendigen Glauben, daß wir auch darinne bis an unser Ende beständig verharren und bleiben mögen, beide im Leben und Sterben, im Guten und Bösen, dazu helfe uns Gott der Vater, Gott der Sohn, Gott der Heilige Geist, als drei Personen, und doch nur Ein einiger, ewiger, allmächtiger Gott, in Einem göttlichen Wesen; dem sei Ehre, Preis, Lob und Dank für diese Wohlthaten in Ewigkeit, Amen.
